% Les grands groupes humains — Géographie Tle E — Cours signé OnBuch+
% Programme officiel MINESEC. Compiler avec : tectonic N02-grands-groupes-humains.tex
\newcommand{\DOCMATIERE}{Géographie}
\newcommand{\DOCNIVEAU}{Tle E}
\newcommand{\DOCMODULE}{Module 1 : Le Cameroun}
\newcommand{\DOCLECON}{N02}
\newcommand{\DOCTITRE}{Les grands groupes humains}
% OnBuch+ — préambule « cours signés » ENRICHI (compilé avec tectonic / XeLaTeX)
% Système visuel « Pop » d'OnBuch+ : fond crème (jamais blanc), bordures encre
% épaisses, ombres « dures » décalées, titres Archivo Black en capitales.
% Version MATHÉMATIQUES : graphes (pgfplots/tikz), boîtes pédagogiques
% (définition, propriété, méthode, attention, le savais-tu, exercice résolu,
% exercice, TP, bilan), page de garde et signature OnBuch+ ancrée en bas.
%
% Macros attendues dans main.tex AVANT \input{preamble} :
%   \DOCMATIERE  \DOCNIVEAU  \DOCMODULE  \DOCLECON (numéro)  \DOCTITRE
\documentclass[11pt]{article}

\usepackage{fontspec}
\usepackage[french]{babel}
\usepackage[a4paper,margin=2cm,top=3.4cm,bottom=2.8cm,headheight=1.4cm,footskip=1.3cm]{geometry}
\usepackage{amsmath,amssymb,amsfonts}
\usepackage{mathtools} % \xmapsto, \coloneqq... souvent utilisés par les modèles
\usepackage{cancel} % simplifications barrées (bilans de forces)
% \abs{z} (module, valeur absolue) et \norm{u} (norme), utilisés par les modèles
\providecommand{\abs}{}\let\abs\relax\DeclarePairedDelimiter{\abs}{\lvert}{\rvert}
\providecommand{\norm}{}\let\norm\relax\DeclarePairedDelimiter{\norm}{\lVert}{\rVert}
\usepackage[table]{xcolor} % [table] : fournit aussi \rowcolors (alternance de lignes)
\usepackage{pagecolor}
\usepackage{tikz}
\usetikzlibrary{arrows.meta,positioning,calc,shapes.geometric,shapes.misc,shapes.arrows,decorations.pathmorphing,decorations.pathreplacing,decorations.markings,patterns,shadows,mindmap,trees,fit,backgrounds,matrix,intersections,angles,quotes}
\usepackage{pgfplots}
\usepackage[version=4]{mhchem} % équations nucléaires \ce{^{235}_{92}U}
\usepackage[european,siunitx]{circuitikz} % schémas électriques
\pgfplotsset{compat=1.18}
\usepgfplotslibrary{fillbetween,groupplots} % aires entre courbes (calcul intégral)
\usepackage{enumitem}
\usepackage{fancyhdr}
% [most] charge toutes les bibliothèques tcolorbox (listings, minted, raster,
% documentation...) et gonfle inutilement la mémoire TeX sur un document long
% (~300 boîtes) : on ne charge que ce qui est réellement utilisé.
\usepackage[skins,breakable]{tcolorbox}
\usepackage{graphicx}
\graphicspath{{/home/user/t-t/geo-onbuch/images/}}
\usepackage{colortbl}
\usepackage{booktabs}
\usepackage{tabularx}
\usepackage{multirow}
\usepackage{diagbox} % case d'en-tête diagonale des tableaux à double entrée
% Colonnes extensibles centrée / gauche / droite (C, L, R), souvent utilisées par les modèles
\newcolumntype{C}{>{\centering\arraybackslash}X}
\newcolumntype{L}{>{\raggedright\arraybackslash}X}
\newcolumntype{R}{>{\raggedleft\arraybackslash}X}
\usepackage{array}
\usepackage{multicol}
\usepackage{siunitx}
% parse-numbers=false : accepte aussi \SI{2,5 \times 10^{-3}}{...}
\sisetup{locale=FR,output-decimal-marker={,},group-minimum-digits=4,parse-numbers=false}
\DeclareSIUnit\ohm{\ensuremath{\symup{\Omega}}} % U+2126 absent de la police
\DeclareSIUnit\division{div} % réglages d'oscilloscope (ms/div, V/div)
\DeclareSIUnit\tr{tr} % tours (vitesse de rotation, ex. \SI{1500}{\tr\per\minute})
\DeclareSIUnit\molar{M} % concentration molaire (ex. \SI{3}{\molar}, \SI{1}{\micro\molar})
\DeclareSIUnit\an{an} % année (le modèle confond parfois avec \year, un compteur TeX) — ex. \SI{5}{\kilo\meter\per\an}
\DeclareSIUnit\FCFA{FCFA} % franc CFA (ex. \SI{500}{\FCFA\per\kilo\gram})
\DeclareSIUnit\degreeBrix{\textdegree Brix} % teneur en sucre d'un sirop/jus (réfractométrie)
\DeclareSIUnit\month{mois} % le modèle utilise parfois \month, un compteur TeX, comme unité de durée
\DeclareSIUnit\microliter{\micro\liter} % le modèle écrit parfois \microliter en un seul token
\DeclareSIUnit\million{millions} % dénombrement (ex. \SI{15}{\million\per\mL} spermatozoïdes)
\usepackage{qrcode}
% \degree : commande fréquemment utilisée par les modèles pour un angle ou une
% température en dehors de \SI{}{} (non fournie par les paquets chargés).
\providecommand{\degree}{^{\circ}}
% \qed (fin de démonstration), utilisé par les modèles sans amsthm
\providecommand{\qed}{\hfill$\square$}
% \barre : double barre des valeurs interdites dans un tableau de variations (tkz-tab)
\providecommand{\barre}{\ensuremath{\|}}
% environnement proof (amsthm non chargé) utilisé par les modèles
\ifdefined\proof\else\newenvironment{proof}[1][Démonstration]{\par\noindent\textit{#1.}\ }{\qed\par}\fi
\usepackage{lastpage}
\usepackage{hyperref}
\hypersetup{hidelinks}

% ============================================================
% Palette « Pop » OnBuch+ — mêmes hex que la classe P (plus_ui.dart)
% ============================================================
\definecolor{poporange}{HTML}{F59321}
\definecolor{popdark}{HTML}{E07A0C}
\definecolor{popcream}{HTML}{FFF6E8}
\definecolor{popink}{HTML}{1C1714}
\definecolor{poppurple}{HTML}{7A5AE0}
\definecolor{popgreen}{HTML}{2E9E6A}
\definecolor{poppink}{HTML}{E0457B}
\definecolor{popblue}{HTML}{2F7DE1}
\definecolor{popgold}{HTML}{C98A12}
\definecolor{popmuted}{HTML}{8A7F76}
\definecolor{popline}{HTML}{EADFD2}
\definecolor{popgray}{HTML}{8A7F76}
\colorlet{popgrey}{popgray}
% Alias tolérants (le modèle nomme parfois les couleurs différemment)
\colorlet{popwhite}{white}
\colorlet{popblack}{popink}
\colorlet{popred}{poppink}
\colorlet{popyellow}{popgold}
% Teintes claires (fonds de tableaux/encadrés) — le modèle nomme parfois
% les couleurs avec un suffixe L (ex. popblueL) sans les avoir définies.
\colorlet{poporangeL}{poporange!20}
\colorlet{popdarkL}{popdark!20}
\colorlet{poppurpleL}{poppurple!20}
\colorlet{popgreenL}{popgreen!20}
\colorlet{poppinkL}{poppink!20}
\colorlet{popblueL}{popblue!20}
\colorlet{popgoldL}{popgold!20}
\colorlet{popredL}{popred!20}
\colorlet{popgrayL}{popgray!20}
% Teintes claires pour fonds de tableaux / zones
\colorlet{popblueL}{popblue!10!white}
\colorlet{poporangeL}{poporange!14!white}
\colorlet{popgreenL}{popgreen!12!white}
\colorlet{poppurpleL}{poppurple!12!white}
\colorlet{poppinkL}{poppink!12!white}
\colorlet{popgoldL}{popgold!14!white}

\pagecolor{popcream}

% ============================================================
% Typographie
% ============================================================
\defaultfontfeatures{Ligatures=TeX}
\setmainfont{PlusJakartaSans-Medium.ttf}[Path=../../../fonts/,BoldFont=PlusJakartaSans-Bold.ttf]
% Maths en sans-serif (Fira Math) avec chiffres et lettres droites de Plus
% Jakarta, pour que les formules se fondent dans le texte.
\usepackage[math-style=ISO,bold-style=ISO]{unicode-math}
\setmathfont{FiraMath-Regular.otf}[Path=../../../fonts/]
\setmathfont{PlusJakartaSans-Medium.ttf}[Path=../../../fonts/,range={up/num,up/{Latin,latin},\mathup}]
\usepackage{icomma} % virgule décimale sans espace en mode math
% Caractères mathématiques Unicode absents des polices de texte (Plus Jakarta,
% Archivo Black) : rendus par leur équivalent mathématique, en texte comme en
% formule — sinon ils s'affichent en carré vide (ex. « Arithmétique dans ℤ »).
\usepackage{newunicodechar}
\newunicodechar{ℝ}{\ensuremath{\mathbb{R}}}
\newunicodechar{ℤ}{\ensuremath{\mathbb{Z}}}
\newunicodechar{ℂ}{\ensuremath{\mathbb{C}}}
\newunicodechar{ℕ}{\ensuremath{\mathbb{N}}}
\newunicodechar{ℚ}{\ensuremath{\mathbb{Q}}}
\newunicodechar{𝔻}{\ensuremath{\mathbb{D}}}
\newunicodechar{α}{\ensuremath{\alpha}}
\newunicodechar{β}{\ensuremath{\beta}}
\newunicodechar{γ}{\ensuremath{\gamma}}
\newunicodechar{δ}{\ensuremath{\delta}}
\newunicodechar{ε}{\ensuremath{\varepsilon}}
\newunicodechar{θ}{\ensuremath{\theta}}
\newunicodechar{λ}{\ensuremath{\lambda}}
\newunicodechar{μ}{\ensuremath{\mu}}
\newunicodechar{σ}{\ensuremath{\sigma}}
\newunicodechar{τ}{\ensuremath{\tau}}
\newunicodechar{φ}{\ensuremath{\varphi}}
\newunicodechar{ψ}{\ensuremath{\psi}}
\newunicodechar{ω}{\ensuremath{\omega}}
\newunicodechar{Σ}{\ensuremath{\Sigma}}
% majuscules produites par \MakeUppercase dans les titres (α-aminés -> Α)
\newunicodechar{Α}{\ensuremath{\alpha}}
\newunicodechar{Β}{\ensuremath{\beta}}
\newunicodechar{Γ}{\ensuremath{\Gamma}}
\newunicodechar{Δ}{\ensuremath{\Delta}}
\newunicodechar{Ε}{\ensuremath{\varepsilon}}
\newunicodechar{Θ}{\ensuremath{\Theta}}
\newunicodechar{Λ}{\ensuremath{\Lambda}}
\newunicodechar{Μ}{\ensuremath{\mu}}
\newunicodechar{Τ}{\ensuremath{\tau}}
\newunicodechar{Φ}{\ensuremath{\Phi}}
\newunicodechar{Ψ}{\ensuremath{\Psi}}
\newunicodechar{Ω}{\ensuremath{\Omega}}
\newunicodechar{↦}{\ensuremath{\mapsto}}
\newunicodechar{∈}{\ensuremath{\in}}
\newunicodechar{∉}{\ensuremath{\notin}}
\newunicodechar{∧}{\ensuremath{\wedge}}
\newunicodechar{⇔}{\ensuremath{\Leftrightarrow}}
\newunicodechar{⟺}{\ensuremath{\Longleftrightarrow}}
\newunicodechar{⇒}{\ensuremath{\Rightarrow}}
\newunicodechar{⟹}{\ensuremath{\Longrightarrow}}
\newunicodechar{∘}{\ensuremath{\circ}}
\newunicodechar{∪}{\ensuremath{\cup}}
\newunicodechar{∩}{\ensuremath{\cap}}
\newunicodechar{≡}{\ensuremath{\equiv}}
\newunicodechar{⊂}{\ensuremath{\subset}}
\newunicodechar{⊥}{\ensuremath{\perp}}
\newunicodechar{∃}{\ensuremath{\exists}}
\newunicodechar{∀}{\ensuremath{\forall}}
\newunicodechar{∛}{\ensuremath{\sqrt[3]{\,}}}
\newunicodechar{∜}{\ensuremath{\sqrt[4]{\,}}}
\newunicodechar{⃗}{\ensuremath{{}^{\to}}}
\newunicodechar{ⁿ}{\ifmmode^{n}\else\textsuperscript{\ensuremath{n}}\fi}
\newunicodechar{ˣ}{\ifmmode^{x}\else\textsuperscript{\ensuremath{x}}\fi}
\newunicodechar{ᵇ}{\ifmmode^{b}\else\textsuperscript{\ensuremath{b}}\fi}
\newunicodechar{ʳ}{\ifmmode^{r}\else\textsuperscript{\ensuremath{r}}\fi}
\newunicodechar{ᵘ}{\ifmmode^{u}\else\textsuperscript{\ensuremath{u}}\fi}
\newunicodechar{ʸ}{\ifmmode^{y}\else\textsuperscript{\ensuremath{y}}\fi}
\newunicodechar{ᵃ}{\ifmmode^{a}\else\textsuperscript{\ensuremath{a}}\fi}
\newunicodechar{ᵉ}{\ifmmode^{e}\else\textsuperscript{\ensuremath{e}}\fi}
\newunicodechar{ᵏ}{\ifmmode^{k}\else\textsuperscript{\ensuremath{k}}\fi}
\newunicodechar{ᵖ}{\ifmmode^{p}\else\textsuperscript{\ensuremath{p}}\fi}
\newunicodechar{ᴺ}{\ifmmode^{N}\else\textsuperscript{\ensuremath{N}}\fi}
\newunicodechar{ᶜ}{\ifmmode^{c}\else\textsuperscript{\ensuremath{c}}\fi}
\newunicodechar{ʰ}{\ifmmode^{h}\else\textsuperscript{\ensuremath{h}}\fi}
\newunicodechar{ᵠ}{\ifmmode^{q}\else\textsuperscript{\ensuremath{q}}\fi}
\newunicodechar{ₙ}{\ifmmode_{n}\else\textsubscript{\ensuremath{n}}\fi}
\newunicodechar{ₐ}{\ifmmode_{a}\else\textsubscript{\ensuremath{a}}\fi}
\newunicodechar{ᵢ}{\ifmmode_{i}\else\textsubscript{\ensuremath{i}}\fi}
\newunicodechar{ᵣ}{\ifmmode_{r}\else\textsubscript{\ensuremath{r}}\fi}
\newunicodechar{ₖ}{\ifmmode_{k}\else\textsubscript{\ensuremath{k}}\fi}
\newunicodechar{ᵦ}{\ifmmode_{\beta}\else\textsubscript{\ensuremath{\beta}}\fi}
\newunicodechar{ₓ}{\ifmmode_{x}\else\textsubscript{\ensuremath{x}}\fi}
\newfontfamily\popdisplay{ArchivoBlack-Regular.ttf}[Path=../../../fonts/]
\newfontfamily\poplogo{SpaceGrotesk-Bold.ttf}[Path=../../../fonts/]
\newfontfamily\popsemi{PlusJakartaSans-SemiBold.ttf}[Path=../../../fonts/]
\newfontfamily\popxbold{PlusJakartaSans-ExtraBold.ttf}[Path=../../../fonts/]
\newfontfamily\popxboldls{PlusJakartaSans-ExtraBold.ttf}[Path=../../../fonts/,LetterSpace=3.0]

\setlist[enumerate]{leftmargin=1.7em,itemsep=5pt,topsep=4pt}
\setlist[itemize]{leftmargin=1.7em,itemsep=4pt,topsep=4pt,label={\color{poporange}\textbullet}}
\setlength{\parindent}{0pt}
\setlength{\parskip}{5pt}
\renewcommand{\arraystretch}{1.25}

% pgfplots : style Pop par défaut
\pgfplotsset{
  every axis/.append style={axis line style={popink,line width=1pt},tick style={popink},
    grid style={popline,line width=0.5pt},label style={font=\small},tick label style={font=\footnotesize}},
  popcurve/.style={poporange,line width=1.8pt,no markers,smooth},
}
% Même style disponible dans un \draw TikZ simple (hors pgfplots)
\tikzset{popcurve/.style={poporange,line width=1.8pt}}
\tikzset{popgrid/.style={popline,very thin}}
\colorlet{popcurve}{poporange} % le nom du style est parfois utilisé comme couleur

% ============================================================
% Pill / squiggle
% ============================================================
\newcommand{\poppill}[3]{%
  \tikz[baseline=-0.55ex]\node[draw=popink,line width=1.2pt,rounded corners=2.6pt,%
    fill=#2,inner xsep=6.5pt,inner ysep=3pt]{%
    \popxboldls\fontsize{7}{7}\selectfont\color{#3}\MakeUppercase{#1}};%
}
\newcommand{\popsquiggle}[1]{%
  \tikz[baseline=-0.4ex]{
    \draw[#1,line width=2.6pt,line cap=round]
      plot [smooth] coordinates {(0,0.09) (0.34,0.01) (0.68,0.21) (1.02,0.01) (1.36,0.10)};
  }%
}
% Mot-clé surligné (marqueur orange sous le texte)
\newcommand{\cle}[1]{{\popxbold\color{popdark}#1}}

% ============================================================
% En-tête / pied
% ============================================================
\pagestyle{fancy}
\fancyhf{}
\renewcommand{\headrulewidth}{0pt}
\renewcommand{\footrulewidth}{0pt}
\lhead{\raisebox{-0.15ex}{\poplogo\fontsize{13}{13}\selectfont\color{popink}OnBuch\color{poporange}+}}
\rhead{\poppill{\DOCMATIERE\ \textbullet\ \DOCNIVEAU\ \textbullet\ Leçon \DOCLECON}{white}{popink}}
\lfoot{\popxbold\fontsize{7.6}{7.6}\selectfont\color{popmuted}\MakeUppercase{Cours signé OnBuch+}\ \textbullet\ {\popsemi\DOCTITRE}}
\rfoot{\tikz[baseline=-0.6ex]\node[rounded corners=8.5pt,fill=popink,minimum height=17pt,minimum width=17pt,inner xsep=5pt,inner sep=0]{\popxbold\fontsize{8}{8}\selectfont\color{popcream}\thepage\,/\,\pageref*{LastPage}};}

% ============================================================
% Titres
% ============================================================
\newcommand{\chaptitre}[2]{%
  \begin{tcolorbox}[enhanced,colback=poporange,colframe=popink,boxrule=2.6pt,arc=11pt,
      drop shadow=popink,left=15pt,right=15pt,top=12pt,bottom=11pt,before skip=3pt,after skip=18pt]
    \poppill{#2}{popink}{popcream}\par\vspace{9pt}
    {\raggedright\hyphenpenalty=10000\exhyphenpenalty=10000\popdisplay\fontsize{22}{24}\selectfont\color{popcream}\MakeUppercase{#1}\par}
  \end{tcolorbox}
}
% Section numérotée : pastille orange avec le numéro + titre Archivo
\newcounter{popsec}
\newcommand{\coursec}[1]{%
  \stepcounter{popsec}\setcounter{popsub}{0}%
  \par\vspace{16pt}\noindent
  \begin{minipage}{\linewidth}
  \tikz[baseline=-0.7ex]\node[draw=popink,line width=1.6pt,fill=poporange,rounded corners=4pt,minimum size=22pt,inner sep=0]{\popdisplay\fontsize{12}{12}\selectfont\color{popcream}\thepopsec};%
  \hspace{8pt}{\popdisplay\fontsize{15}{17}\selectfont\color{popink}\MakeUppercase{#1}}\par
  \vspace{4pt}\noindent{\color{poporange}\rule{36pt}{3.6pt}}
  \end{minipage}\par\nopagebreak\vspace{8pt}
}
\newcounter{popsub}
\newcommand{\courssub}[1]{%
  \stepcounter{popsub}%
  \par\vspace{9pt}\noindent{\popxbold\fontsize{11.5}{13}\selectfont\color{popink}{\color{poporange}\thepopsec.\thepopsub}\hspace{6pt}#1}\par\nopagebreak\vspace{4pt}
}

% ============================================================
% Boîtes pédagogiques — même recette : fond blanc, bordure épaisse colorée,
% ombre dure encre, badge plein. Titre optionnel [#1].
% ============================================================
\tcbset{popbase/.style={breakable,enhanced,colback=white,boxrule=2.3pt,arc=9pt,
  shadow={3pt}{-3pt}{0pt}{popink},left=14pt,right=14pt,top=11pt,bottom=11pt,before skip=11pt,after skip=15pt,
  fontupper=\popsemi\fontsize{10.6}{14.5}\selectfont\color{popink},
  fontlower=\popsemi\fontsize{10.6}{14.5}\selectfont\color{popink}}}
% Coche dessinée (le glyphe ✓ n'existe pas dans les polices du document)
\newcommand{\popcheck}[1]{\tikz[baseline=-0.5ex]\draw[#1,line width=1.6pt,line cap=round,line join=round](-0.13,0.0)--(-0.04,-0.09)--(0.13,0.1);}
\providecommand{\coche}{\popcheck{popgreen}}
\providecommand{\resultat}[1]{\par\smallskip\noindent\fcolorbox{popgreen}{popgreenL}{\parbox{\dimexpr\linewidth-2\fboxsep-2\fboxrule\relax}{\textbf{Résultat.} #1}}\par\smallskip}
\newcommand{\popboxhead}[3]{\poppill{#1}{#2}{white}\ifx\relax#3\relax\else\hspace{6pt}{\popxbold\fontsize{10.6}{12}\selectfont\color{popink}#3}\fi\par\vspace{7pt}}

\newtcolorbox{aretenir}[1][]{popbase,colframe=popgreen,before upper={\popboxhead{À retenir}{popgreen}{#1}}}
\newtcolorbox{exemplebox}[1][]{popbase,colframe=poporange,before upper={\popboxhead{Exemple}{poporange}{#1}}}
\newtcolorbox{definition}[1][]{popbase,colframe=popblue,before upper={\popboxhead{Définition}{popblue}{#1}}}
\newtcolorbox{propriete}[1][]{popbase,colframe=poppurple,before upper={\popboxhead{Propriété}{poppurple}{#1}}}
\newtcolorbox{methode}[1][]{popbase,colframe=popgold,before upper={\popboxhead{Méthode}{popgold}{#1}}}
\newtcolorbox{attention}[1][]{popbase,colframe=poppink,before upper={\popboxhead{Attention}{poppink}{#1}}}
\newtcolorbox{experience}[1][]{popbase,colframe=popblue,colback=popblueL,before upper={\popboxhead{Expérience}{popblue}{#1}}}
% « Le savais-tu ? » : carte violette pleine (contexte, culture, Cameroun)
\newtcolorbox{savaistu}[1][]{popbase,colback=poppurple,colframe=popink,
  fontupper=\popsemi\fontsize{10.4}{14.2}\selectfont\color{white},
  before upper={\poppill{Le savais-tu\,?}{white}{poppurple}\ifx\relax#1\relax\else\hspace{6pt}{\popxbold\color{white}#1}\fi\par\vspace{7pt}}}
% Exercice résolu : énoncé en haut, \tcblower puis la solution sur fond crème
\newcounter{popexo}\newcounter{popexr}
\newtcolorbox{exoresolu}[1][]{popbase,colframe=poporange,colbacklower=popcream,
  segmentation style={popink,line width=1.2pt,dashed},
  before upper={\stepcounter{popexr}\popboxhead{Exercice résolu \thepopexr}{poporange}{#1}},
  before lower={\poppill{Solution}{popink}{popcream}\par\vspace{6pt}}}
% Exercice (non résolu en place) — la correction est dans la partie Corrigés
\newtcolorbox{exercice}[1][]{popbase,colframe=popink,
  before upper={\stepcounter{popexo}\popboxhead{Exercice \thepopexo}{popink}{#1}}}
% Corrigé
\newtcolorbox{corrige}[1][]{popbase,colframe=popgreen,colback=popgreenL,
  before upper={\popboxhead{Corrigé}{popgreen}{#1}}}
% Objectifs / prérequis / bilan
\newtcolorbox{objectifs}{popbase,colframe=popink,colback=poporangeL,
  before upper={\popboxhead{Objectifs de la leçon}{popink}{}}}
\newtcolorbox{prerequis}{popbase,colframe=popink,colback=popblueL,
  before upper={\popboxhead{Prérequis}{popblue}{}}}
\newtcolorbox{bilan}{popbase,colframe=popink,colback=white,boxrule=2.8pt,
  before upper={\popboxhead{Fiche bilan}{poporange}{L'essentiel à connaître pour l'examen}}}
% Figure encadrée avec légende
\newtcolorbox{popfigure}[1][]{enhanced,breakable=false,colback=white,colframe=popline,boxrule=1.4pt,arc=8pt,
  left=10pt,right=10pt,top=10pt,bottom=8pt,before skip=10pt,after skip=12pt,halign=center,
  lower separated=false,halign lower=center,
  fontlower=\popsemi\fontsize{9}{11.5}\selectfont\color{popmuted},
  #1}
\newcommand{\legende}[1]{\par\vspace{6pt}{\popsemi\fontsize{9}{11.5}\selectfont\color{popmuted}#1}}

% ============================================================
% Signature OnBuch+
% ============================================================
\newcommand{\poplogoinline}[1]{{\poplogo\fontsize{#1}{#1}\selectfont\color{popink}OnBuch\color{poporange}+}}
\newcommand{\signatureonbuch}{%
  \begin{tcolorbox}[enhanced,colback=popink,colframe=popink,boxrule=2.2pt,arc=10pt,
      drop shadow=poporange,left=14pt,right=14pt,top=10pt,bottom=10pt,before skip=0pt,after skip=0pt]
    \tikz[baseline=-0.7ex]\node[circle,fill=popgreen,draw=popcream,line width=1.3pt,minimum size=22pt,inner sep=0]{\popcheck{white}};%
    \hspace{9pt}\begin{minipage}[c]{0.62\linewidth}
      {\popxbold\fontsize{12}{13}\selectfont\color{popcream}Signé\ \textbullet\ {\poplogo OnBuch\color{poporange}+}}\par\vspace{2pt}
      {\raggedright\popsemi\fontsize{8}{10.5}\selectfont\color{popline}\MakeUppercase{Équipe pédagogique OnBuch+}\\\MakeUppercase{Conforme au programme officiel MINESEC}\par}
    \end{minipage}\hfill
    {\poplogo\fontsize{22}{22}\selectfont\color{popcream}OnBuch\color{poporange}+}
  \end{tcolorbox}%
}

% ============================================================
% Page de garde
% #1 titre · #2 matière · #3 niveau · #4 module · #5 n° leçon · #6 durée ·
% #7 description / objectifs (texte ou itemize)
% ============================================================
\newcommand{\pagegarde}[7]{%
  \thispagestyle{empty}
  \newgeometry{a4paper,margin=1.7cm,top=1.6cm,bottom=1.4cm}%
  \begin{tikzpicture}[remember picture,overlay]
    \fill[poporange!18] ([xshift=-3.2cm,yshift=-3.6cm]current page.north east) circle (4.6cm);
    \fill[poppurple!14] ([xshift=2.4cm,yshift=7.6cm]current page.south west) circle (3.8cm);
  \end{tikzpicture}%
  {\poplogo\fontsize{30}{30}\selectfont\color{popink}OnBuch\color{poporange}+}\hfill
  \raisebox{0.6ex}{\poppill{Cours signé \textbullet\ Premium}{popink}{popcream}}\par
  \vspace{1pt}\popsquiggle{poppurple}\par
  \vspace{8pt}
  \begin{tcolorbox}[enhanced,colback=poporange,colframe=popink,boxrule=2.8pt,arc=13pt,
      drop shadow=popink,left=18pt,right=18pt,top=12pt,bottom=13pt,after skip=10pt]
    \poppill{#2 \textbullet\ #3}{popink}{popcream}\hspace{5pt}\poppill{#4}{white}{popink}\par\vspace{8pt}
    {\popdisplay\fontsize{14}{14}\selectfont\color{popink}LEÇON #5}\par\vspace{4pt}
    {\raggedright\hyphenpenalty=10000\exhyphenpenalty=10000\popdisplay\fontsize{25}{27}\selectfont\color{popcream}\MakeUppercase{#1}\par}
    \vspace{8pt}
    {\popxbold\fontsize{9.5}{9.5}\selectfont\color{popink}\MakeUppercase{Durée conseillée : #6}}
  \end{tcolorbox}
  \begin{tcolorbox}[enhanced,colback=white,colframe=popink,boxrule=2.2pt,arc=10pt,
      drop shadow=popink,left=16pt,right=16pt,top=10pt,bottom=10pt,after skip=8pt]
    \poppill{Dans cette leçon}{poppurple}{white}\par\vspace{5pt}
    \popsemi\fontsize{9.3}{12.6}\selectfont\color{popink}#7
  \end{tcolorbox}
  \noindent\begin{minipage}[t]{0.487\linewidth}
    \begin{tcolorbox}[enhanced,colback=popgreen,colframe=popink,boxrule=2.2pt,arc=10pt,
        drop shadow=popink,left=11pt,right=11pt,top=10pt,bottom=10pt,height=3.1cm,valign=center]
      \tcbox[colback=white,colframe=popink,boxrule=1.3pt,arc=5pt,boxsep=0pt,nobeforeafter,
        left=3pt,right=3pt,top=3pt,bottom=3pt,valign=center]{\qrcode[height=1.9cm]{https://play.google.com/store/apps/details?id=com.luvvix.onbuch&hl=fr}}%
      \hspace{8pt}\begin{minipage}[c]{0.5\linewidth}
        {\popdisplay\fontsize{11}{12.5}\selectfont\color{white}\MakeUppercase{Télécharge OnBuch}}\par\vspace{4pt}
        {\raggedright\popsemi\fontsize{8.4}{11}\selectfont\color{white}Gratuit \textbullet\ Google Play\\Tuteur IA, annales, concours, résultats\par}
      \end{minipage}
    \end{tcolorbox}
  \end{minipage}\hfill
  \begin{minipage}[t]{0.487\linewidth}
    \begin{tcolorbox}[enhanced,colback=poppurple,colframe=popink,boxrule=2.2pt,arc=10pt,
        drop shadow=popink,left=11pt,right=11pt,top=10pt,bottom=10pt,height=3.1cm,valign=center]
      \tikz[baseline=-0.7ex]\node[circle,fill=white,draw=popink,line width=1.3pt,minimum size=1.6cm,inner sep=0]{\popdisplay\fontsize{24}{24}\selectfont\color{poppurple}+};%
      \hspace{8pt}\begin{minipage}[c]{0.6\linewidth}
        {\popdisplay\fontsize{11}{12.5}\selectfont\color{white}\MakeUppercase{Passe à OnBuch+}}\par\vspace{4pt}
        {\raggedright\popsemi\fontsize{8.4}{11}\selectfont\color{white}Tous les cours signés, Tuteur IA illimité, fiches PDF — onglet OnBuch+ de l'app\par}
      \end{minipage}
    \end{tcolorbox}
  \end{minipage}
  \vspace{8pt}
  \popcontenu
  \vfill
  \signatureonbuch
  \restoregeometry
  \newpage
}

% Rangée « Ce cours contient » (page de garde)
\newcommand{\poptile}[3]{% #1 couleur #2 gros texte #3 légende
  \begin{tcolorbox}[enhanced,colback=#1,colframe=popink,boxrule=2pt,arc=9pt,shadow={2.5pt}{-2.5pt}{0pt}{popink},
      left=6pt,right=6pt,top=9pt,bottom=9pt,halign=center,valign=center,height=1.75cm,width=\linewidth,nobeforeafter]
    {\popdisplay\fontsize{12.5}{13}\selectfont\color{white}\MakeUppercase{#2}}\par\vspace{3pt}
    {\centering\popsemi\fontsize{7.8}{9.5}\selectfont\color{white}#3\par}
  \end{tcolorbox}}
\newcommand{\popcontenu}{%
  \par\noindent{\popxboldls\fontsize{7.5}{7.5}\selectfont\color{popmuted}CE COURS SIGNÉ CONTIENT}\par\vspace{7pt}
  \noindent\begin{minipage}[t]{0.235\linewidth}\poptile{popblue}{Cours}{complet, illustré, pas à pas}\end{minipage}\hfill
  \begin{minipage}[t]{0.235\linewidth}\poptile{poporange}{Méthodes}{et exercices résolus}\end{minipage}\hfill
  \begin{minipage}[t]{0.235\linewidth}\poptile{poppink}{Exercices}{type examen + corrigés}\end{minipage}\hfill
  \begin{minipage}[t]{0.235\linewidth}\poptile{popgreen}{Fiche bilan}{l'essentiel à retenir}\end{minipage}\par}

% Sommaire
\newtcolorbox{sommaire}{enhanced,colback=white,colframe=popink,boxrule=2.2pt,arc=10pt,
  drop shadow=popink,left=18pt,right=18pt,top=13pt,bottom=13pt,before skip=6pt,after skip=20pt,
  before upper={\poppill{Sommaire}{poppurple}{white}\par\vspace{9pt}\popsemi\fontsize{10.6}{14.5}\selectfont\color{popink}}}

% Signature de fin de cours — poussée tout en bas de la dernière page
\newcommand{\findecours}{%
  \par\vfill\nopagebreak
  \begin{center}\popsquiggle{poporange}\end{center}\vspace{4pt}
  \signatureonbuch
}
\AtBeginDocument{\def\llbracket{\mathopen{[\mkern-3.5mu[}}\def\rrbracket{\mathclose{]\mkern-3.5mu]}}} % intervalles d'entiers (glyphe ⟦ absent de la police)
% Glyphes absents de Fira Math : redéfinis à partir de glyphes disponibles
\AtBeginDocument{%
  \def\setminus{\mathbin{\backslash}}\let\smallsetminus\setminus
  \def\vdots{\mathord{\vbox{\baselineskip3.2pt\lineskiplimit0pt\kern4pt\hbox{.}\hbox{.}\hbox{.}}}}%
  \def\ddots{\mathinner{\mkern1mu\raise6.5pt\hbox{.}\mkern2mu\raise3.6pt\hbox{.}\mkern2mu\raise0.7pt\hbox{.}\mkern1mu}}%
  \def\triangle{\mathord{\scalebox{0.9}{$\symup{\Delta}$}}}%
  \def\nabla{\mathord{\raisebox{\depth}{\scalebox{1}[-1]{$\symup{\Delta}$}}}}%
  \def\checkmark{\popcheck{popdark}}%
  \def\star{\mathbin{\ast}}\def\bigstar{\mathord{\ast}}%
  \def\diamond{\mathbin{\scalebox{0.6}{$\Diamond$}}}%
  \def\bigcup{\mathop{\mathchoice{\raisebox{-0.3em}{\scalebox{1.7}{$\cup$}}}{\scalebox{1.25}{$\cup$}}{\cup}{\cup}}\displaylimits}%
  \def\bigcap{\mathop{\mathchoice{\raisebox{-0.3em}{\scalebox{1.7}{$\cap$}}}{\scalebox{1.25}{$\cap$}}{\cap}{\cap}}\displaylimits}%
  \def\lesseqqgtr{\mathrel{\lessgtr}}\def\bigoplus{\mathop{\oplus}\displaylimits}%
}
\providecommand{\ssi}{si et seulement si} % abréviation fréquente
\AtBeginDocument{\providecommand{\wideparen}{\overparen}} % arc de cercle

\begin{document}

\pagegarde{Les grands groupes humains}{Géographie}{Tle E}{Module 1 : Le Cameroun}{N02}{2 h + TP 2 (2 h)}{Cette leçon présente les grands groupes humains du Cameroun (soudanais, bantou, semi-bantou, pygmées), leurs localisations, leurs organisations socio-culturelles et religieuses, ainsi que l'inégale répartition de la population sur le territoire national. Elle s'achève par un travail pratique de construction et de lecture de la carte des densités de population.
\vspace{4pt}
\begin{multicols}{2}\small
\begin{itemize}
\item À la fin de cette leçon, je sais localiser sur une carte les grands groupes humains du Cameroun (soudanais, bantou, semi-bantou, pygmées)
\item À la fin de cette leçon, je sais distinguer les paléo-soudanais des néo-soudanais et situer le peuplement bantou
\item À la fin de cette leçon, je sais décrire les organisations socio-culturelles et religieuses propres à chaque grand groupe humain
\item À la fin de cette leçon, je sais expliquer l'inégale répartition de la population camerounaise à partir des différences de densités
\item À la fin de cette leçon, je sais analyser la grande mobilité et le brassage des populations comme facteurs d'intégration nationale
\item À la fin de cette leçon, je sais observer, localiser et décrire des phénomènes géographiques à partir de cartes, photographies et tableaux statistiques
\end{itemize}\end{multicols}}

\begin{sommaire}
\begin{multicols}{2}
\begin{enumerate}
\item Introduction : définitions et cadre d'étude
\item Les groupes soudanais : paléo-soudanais et néo-soudanais
\item Les groupes bantou et semi-bantou ; les pygmées
\item L'inégale répartition de la population : les différences de densités
\item Mobilité, brassage des populations et intégration nationale
\item TP 2 : Construction et lecture de la carte de répartition de la population
\item Activité d'intégration
\item Méthodes et exercices résolus
\item Exercices
\item Corrigés des exercices
\item Fiche bilan
\end{enumerate}
\end{multicols}
\end{sommaire}

\begin{objectifs}
\begin{itemize}
\item À la fin de cette leçon, je sais localiser sur une carte les grands groupes humains du Cameroun (soudanais, bantou, semi-bantou, pygmées)
\item À la fin de cette leçon, je sais distinguer les paléo-soudanais des néo-soudanais et situer le peuplement bantou
\item À la fin de cette leçon, je sais décrire les organisations socio-culturelles et religieuses propres à chaque grand groupe humain
\item À la fin de cette leçon, je sais expliquer l'inégale répartition de la population camerounaise à partir des différences de densités
\item À la fin de cette leçon, je sais analyser la grande mobilité et le brassage des populations comme facteurs d'intégration nationale
\item À la fin de cette leçon, je sais observer, localiser et décrire des phénomènes géographiques à partir de cartes, photographies et tableaux statistiques
\item À la fin de cette leçon, je sais construire et légender une carte de répartition de la population (zones densément, moyennement et faiblement peuplées)
\item À la fin de cette leçon, je sais construire un diagramme et un croquis simples pour représenter des données démographiques
\end{itemize}
\end{objectifs}

\begin{prerequis}
\begin{itemize}
\item Notion de population, densité de population et croissance démographique (classe de Première)
\item Localisation des dix régions du Cameroun et de leurs chefs-lieux
\item Notions de migration, d'exode rural et d'urbanisation
\item Lecture élémentaire d'une carte (échelle, légende, orientation)
\item Grandes divisions physiques du Cameroun (relief, climats) vues en début d'année
\end{itemize}
\end{prerequis}

\newpage

\coursec{Introduction : définitions et cadre d'étude}

\courssub{Situation d'accroche : le grand départ vers les villages}

Chaque mois de décembre, le même phénomène se répète au Cameroun : les gares routières de Douala (Bépanda, Bonabéri) et de Yaoundé (Mvan, Nsam) débordent de voyageurs chargés de bagages et de cadeaux. Les cars des agences de voyage partent pleins vers Bafoussam, Ebolowa, Maroua, Bertoua ou Ngaoundéré. Bamiléké, Beti, Fulbé, Bassa, Ewondo, Kotoko, Mafa... chacun \og retourne au village \fg{} pour les fêtes de fin d'année. Pourtant, ces mêmes familles vivent toute l'année dans des quartiers cosmopolites comme New-Bell à Douala ou Briqueterie à Yaoundé, où l'on entend, dans la même rue, cinq ou six langues différentes, où le voisin de palier vient d'une autre région, où les mariages croisent les origines.

Cette observation simple, que tout élève camerounais a faite, pose deux grandes questions géographiques. D'une part, comment un pays comptant plus de 250 groupes ethniques, héritiers de grandes vagues de migrations venues du nord (migrations soudanaises) et du sud (expansion bantoue), parvient-il à construire une nation unie ? D'autre part, pourquoi la population se répartit-elle si inégalement : pourquoi le Littoral et l'Ouest concentrent-ils les hommes tandis que l'Est ou certaines parties de l'Extrême-Nord restent faiblement peuplés ?

\begin{aretenir}[Problématique de la leçon]
Comment la diversité humaine du Cameroun et la répartition inégale de sa population influencent-elles l'intégration nationale ?
\end{aretenir}

Pour répondre à cette problématique, la leçon suivra le plan suivant : après cette introduction qui définit les notions et le cadre d'étude, nous étudierons les groupes soudanais (paléo-soudanais et néo-soudanais), puis les groupes bantou, semi-bantou et les pygmées, ensuite l'inégale répartition de la population et ses différences de densités, enfin la mobilité, le brassage des populations et l'intégration nationale. La leçon se clôturera par le TP 2 : construction et lecture de la carte de répartition de la population.

\courssub{Qu'est-ce qu'un groupe humain ? Définition et premiers repères}

\textbf{L'intuition d'abord.} Quand tu dis \og je suis Maka \fg{}, \og je suis Bamoun \fg{} ou \og je suis Douala \fg{}, tu affirmes une appartenance : tu partages avec d'autres personnes une langue, des coutumes (mariage, funérailles, danses, cuisine), une histoire commune (souvent racontée sous forme de récits d'origine et de migrations) et, le plus souvent, un territoire d'attache, le \og village \fg{}. C'est exactement ce que le géographe appelle un groupe humain.

\begin{definition}[Groupe humain (ou ethnie)]
Un \cle{groupe humain} (ou \cle{ethnie}) est un ensemble de personnes qui partagent une \textbf{langue}, une \textbf{culture} (coutumes, croyances, arts, modes de vie), une \textbf{histoire} commune (origine, migrations) et, le plus souvent, un \textbf{territoire} commun. L'appartenance à un groupe humain est culturelle et historique : elle se transmet par l'éducation, la langue et la vie sociale.
\end{definition}

Le Cameroun compte \textbf{plus de 250 groupes ethniques}, ce qui en fait l'un des pays les plus diversifiés d'Afrique. Cette diversité n'est pas un désordre : elle s'explique par la position géographique du pays et par son histoire, faite de grandes migrations successives.

\begin{attention}[Ne pas confondre groupe ethnique et race]
Le mot \og \cle{race} \fg{} n'a aucune valeur scientifique pour classer les êtres humains : la biologie a montré que l'espèce humaine ne se divise pas en races. La géographie étudie des \textbf{groupes culturels et historiques}, définis par la langue, les coutumes et l'histoire, et non des catégories biologiques fondées sur l'apparence physique. Dire \og les Bamiléké sont une ethnie \fg{} est correct ; parler de \og race bamiléké \fg{} est une erreur grave, à la fois scientifique et de rédaction. Au Baccalauréat, cette confusion coûte des points.
\end{attention}

\courssub{Le Cameroun, \og Afrique en miniature \fg{} : un carrefour humain}

\textbf{Observation.} Regarde une carte d'Afrique : le Cameroun se situe exactement à la charnière entre l'Afrique de l'Ouest et l'Afrique centrale, entre le monde sahélien (savanes sèches, islam, grands empires historiques) et le monde forestier équatorial (forêt dense, agriculture sur brûlis, sociétés villageoises). Le pays s'étire sur plus de \num{1200} km du lac Tchad au golfe de Guinée, et traverse presque tous les milieux naturels africains : sahel, savane, hautes terres, forêt dense, littoral.

\textbf{Conséquence humaine.} Ce carrefour naturel a été aussi un carrefour humain. Pendant des siècles, des peuples ont migré vers et à travers le territoire camerounais :

\begin{itemize}
\item depuis le \textbf{nord-est} (région du Tchad), des populations dites \emph{soudanaises} se sont installées dans les savanes du Nord-Cameroun ;
\item depuis le \textbf{nord-ouest} (actuels Nigeria et plateau de l'Adamaoua), les \emph{Peuls (Fulbé)} ont progressé avec leurs troupeaux, notamment lors du jihad du XIX\textsuperscript{e} siècle ;
\item depuis le \textbf{sud et le sud-est} (bassin du Congo), le vaste mouvement d'\emph{expansion bantoue}, amorcé il y a environ \num{2000} à \num{3000} ans (ordre de grandeur admis par les historiens), a occupé la zone forestière et les hautes terres de l'Ouest ;
\item les \emph{pygmées}, considérés comme les premiers habitants de la forêt équatoriale, ont été progressivement refoulés vers les massifs forestiers du Sud et de l'Est.
\end{itemize}

C'est pourquoi on dit que le Cameroun est une \og \cle{Afrique en miniature} \fg{} : on y trouve, concentrés sur \SI{475442}{km^{2}}, presque tous les milieux naturels et presque tous les grands types de sociétés du continent.

\begin{savaistu}[Un record africain de diversité linguistique]
Avec \textbf{plus de 250 langues} parlées sur son territoire (certains recensements en dénombrent près de 280), le Cameroun est l'un des pays les plus diversifiés linguistiquement d'Afrique, aux côtés du Nigeria et de la République démocratique du Congo. Le français et l'anglais, langues officielles héritées de la colonisation, ainsi que des langues véhiculaires comme le fulfuldé au Nord, l'ewondo autour de Yaoundé ou le pidgin english dans les régions anglophones, servent de ponts de communication entre les groupes. C'est un enjeu majeur d'intégration nationale.
\end{savaistu}

\courssub{Les critères de classement des grands groupes humains}

Comment le géographe s'y prend-il pour ranger plus de 250 ethnies en quelques grands ensembles lisibles ? Il croise plusieurs critères, qu'il faut connaître précisément.

\begin{definition}[Critères de classement des groupes humains]
Les groupes humains du Cameroun sont classés selon quatre critères combinés :
\begin{itemize}
\item l'\cle{aire géographique} : la région d'implantation dominante (savanes du Nord, hautes terres de l'Ouest, forêt du Centre-Sud-Est, littoral) ;
\item la \cle{langue} : l'appartenance à une même famille linguistique (langues tchadiques, adamawa-oubanguiennes, bantoues, etc.) ;
\item l'\cle{histoire migratoire} : l'origine commune et les grands mouvements de peuplement (migrations soudanaises, expansion bantoue, avancée peule) ;
\item les \cle{modes de vie} : organisation sociale et politique (chefferies centralisées ou sociétés villageoises), activités (agriculture, élevage, cueillette), religion (islam, christianisme, religions traditionnelles).
\end{itemize}
\end{definition}

\begin{methode}[Étudier un groupe humain en géographie : les trois questions]
Pour décrire n'importe quel groupe humain du Cameroun (ou d'ailleurs), réponds toujours, dans cet ordre, à trois questions :
\begin{enumerate}
\item \textbf{Où vit-il ?} Localise son aire géographique (région, milieu naturel, villes principales) et représente-la sur une carte ou un croquis.
\item \textbf{Comment est-il organisé ?} Décris son organisation sociale, politique et religieuse (chefferie, lignages, sultanat, associations, cultes).
\item \textbf{Quelles sont ses activités ?} Précise ses modes de vie et ses activités économiques (cultures, élevage, artisanat, commerce, migrations de travail).
\end{enumerate}
Cette démarche en trois questions sera appliquée dans les sections suivantes à chaque grand groupe étudié : soudanais, bantou et semi-bantou, pygmées.
\end{methode}

\courssub{Les trois grands ensembles humains du Cameroun}

En croisant ces critères, les géographes regroupent les ethnies camerounaises en \textbf{trois grands ensembles}.

\textbf{1. Les groupes soudanais.} Ils occupent les savanes des régions du \textbf{Nord, de l'Extrême-Nord et de l'Adamaoua}. On distingue les \emph{paléo-soudanais} (anciens habitants des savanes et zones de refuge montagneuses : Mafa, Mofu, Kotoko, Massa, Toupouri, Mundang...) et les \emph{néo-soudanais} (arrivés plus récemment, souvent islamisés, occupants les plaines : Fulbé/Peuls, Arabes Choa, Kanouri, Haoussa...). Leurs sociétés sont souvent organisées en chefferies ou sultanats hiérarchisés (lamidats), avec l'élevage et l'agriculture de savane (mil, sorgho, maïs, coton) comme activités de base.

\textbf{2. Les groupes bantou et semi-bantou.} Issus de la grande expansion bantoue, ils occupent le \textbf{Centre, le Sud, le Littoral, l'Est, le Nord-Ouest, le Sud-Ouest et l'Ouest}. Les \emph{bantou} proprement dits (Beti : Ewondo, Bulu, Eton ; Bassa, Douala, Bakweri, Maka, Kako...) vivent dans la zone forestière et littorale ; les \emph{semi-bantou} (Bamiléké, Bamoun, Tikari, Kom, Bafut, Nso...) occupent les hautes terres de l'Ouest et du Nord-Ouest (Grassfields), à la charnière entre forêt et savane, avec de puissantes chefferies centralisées (fonctions, sultanats).

\textbf{3. Les pygmées.} Premiers habitants de la forêt équatoriale, ils vivent aujourd'hui dans les massifs forestiers du \textbf{Sud et de l'Est} (Baka à l'Est, Bakola/Bagyeli au Sud, Bedzang au Centre), en petits groupes mobiles traditionnellement fondés sur la chasse, la cueillette et la pêche, en relation d'échange avec les villages bantou voisins.

\begin{exemplebox}[Appliquer la méthode des trois questions aux Baka]
\begin{itemize}
\item \textbf{Où ?} Forêts denses de l'Est (environs de Lomié, Moloundou, Yokadouma) et du Sud (Campo, Ma'an).
\item \textbf{Comment organisés ?} Petits campements familiaux mobiles (mongulu), société égalitaire sans chefferie centralisée, forte connaissance du milieu forestier, rites d'initiation (Jengi).
\item \textbf{Quelles activités ?} Chasse au filet/arc, cueillette (miel, fruits, champignons), pêche, échanges de produits forestiers contre produits agricoles (manioc, plantain) avec les villages Bantu voisins, de plus en plus agriculture d'appoint et travail salarié (guides, pisteurs).
\end{itemize}
\end{exemplebox}

\courssub{Cartographie et classification : les deux outils du géographe}

Pour étudier ces groupes, le géographe dispose de deux représentations fondamentales : la \textbf{carte} (où sont les groupes ?) et le \textbf{schéma de classification} (comment se rattachent-ils les uns aux autres ?). La carte muette ci-dessous servira de support tout au long de la leçon ; l'arbre de classification résume l'organisation générale du peuplement.

\begin{popfigure}
\begin{center}
\includegraphics[width=\linewidth,height=11cm,keepaspectratio]{N02/cmr_reg_muette.jpg}
\end{center}
\legende{Carte muette du Cameroun (limites des régions en trait épais, départements en trait fin, fleuves et lac Tchad en bleu). Travail à faire en TP : placer et nommer les dix régions (Adamaoua, Centre, Est, Extrême-Nord, Littoral, Nord, Nord-Ouest, Ouest, Sud, Sud-Ouest), repérer les grandes villes, puis hachurer les aires des grands groupes humains. \\ Source du fond de carte : Wikimedia Commons, Flappiefh, CC BY-SA 4.0.}
\end{popfigure}

\begin{popfigure}
\begin{center}
\begin{tikzpicture}[
  grow=down,
  level 1/.style={sibling distance=5.6cm, level distance=1.6cm},
  level 2/.style={sibling distance=2.9cm, level distance=1.7cm},
  every node/.style={draw, rounded corners=2pt, align=center, font=\scriptsize, fill=popcream, inner sep=3pt},
  edge from parent/.style={draw=popmuted, thick}
]
\node[fill=popblueL, font=\scriptsize\bfseries, align=center]{Population du Cameroun\\ (plus de 250 ethnies)}
  child { node[fill=popgoldL, align=center]{\textbf{Groupes soudanais}\\ Nord, Extrême-Nord,\\ Adamaoua}
    child { node[align=center]{Paléo-soudanais\\ Mafa, Mofu,\\ Kotoko, Massa...} }
    child { node[align=center]{Néo-soudanais\\ Fulbé (Peuls),\\ Arabes Choa...} }
  }
  child { node[fill=popgreenL, align=center]{\textbf{Bantou et semi-bantou}\\ Centre, Sud, Littoral,\\ Est, Ouest, NO, SO}
    child { node[align=center]{Bantou\\ Beti, Bulu, Bassa,\\ Douala, Bakweri...} }
    child { node[align=center]{Semi-bantou\\ Bamiléké, Bamoun,\\ Tikari, Kom...} }
  }
  child { node[fill=poppinkL, align=center]{\textbf{Pygmées}\\ Forêts du Sud\\ et de l'Est}
    child { node[align=center]{Baka, Bakola,\\ Bagyeli...} }
  };
\end{tikzpicture}
\end{center}
\legende{Arbre de classification des grands groupes humains du Cameroun. Le classement croise quatre critères : aire géographique, famille linguistique, histoire migratoire et modes de vie.}
\end{popfigure}

\courssub{Exercice résolu et premières mises en garde}

\begin{exoresolu}[Lire un tableau statistique : la diversité en chiffres]
\textbf{Énoncé.} Voici des données simplifiées sur le peuplement camerounais (ordres de grandeur) : population totale environ 28 millions d'habitants (estimations récentes) ; plus de 250 ethnies ; trois grandes familles humaines (soudanaise, bantou et semi-bantou, pygmée) ; les pygmées représentent moins de 1 \% de la population.
1) Calcule le nombre moyen d'ethnies par million d'habitants.
2) Estime le nombre maximal de Pygmées au Cameroun.
3) Que révèle la comparaison de ces deux résultats ?
\tcblower
\textbf{Solution détaillée.}
1) Nombre moyen d'ethnies par million d'habitants :
\[ \frac{250}{28} \approx 8,9 \text{ ethnies par million d'habitants.} \]
C'est un chiffre très élevé, qui traduit une mosaïque humaine exceptionnelle.
2) Les Pygmées représentent moins de 1 \% de 28 millions :
\[ \frac{1}{100} \times 28\,000\,000 = 280\,000 \text{ personnes.} \]
Les Pygmées sont donc moins de \num{280000} (les estimations courantes donnent plutôt entre \num{50000} et \num{100000} personnes, ordre de grandeur à retenir avec prudence).
3) \textbf{Interprétation :} le Cameroun combine une extrême diversité (environ 9 ethnies par million d'habitants) et une forte inégalité numérique entre les groupes : certains comptent plusieurs millions de membres (Bamiléké, Beti, Fulbé), d'autres quelques milliers seulement (Pygmées). Cette diversité inégalement répartie est au cœur de la problématique de l'intégration nationale : faire vivre ensemble, sur un même territoire, des groupes très différents par la taille, la langue et les modes de vie.
\end{exoresolu}

\begin{attention}[Pièges fréquents dans cette introduction]
\begin{itemize}
\item \textbf{Confondre ethnie et race} : on l'a dit, la géographie classe des cultures, pas des corps.
\item \textbf{Donner un chiffre faux} : retiens \og plus de 250 ethnies \fg{} et \og plus de 250 langues \fg{} ; évite les chiffres précis inventés.
\item \textbf{Croire que les groupes sont figés} : les aires des groupes se chevauchent et les migrations brassent en permanence les populations ; une carte des ethnies est toujours une simplification.
\item \textbf{Oublier les semi-bantou} : ils forment un ensemble à part entière, à la charnière des mondes bantou et soudanais ; ne les fonds pas dans les bantou sans explication.
\end{itemize}
\end{attention}

\begin{aretenir}[L'essentiel de l'introduction]
\begin{itemize}
\item Un \cle{groupe humain} (ethnie) partage langue, culture, histoire et souvent territoire.
\item Le Cameroun, \og \cle{Afrique en miniature} \fg{}, est un carrefour entre l'Afrique soudano-sahélienne et l'Afrique centrale bantoue : plus de 250 ethnies et plus de 250 langues.
\item On classe les groupes selon quatre critères : \cle{aire géographique}, \cle{langue}, \cle{histoire migratoire}, \cle{modes de vie}.
\item Trois grands ensembles : les \cle{soudanais} (paléo et néo-soudanais) au Nord, les \cle{bantou et semi-bantou} au Centre-Sud-Ouest, les \cle{pygmées} dans les forêts du Sud-Est.
\item Problématique directrice : comment cette diversité et cette répartition inégale influencent-elles l'\cle{intégration nationale} ?
\end{itemize}
\end{aretenir}

Nous pouvons maintenant entrer dans le vif du sujet en remontant vers les savanes du Nord : qui sont les groupes soudanais, paléo-soudanais et néo-soudanais, et comment se sont-ils installés au Cameroun ? C'est l'objet de la section suivante.

\coursec{Les grands groupes humains : localisation, organisations et répartition}

\courssub{Introduction : définitions et cadre d'étude}

Le Cameroun, souvent qualifié d'\textbf{« Afrique en miniature »} par sa diversité géographique, l'est tout autant par sa mosaïque humaine. Avec une population estimée à \textbf{28,6 millions d'habitants en 2023} (INS, projections RGPH 2005), le pays compte plus de \textbf{250 groupes ethniques} recensés, regroupés en grands ensembles linguistiques et historiques. Cette leçon s'inscrit dans la \textbf{Famille de situations : L'intégration nationale}. Comprendre qui sont les Camerounais, où ils vivent, comment ils s'organisent et se déplacent, c'est saisir les fondements et les défis de la construction de l'unité nationale.

\begin{definition}[Concepts clés : Peuple, Ethnie, Groupe, Population]
\begin{itemize}
    \item \textbf{Peuple :} Ensemble d'individus unis par une histoire commune, une culture, une langue, une conscience d'appartenance et un territoire (dimension politique).
    \item \textbf{Ethnie (ou groupe ethnique) :} Ensemble d'individus partageant une langue, une culture, des traditions, une origine mythique commune, sans nécessairement former un État (dimension socio-culturelle). Au Cameroun, on préfère souvent le terme \textbf{groupe ethnique} ou \textbf{groupe humain}.
    \item \textbf{Population :} Ensemble des habitants d'un territoire donné, approche quantitative (démographie) et qualitative (structure par âge, sexe, activité).
\end{itemize}
\end{definition}

\begin{aretenir}[Cadre spatial et démographique national]
\begin{itemize}
    \item \textbf{Superficie :} \num{475 442} km².
    \item \textbf{Population (2023) :} \textasciitilde 28,6 millions d'hab. (Densité moyenne : \textasciitilde 60 hab/km²).
    \item \textbf{Taux d'accroissement :} \textasciitilde 2,6 \% / an (forte croissance).
    \item \textbf{Structure par âge :} Population jeune (\textasciitilde 60 \% < 25 ans).
    \item \textbf{Urbanisation :} \textasciitilde 58 \% (seuil franchi vers 2015), forte primatie de Douala (\textasciitilde 3,9 M) et Yaoundé (\textasciitilde 4,1 M).
    \item \textbf{Langues officielles :} Français, Anglais. \textbf{Langues véhiculaires :} Fulfulde (Nord), Ewondo/Beti (Centre-Sud-Est), Duala/Bassa (Littoral), Pidgin English (Nord-Ouest/Sud-Ouest).
\end{itemize}
\end{aretenir}

\noindent Le programme distingue trois grands ensembles géo-historiques : les \textbf{Groupes soudanais} (Grand Nord), les \textbf{Groupes bantou et semi-bantou} (Sud et Ouest), et les \textbf{Pygmées} (forêts du Sud-Est). Nous commençons par les groupes soudanais.

\coursec{Les groupes soudanais : paléo-soudanais et néo-soudanais}

\courssub{Localisation générale : le Grand Nord camerounais}

L'ensemble des groupes soudanais occupe les trois régions septentrionales du Cameroun : le \textbf{Nord}, l'\textbf{Adamaoua} et l'\textbf{Extrême-Nord}. Ce vaste territoire, qui s'étend du 7\textsuperscript{e} parallèle nord (frontière naturelle avec la zone de transition soudanais-guinéenne) jusqu'au lac Tchad au nord, forme une unité géographique et historique forte : le \textbf{Grand Nord}.

\begin{popfigure}
\begin{center}
\includegraphics[width=\linewidth,height=10cm,keepaspectratio]{N02/cmr_regions.jpg}
\end{center}
\legende{Fond de carte pour localiser les groupes soudanais au Cameroun : les ensembles du Nord (Extrême-Nord et Nord, en jaune-vert au nord de l'Adamaoua) et l'Adamaoua (en rose), avec Maroua, Garoua et Ngaoundéré (noms en anglais). Au nord du 7\textsuperscript{e} parallèle, les zones paléo-soudanaises (refuges montagnards et plaines inondables) s'imbriquent avec les zones néo-soudanaises (plaines ouvertes propices à l'élevage et aux chefferies centralisées) : à reporter sur le croquis par l'élève. \\ Source du fond de carte : Wikimedia Commons, Peter Fitzgerald, CC BY 3.0.}
\end{popfigure}

\noindent
Cette zone se caractérise par un climat de type \textbf{soudanais} (une saison des pluies courte, 5 à 7 mois, une saison sèche longue et marquée par l'harmattan) et une végétation de savane arborée ou herbeuse. C'est dans ce cadre que deux vagues historiques majeures se sont superposées : les populations « autochtones » de longue date (paléo-soudanais) et les populations conquérantes venues de l'Est et du Nord-Est à partir du XIXe siècle (néo-soudanais).

\begin{aretenir}[Le Grand Nord : cadre commun, destins croisés]
Le Grand Nord (régions du Nord, de l'Adamaoua et de l'Extrême-Nord) est l'espace géographique des groupes soudanais.
\begin{itemize}
    \item \textbf{Climat :} Soudanais (alternance nette saison sèche / saison des pluies).
    \item \textbf{Végétation :} Savane (herbacée, arbustive, arborée selon l'altitude et l'hydromorphie).
    \item \textbf{Division historique :} \textbf{Paléo-soudanais} (installation ancienne, refuge) vs \textbf{Néo-soudanais} (arrivée tardive, conquête, islamisation).
\end{itemize}
\end{aretenir}

\courssub{Les paléo-soudanais : les « anciens » soudanais}

\begin{definition}[Paléo-soudanais et Néo-soudanais]
Les \textbf{paléo-soudanais} (ou \emph{anciens soudanais}) sont les populations soudanaises installées de longue date dans la région, souvent réfugiées dans les zones de relief accidenté (monts Mandara) ou les plaines inondables difficiles d'accès (Logone), qui ont conservé des organisations sociales segmentaires et des religions traditionnelles (animistes), même si le christianisme y a progressé.
Les \textbf{néo-soudanais} (ou \emph{nouveaux soudanais}) sont les populations arrivées plus tardivement, principalement au XIXe siècle, dans le sillage des conquêtes peules (jihad d'Ousman dan Fodio et lamidat d'Adamaoua) et des migrations arabes (Choa). Ils ont imposé une organisation politique centralisée (les \emph{lamidats}) et l'islam.
\end{definition}

\paragraph{Principaux groupes et localisation.}
On distingue plusieurs ensembles ethno-linguistiques (familles tchadique, adamawa-oubanguienne, nilo-saharienne) :
\begin{itemize}
    \item \textbf{Dans les Monts Mandara (Extrême-Nord, dép. Mayo-Tsanaga, Mayo-Sava) :} \textbf{Mafa}, \textbf{Mofu}, \textbf{Mada}, \textbf{Kapsiki}, \textbf{Higi}, \textbf{Gude}, \textbf{Zulgo}, \textbf{Gemzek}. Ces groupes sont souvent regroupés sous l'exonyme péjoratif de \emph{Kirdi} (voir \textbf{Attention} ci-dessous).
    \item \textbf{Dans la plaine du Diamaré et le bassin du Bénoué (Nord, dép. Bénoué, Faro, Mayo-Louti) :} \textbf{Daba}, \textbf{Fali}, \textbf{Namchi}, \textbf{Guider}, \textbf{Dowayo} (Namji).
    \item \textbf{Dans les plaines du Logone et du Chari / Yaérés (Extrême-Nord, dép. Logone-et-Chari, Mayo-Danay) :} \textbf{Massa}, \textbf{Mousseye}, \textbf{Toupouri}, \textbf{Moundang} (aussi au Tchad), \textbf{Kotoko} (sédentaires islamisés anciennement, mais d'origine paléo-soudanaise, pêcheurs/agriculteurs), \textbf{Arabes Choa} (voir néo-soudanais mais installés anciennement).
    \item \textbf{Dans l'Adamaoua (montagnes et plateau, dép. Vina, Mbéré, Djérem, Faro-et-Déo) :} \textbf{Mboum}, \textbf{Ndoro}, \textbf{Tikar} (frange sud, transition bantoue), \textbf{Gbaya} (frange est, transition bantoue), \textbf{Doayo} (Namchi).
\end{itemize}

\paragraph{Habitat et adaptation au milieu : le génie du banco.}
L'habitat paléo-soudanais est une réponse directe aux contraintes du milieu (chaleur, sécheresse, insécurité historique).
\begin{itemize}
    \item \textbf{Cases en banco (terre crue) :} Murs épais (isolant thermique), toits terrasses (récupération eau de pluie, espace de vie nocturne), peu d'ouvertures.
    \item \textbf{Habitat groupé / villages fortifiés :} Dans les Monts Mandara, les villages s'accrochent aux flancs des montagnes (ex : \textbf{Mora}, \textbf{Koza}, \textbf{Mokolo}, \textbf{Rumsiki}), parfois ceints de murailles, offrant une défense naturelle contre les razzias esclavagistes peules.
    \item \textbf{L'architecture Mousgoum (Pouss) :} Un cas d'école unique au monde.
\end{itemize}

\begin{savaistu}[Les cases obus des Mousgoum : un chef-d'œuvre d'architecture durable]
Les \textbf{Mousgoum} (groupe massa, plaine du Logone, arrondissement de Pouss, près de Kousséri) construisent des cases en forme d'obus (\emph{tolék}), entièrement en banco, sans armature bois, ni ciment.
\begin{itemize}
    \item \textbf{Technique :} Modelage à la main, couches successives séchées au soleil. Forme parabolique (catenaire inversée) assurant une stabilité mécanique parfaite sans contreforts.
    \item \textbf{Performance :} Température intérieure constante (≈ 25°C) malgré 45°C dehors. Étanchéité parfaite. Durabilité : 50 ans si entretenues (re-crépissage annuel).
    \item \textbf{Reconnaissance :} Classées au patrimoine mondial de l'UNESCO (liste indicative), étudiées par des architectes du monde entier (ex : Hassan Fathy en Égypte) comme modèle d'habitat bioclimatique low-tech.
\end{itemize}
\end{savaistu}

\paragraph{Organisation socio-politique : la société segmentaire.}
Contrairement aux néo-soudanais, les paléo-soudanais ne connaissent pas l'État centralisé ni la chefferie héréditaire unique. Leur organisation est \textbf{segmentaire} :
\begin{itemize}
    \item \textbf{Unité de base :} Le \emph{clan} (lignage patrilinéaire) ou le \emph{village} autonome.
    \item \textbf{Pouvoir :} Exercé par un \textbf{conseil d'anciens} (gérontocratie) ou un \emph{chef de terre} (premier occupant, garant du lien avec les ancêtres et la terre), dont l'autorité est morale et religieuse, non coercitive.
    \item \textbf{Solidarité :} Mécanisme de \emph{segmentation} : en cas de conflit, les segments (lignages) s'unissent au niveau supérieur (clan, tribu) face à l'ennemi extérieur, puis se séparent une fois la paix revenue. « Moi contre mon frère ; moi et mon frère contre mon cousin ; moi, mon frère et mon cousin contre l'étranger. »
    \item \textbf{Classes d'âge :} Rôle majeur dans la défense, les travaux communautaires, l'initiation.
\end{itemize}

\begin{exemplebox}[Densité exceptionnelle des Monts Mandara]
\begin{exoresolu}[Analyse d'un chiffre clé : la densité des Monts Mandara]
\textbf{Document :} Extrait d'une étude démographique (INS, RGPH 2005 / projections 2020). « L'arrondissement de Koza (Monts Mandara, département du Mayo-Tsanaga) affiche une densité rurale de \textbf{142 hab/km²} sur une superficie de 350 km², soit près de 50 000 habitants. Le département du Mayo-Tsanaga (1 900 km²) dépasse \textbf{100 hab/km²} en moyenne. »

\textbf{Question :} Comment expliquer une telle densité dans un milieu montagneux, sec et à agriculture pluviale extensive ?

\tcblower
\textbf{Solution détaillée :}
\begin{enumerate}
    \item \textbf{Facteur historique (refuge) :} Les populations ont fui les razzias peules (XIXe s.) pour se réfugier dans les montagnes inaccessibles. L'exiguïté de l'espace vital a forcé l'intensification.
    \item \textbf{Facteur technique (agriculture intensive) :} Mise en place de \textbf{terrasses} (cultures en marches d'escalier) sur les pentes, amendement intensif par le fumier (élevage porcin, petit ruminant en stabulation), cultures associées (mil, sorgho, niébé, arachide, moringa).
    \item \textbf{Facteur social :} Forte cohésion clanique (entraide aux champs), main-d'œuvre familiale nombreuse, faible exode rural historique (attachement à la terre des ancêtres).
    \item \textbf{Conséquence actuelle :} Surpâturage, érosion, conflits fonciers latents, début d'exode vers les villes (Maroua, Mokolo) et les plaines du Diamaré.
\end{enumerate}
\textbf{Conclusion :} Cette densité (\textgreater 100 hab/km²) est une \textbf{anomalie géographique} (milieu défavorable + forte population) explicable par l'histoire (refuge) et la technique (terrasses/fumure). Elle illustre la capacité d'adaptation des sociétés paléo-soudanaises.
\end{exoresolu}
\end{exemplebox}

\paragraph{Religions : tradition et christianisme.}
Les paléo-soudanais ont longtemps résisté à l'islamisation. Leurs religions traditionnelles (culte des ancêtres, génies de la nature, terre) restent vivaces. Cependant, depuis la colonisation allemande puis française, et surtout après l'indépendance, le \textbf{christianisme (catholique, protestant, évangélique)} a connu une expansion massive dans les Monts Mandara (diocèses de Maroua-Mokolo, Yagoua), devenant souvent majoritaire chez les Mafa, Mofu, Kapsiki, Toupouri. Cela crée une frontière religieuse nette avec les néo-soudanais musulmans.

\begin{attention}[Le piège du terme « Kirdi »]
\textbf{Ne jamais écrire : « Le peuple Kirdi » ou « L'ethnie Kirdi ».}
\begin{itemize}
    \item \textbf{Étymologie :} Du peul \emph{kirdii} = « païen », « non musulman », « infidèle ». Terme à connotation méprisante, utilisé historiquement pour justifier la razzia (esclavage).
    \item \textbf{Réalité :} Il regroupe une \textbf{vingtaine de groupes distincts} (Mafa, Mofu, Kapsiki, Mada, Fali, etc.) aux langues, cultures et histoires différentes.
    \item \textbf{Usage correct :} « Les populations dites Kirdi » (guillemets obligatoires) ou, mieux, « les groupes paléo-soudanais non islamisés des Monts Mandara ». Au Bac, l'emploi non critique de « Kirdi » comme ethnie unique est une \textbf{erreur grave} sanctionnée.
\end{itemize}
\end{attention}

\courssub{Les néo-soudanais : les « nouveaux » soudanais}

\paragraph{Origines et vague migratoire (XIXe siècle).}
L'arrivée des néo-soudanais bouleverse l'équilibre du Grand Nord au XIXe siècle. Deux flux majeurs :
\begin{enumerate}
    \item \textbf{Les Peuls (Fulbé) :} Originaires du Fouta-Toro (Sénégal) et du Macina (Mali), ils descendent vers l'Est et le Sud. Au Cameroun, c'est \textbf{Modibo Adama} (disciple d'Ousman dan Fodio, calife de Sokoto) qui lance le \textbf{jihad} vers 1810-1847. Il fonde l'\textbf{Empire peul d'Adamaoua} (émirat de Yola), subdivisé en \textbf{lamidats}.
    \item \textbf{Les Arabes Choa (Choua / Shuwa) :} Venant du Kanem-Bornou (Nord-Est du lac Tchad), ils s'installent dans la plaine du Logone-Chari (Kousséri, Logone-Birni, Darak) dès le XVIe siècle mais massivement au XIXe. Éleveurs de zébus et de dromadaires, commerçants transsahariens.
\end{enumerate}

\begin{definition}[Le Lamidat]
Un \textbf{lamidat} est une \textbf{chefferie traditionnelle musulmane} dirigée par un \textbf{lamido} (terme peul signifiant « seigneur », « gouverneur »). C'est l'unité politico-administrative de base de l'ancien empire peul d'Adamaoua.
\begin{itemize}
    \item \textbf{Légitimité :} Religieuse (représentant du calife de Sokoto puis de l'émir de Yola) et héréditaire (choix parmi les fils du lamido défunt par un collège d'électeurs / conseil de la cour).
    \item \textbf{Fonctions :} Chef politique, chef militaire, juge suprême (charia), percepteur de l'impôt (\emph{jangali} sur le bétail, \emph{zakkat} sur les récoltes), distributeur des terres.
    \item \textbf{Exemples majeurs :} Lamidat de \textbf{Ngaoundéré} (Adamaoua), de \textbf{Rey-Bouba}, de \textbf{Tibati}, de \textbf{Maroua}, de \textbf{Garoua}, de \textbf{Mokolo} (créé plus tard sur territoire paléo-soudanais).
\end{itemize}
\end{definition}

\paragraph{Organisation politique : une société hiérarchisée et islamisée.}
La société néo-soudanaise (peule surtout) est pyramidale, rigide, contrastant avec le segmentarisme paléo-soudanais.

\begin{popfigure}
\begin{tikzpicture}[scale=0.9, every node/.style={font=\small}]
% Colonne Gauche : Paléo-soudanais (Segmentaire)
\node[draw=popgreen, thick, rounded corners, fill=popgreenL, minimum width=6cm, minimum height=8cm] (boxP) at (-4,0) {};
\node[anchor=north, font=\bfseries\color{popgreen!80!black}] at (boxP.north) {ORGANISATION PALÉO-SOUDANAISE (SEGMENTAIRE)};

\node[draw, circle, fill=popgreen, text=white, minimum size=1cm] (clan1) at (-6,-1) {Clan A};
\node[draw, circle, fill=popgreen, text=white, minimum size=1cm] (clan2) at (-4,-1) {Clan B};
\node[draw, circle, fill=popgreen, text=white, minimum size=1cm] (clan3) at (-2,-1) {Clan C};

\node[draw, diamond, fill=popgreen!50, text=white, aspect=2, align=center] (conseil) at (-4,-3){Conseil des Anciens \\ (Gérontocratie)};
\node[draw, rectangle, rounded corners, fill=popgreen!30, align=center] (chefTerre) at (-4,-4.5){Chef de Terre \\ (Premier occupant,\\ autorité rituelle)};

\draw[->, thick, popgreen] (clan1) -- (conseil);
\draw[->, thick, popgreen] (clan2) -- (conseil);
\draw[->, thick, popgreen] (clan3) -- (conseil);
\draw[->, thick, popgreen] (conseil) -- (chefTerre);
\draw[<->, dashed, thick, popmuted] (clan1) -- (clan2) node[midway, above, sloped] {Alliance / Conflit};
\draw[<->, dashed, thick, popmuted] (clan2) -- (clan3) node[midway, above, sloped] {Alliance / Conflit};

\node[anchor=north, text width=5.5cm, align=center] at (-4,-6.5) {
    \textbf{Caractères :} \newline
    $\bullet$ Pas de chef unique héréditaire \newline
    $\bullet$ Pouvoir diffus, consensuel \newline
    $\bullet$ Unité = Village / Clan \newline
    $\bullet$ Défense : Classes d'âge \newline
    $\bullet$ Religion : Traditionnelle / Chrétienne
};

% Colonne Droite : Néo-soudanais (Hiérarchique / Lamidat)
\node[draw=poporange, thick, rounded corners, fill=poporangeL, minimum width=6cm, minimum height=8cm] (boxN) at (4,0) {};
\node[anchor=north, font=\bfseries\color{poporange!80!black}] at (boxN.north) {ORGANISATION NÉO-SOUDANAISE (LAMIDAT)};

\node[draw, circle, fill=poporange, text=white, minimum size=1.2cm, align=center] (lamido) at (4,2.5){LAMIDO \\ (Chef suprême,\\Chef religieux)};
\node[draw, rectangle, rounded corners, fill=poporange!40, text=white, align=center] (conseilL) at (4,1){Conseil de la Cour \\ (Waziri, Galadima,\\Sarkin Yaki, etc.)};
\node[draw, diamond, fill=poporange!50, text=white, aspect=2, align=center] (ardo) at (4,-0.5){Ardo'en \\ (Chefs de districts /\\chefs de clans peuls)};
\node[draw, circle, fill=poporange!60, text=white, align=center] (chefV) at (4,-2){Chefs de villages\\(Nommés par Lamido)};
\node[draw, rectangle, fill=poporange!30] (peuple) at (4,-3.5) {Population : \newline Fulbé (nobles/éleveurs) \newline Rimaybé (anciens captifs) \newline Maccuɓe (artisans) \newline Paléo-soudanais (sujets)};

\draw[->, thick, poporange] (lamido) -- (conseilL);
\draw[->, thick, poporange] (conseilL) -- (ardo);
\draw[->, thick, poporange] (ardo) -- (chefV);
\draw[->, thick, poporange] (chefV) -- (peuple);

\node[anchor=north, text width=5.5cm, align=center] at (4,-5.5) {
    \textbf{Caractères :} \newline
    $\bullet$ Pouvoir centralisé, héréditaire \newline
    $\bullet$ Hiérarchie administrative stricte \newline
    $\bullet$ Unité = Lamidat (État) \newline
    $\bullet$ Défense : Cavalerie (Sarkin Yaki) \newline
    $\bullet$ Religion : Islam (Charia)
};

% Flèche de comparaison centrale
\draw[<->, very thick, popblue, double distance=1pt] (-1.2,0) -- (1.2,0) node[midway, above, fill=popcream, rounded corners, inner sep=2pt] {CONTRASTE};
\end{tikzpicture}
\legende{Schéma comparatif des organisations socio-politiques. À gauche, la société segmentaire paléo-soudanaise : pouvoir horizontal, clans autonomes unis par un conseil d'anciens. À droite, la société hiérarchisée néo-soudanaise (lamidat) : pouvoir vertical, centralisé dans les mains du Lamido, relayé par une administration de cour et des chefs de districts (Ardo).}
\end{popfigure}

\noindent
La société peule (Fulbé) est divisée en \textbf{castes / groupes statutaires} (système endogame) :
\begin{itemize}
    \item \textbf{Fulbé (nobles) :} Éleveurs, guerriers, détenteurs du pouvoir politique.
    \item \textbf{Rimaybé (ou \emph{Black Fulani}) :} Descendants d'esclaves capturés chez les paléo-soudanais (Kirdi), intégrés comme pasteurs ou cultivateurs, musulmans, mais statut inférieur.
    \item \textbf{Maccuɓe :} Artisans (forgerons, tisserons, cordonniers, griots), castes héréditaires, statut très bas (interdiction de mariage avec les Fulbé).
    \item \textbf{Sujets paléo-soudanais :} Populations vaincues, payant l'impôt (\emph{haradj}), corvéables, convertis à l'islam pour certains (ex : Kotoko, Massa islamisés), mais restant en marge de la société peule.
\end{itemize}

\paragraph{Activités économiques : l'élevage roi, l'agriculture et la pêche.}
\begin{itemize}
    \item \textbf{Élevage bovin (Zébu) :} Activité emblématique des Peuls. \textbf{Transhumance} Nord-Sud saisonnière (saison sèche : plaines du Logone/Bénoué ; saison des pluies : plateaux de l'Adamaoua). Cheptel national estimé à \textbf{6 à 7 millions de têtes} (majorité au Nord). Production de lait (vente locale), de viande (export vers le Sud, Douala, Yaoundé, Gabon, Congo), de cuir.
    \item \textbf{Agriculture :} Mil (pennisetum), sorgho (sorgho rouge/blanc), maïs (extension récente), riz (plaines inondables du Logone, SEMRY), coton (culture de rente imposée par la colonisation, gérée par la \textbf{SODECOTON}). \textbf{Association mil/niébé/arachide} classique.
    \item \textbf{Pêche :} Activité majeure pour les \textbf{Kotoko} (pêcheurs spécialisés du Logone/Chari, utilisant filets, nasses, barrages) et les \textbf{Massa} (pêche de crue dans les yaérés). Poisson séché (banda) exporté vers le Sud.
\end{itemize}

\begin{exemplebox}[La filière coton : intégration Nord-Sud]
La culture du coton (variété \emph{Allen 333} puis variétés améliorées IRMA) structure l'espace nordiste depuis les années 1950 (CFDT puis \textbf{SODECOTON}).
\begin{itemize}
    \item \textbf{Aménagement :} Création de \textbf{Secteurs de Développement Cotonnier (SDC)} quadrillant le Nord et l'Adamaoua. Pistes, magasins de stockage, usines d'égrenage (Garoua, Maroua, Guider, Ngaoundéré).
    \item \textbf{Encadrement :} Crédit intrants (engrais, herbicides, semences), conseil technique, prix plancher garanti.
    \item \textbf{Impacts :} Monétarisation de l'économie paysanne, scolarisation, mais aussi \textbf{érosion des sols}, \textbf{réduction des parcours pastoraux} (conflits éleveurs/agriculteurs), \textbf{endettement} des producteurs.
\end{itemize}
\end{exemplebox}

\paragraph{Religion : l'islam, ciment identitaire et politique.}
L'islam (sunnite, malékite, confrérie \textbf{Tidjaniya} majoritaire, \textbf{Qadiriya} minoritaire) n'est pas seulement une foi : c'est le \textbf{fondement de la légitimité politique} du Lamido (Commandeur des Croyants local). La \textbf{charia} (loi islamique) régit le statut personnel (mariage, héritage, successions) et foncier (terre = propriété de la communauté/État, gérée par le Lamido). L'arabe est la langue liturgique et administrative historique (ajami : peul écrit en caractères arabes). Aujourd'hui, le français et l'anglais (officiels) coexistent, mais l'arabe reste langue de culture et de commerce transfrontalier (Nigéria, Tchad, RCA).

\courssub{Coexistence, tensions et brassage : vers une société nordiste recomposée}

La distinction paléo/néo n'est pas figée. Le Grand Nord est un laboratoire de \textbf{brassage} :
\begin{itemize}
    \item \textbf{Islamisation par le bas :} De nombreux paléo-soudanais (Massa, Toupouri, Moundang, Fali) se convertissent à l'islam pour s'affranchir du statut de « Kirdi », accéder au commerce, à l'administration, ou par conviction. Ils deviennent « Peuls » sociologiquement (langue, habit, nom) en quelques générations (\emph{fulɓe-isation}).
    \item \textbf{Mariages interethniques :} Fréquents en milieu urbain (Maroua, Garoua, Ngaoundéré, Kousséri).
    \item \textbf{Conflits fonciers persistants :} Le code foncier de 1974 (terres = domaine de l'État) a légalisé l'appropriation des terres de chefferie (paléo) par l'État, souvent gérées par des lamidos (néo). Les paléo-soudanais réclament la reconnaissance du \textbf{droit coutumier} (chef de terre).
    \item \textbf{Crise Boko Haram (depuis 2014) :} A touché indistinctement les deux groupes dans l'Extrême-Nord, forçant des déplacements massifs (PDI : \emph{Personnes Déplacées Internes}) vers le Sud (Mokolo, Maroua) et renforçant par la tragédie un sentiment d'appartenance commune « Nordiste » face à l'ennemi extérieur.
\end{itemize}

\begin{methode}[Comment analyser une carte de répartition des groupes ethniques au Bac ?]
\begin{enumerate}
    \item \textbf{Légende :} Identifier les codes couleurs et hachures de la carte étudiée (un code pour les sociétés paléo-soudanaises, un autre pour les néo-soudanaises). Repérer les symboles (points = villages, triangles = marchés, traits = axes/transhumance).
    \item \textbf{Localisation :} Citer les régions, départements, reliefs (Monts Mandara, Plaine du Bénoué, Yaérés), parallèles (7°N, 13°N).
    \item \textbf{Description :} Noter l'imbrication (paléo en enclave dans le néo, ou néo en couloir dans le paléo). Citer 2-3 groupes par catégorie.
    \item \textbf{Explication (Le « Pourquoi ») :} Lier localisation / milieu / histoire.
        \begin{itemize}
            \item Paléo en montagne (refuge, défense) ou plaines inondables (pêche, riz, isolement).
            \item Néo en plaines ouvertes (élevage, cavalerie, chefferie centralisée).
        \end{itemize}
    \item \textbf{Conséquences / Enjeux :} Densités (forte en montagne), conflits (fonciers, transhumance), intégration (islamisation, urbanisation, Boko Haram).
    \item \textbf{Conclusion :} Synthèse sur la dynamique actuelle (métissage, urbanisation, tensions).
\end{enumerate}
\end{methode}

\begin{exoresolu}[Sujet type Bac : « Montrez que l'espace nord-camerounais est structuré par l'opposition/complémentarité entre paléo-soudanais et néo-soudanais. »]
\tcblower
\textbf{Solution détaillée (Plan détaillé) :}

\textbf{Introduction :}
\begin{itemize}
    \item Accroche : Le Grand Nord (3 régions, 1/3 du territoire, ~30\% pop.) est un carrefour sahélien.
    \item Définitions : Paléo-soudanais (anciens, segmentaires, non islamisés à l'origine) vs Néo-soudanais (Peuls/Arabes, XIXe s., lamidats, islamisés).
    \item Problématique : Comment l'histoire et le milieu ont-ils produit une structure spatiale duale qui évolue vers le brassage ?
    \item Annonce du plan : 1. Une structuration spatiale duale héritée de l'histoire. 2. Des organisations sociales contrastées. 3. Une recomposition actuelle par le brassage et les tensions.
\end{itemize}

\textbf{I. Une structuration spatiale duale : Refuges vs Plaines ouvertes}
\begin{itemize}
    \item \textbf{Paléo-soudanais :} Enclaves de refuge. \textbf{Monts Mandara} (Mafa, Mofu, Kapsiki) : densités \textgreater 100 hab/km², terrasses, cases en banco. \textbf{Yaérés du Logone} (Massa, Kotoko, Toupouri) : plaines inondables, pêche/riz, isolement saisonnier.
    \item \textbf{Néo-soudanais :} Domaines de l'élevage et du pouvoir. \textbf{Plaines du Bénoué} (Garoua, Rey-Bouba) + \textbf{Plateaux de l'Adamaoua} (Ngaoundéré, Tibati) : savane herbeuse, transhumance, lamidats. \textbf{Bassin du Lac Tchad} (Arabes Choa) : zone sahélienne, dromadaires, commerce transfrontalier.
    \item \textbf{Frontière mobile :} Le 7\textsuperscript{e} parallèle n'est pas une ligne nette. Imbrication : lamidats créés dans zones paléo (Mokolo, Koza) ; poches paléo dans zones néo (Fali dans le Bénoué).
\end{itemize}

\textbf{II. Des organisations socio-politiques et religieuses antagonistes}
\begin{itemize}
    \item \textbf{Paléo :} Segmentarisme. Village/Clan autonome. Chef de terre (autorité rituelle). Conseil d'anciens. Religion traditionnelle / Christianisme (post-colonisation). Égalitarisme relatif.
    \item \textbf{Néo :} Hiérarchie étatique (Lamidat). Lamido (pouvoir politique + religieux + foncier). Administration de cour (Waziri, Galadima). Castes (Fulbé, Rimaybé, Maccuɓe). Islam (Charia, Tidjaniya). Inégalités statutaires fortes.
    \item \textbf{Rapport de domination historique :} Razzias (esclaves) $\rightarrow$ Impôt (haradj) $\rightarrow$ Administration coloniale (indirect rule via lamidos) $\rightarrow$ Post-indépendance (élites peules au pouvoir).
\end{itemize}

\textbf{III. Une recomposition accélérée : Brassage, urbanisation, crises}
\begin{itemize}
    \item \textbf{Brassage culturel :} \emph{Fulɓe-isation} (conversion à l'islam, adoption langue peul, noms peuls) des Massa, Toupouri, Moundang, Fali. Mariages mixtes en ville.
    \item \textbf{Urbanisation :} Maroua, Garoua, Ngaoundéré, Kousséri = creusets. Émergence d'une identité « Nordiste » commune (langue véhiculaire : \textbf{Fulfulde}).
    \item \textbf{Tensions persistantes :} Fonciers (Code 1974 vs droit coutumier chef de terre). Pastoraux (transhumance vs champs coton). Sécuritaires (Boko Haram : stigmatisation des Kanouri/Shuwa, déplacés Kirdi et Peuls mélangés dans camps).
    \item \textbf{Intégration nationale :} Élites nordistes (toutes origines) au gouvernement. SODECOTON, SEMRY, Université de Ngaoundéré/Maroua comme vecteurs de mixité.
\end{itemize}

\textbf{Conclusion :}
L'opposition paléo/néo structure encore l'espace (carte) et les mentalités, mais le brassage (islamisation, urbanisation, école, crise sécuritaire) forge une nouvelle identité régionale plurielle. Le défi de l'intégration nationale passe par la reconnaissance juridique du pluralisme foncier et la gestion équitable des ressources (eau, pâture, terre).

\textbf{Ouverture :} Comparer avec la dynamique Sud (Bantou/Bamiléké) : dualité « autochtones/allogènes » aussi présente mais formes différentes (chefferies traditionnelles reconnues vs lamidats).
\end{exoresolu}

\coursec{Les groupes bantou, semi-bantou et les pygmées}

\courssub{Localisation : le Sud-Cameroun (au sud du 7\textsuperscript{e} parallèle)}

Au sud du 7\textsuperscript{e} parallèle, le climat devient de type \textbf{guinéen} (bimodal ou équatorial de transition) puis \textbf{équatorial} (forêt dense). La végétation passe de la savane boisée à la forêt dense humide. C'est l'espace des langues \textbf{bantoues} (famille nigéro-congolaise) et \textbf{semi-bantoues} (bantoïdes, Grassfields).

\begin{popfigure}
\begin{center}
\includegraphics[width=\linewidth,height=10cm,keepaspectratio]{N02/nigercongo.png}
\end{center}
\legende{Les langues nigéro-congolaises en Afrique (légende en anglais). L'aire bantoue (beige) couvre le sud du Cameroun et toute l'Afrique centrale, orientale et australe ; au nord-ouest de cette aire, les langues bantoïdes (orange) et de l'aire Bénoué ; au nord, l'aire adamawa-oubanguienne (vert). Elle montre que le Cameroun se trouve à la jonction entre le domaine bantou, au sud, et les domaines soudanais et adamawa, au nord. Les Pygmées, non figurés, vivent dans la forêt du Sud et de l'Est. \\ Source : Wikimedia Commons, Kimdime et Ulamm, CC BY-SA 3.0.}
\end{popfigure}

\courssub{Les groupes bantous : diversité forestière et côtière}

Les populations bantoues occupent la façade atlantique (Littoral, Sud-Ouest), le Centre, le Sud et l'Est. Elles sont issues de la grande migration bantoue (origines Nigeria/Cameroun, il y a 3000-4000 ans).

\begin{itemize}
    \item \textbf{Groupe Beti-Fang-Bulu (Centre, Sud, Est) :} \textbf{Ewondo} (Yaoundé), \textbf{Bane}, \textbf{Eton}, \textbf{Manguissa}, \textbf{Bulu} (Sud, Ebolowa), \textbf{Fang} (Sud, Est, Gabon/Congo), \textbf{Ntumu}, \textbf{Mvae}. Organisation en \textbf{clans patrilinéaires}, chefferie de clan/village, société secrète \emph{Ngil} (juridique/religieuse, interdite coloniale), culte des ancêtres (\emph{Byeri} : statues gardiennes).
    \item \textbf{Groupe Bassa-Duala (Littoral, Sanaga maritime) :} \textbf{Bassa} (Nord Littoral), \textbf{Duala} (Wouri, Douala), \textbf{Bakoko}, \textbf{Mbo}, \textbf{Bakweri} (Mont Cameroun). Organisation en \textbf{chefferies traditionnelles} (Bell, Akwa, Deido pour Duala), société \emph{Ngondo} (culte de l'eau, jumeaux), forte tradition marchande et côtière.
    \item \textbf{Groupe Makaa-Njem-Kako (Est, Sud-Est) :} \textbf{Makaa}, \textbf{Njem}, \textbf{Kako}, \textbf{Bomwali}, \textbf{Konabembe}. Chasseurs-cueilleurs-agriculteurs, forêt dense, chefferies de village.
    \item \textbf{Autres :} \textbf{Bakundu}, \textbf{Balong}, \textbf{Bafaw} (Sud-Ouest), \textbf{Tikar} (frange Nord-Ouest/Adamaoua, groupe à part, roi \emph{Fon}).
\end{itemize}

\paragraph{Organisation socio-politique : chefferies et sociétés secrètes.}
\begin{itemize}
    \item \textbf{Chefferie traditionnelle :} Héréditaire (souvent primogéniture patrilinéaire), pouvoir politique, judiciaire, foncier (terre des ancêtres), religieux (intermédiaire avec les ancêtres). Ex : \emph{Fon} (Bamiléké/Bamoun/Tikar), \emph{Paramount Chief} (Anglophone), \emph{Chef supérieur} (Francophone).
    \item \textbf{Sociétés secrètes / Initiatiques :} \emph{Ngil} (Fang/Beti, police religieuse), \emph{Ngondo} (Duala/Sawa, eau), \emph{Kwifoyn} (Grassfields, régulation politique), \emph{Male} (Bassa). Rôle de contre-pouvoir, cohésion, justice.
    \item \textbf{Lignage / Clan :} Unité de base patrilinéaire (sauf exceptions matrilinéaires rares). Solidarité forte (funérailles, dot, terres).
\end{itemize}

\paragraph{Activités économiques : agriculture de rente, forêt, ville.}
\begin{itemize}
    \item \textbf{Agriculture d'exportation :} \textbf{Cacao} (Sud, Centre, Est, Sud-Ouest - 1er exportateur agricole), \textbf{Café} (robusta littoral, arabica Ouest), \textbf{Palmier à huile} (Petits planteurs + agro-industrie SOCAPALM, CDC), \textbf{Hévéa} (Sud, HEVECAM), \textbf{Banane} (Plantations du Haut Penja - PHP, CDC).
    \item \textbf{Agriculture vivrière :} Plantain, manioc, macabo, igname, taro, maïs. Système \emph{abattis-brûlis} itinérant (forêt) ou intensif (volcan).
    \item \textbf{Exploitation forestière :} Bois d'œuvre (Ayous, Sapelli, Azobé) - concessions (UFA) gérées par des entreprises (souvent étrangères), redevances pour l'État et communes. Braconnage, cueillette (Gnetum/okok, champignons, miel).
    \item \textbf{Pêche artisanale et industrielle :} Côte (Kribi, Douala, Limbé, Idenau), fleuves (Sanaga, Wouri, Nyong).
    \item \textbf{Industrie / Services :} Douala (poumon économique, port, raffinerie SONARA à Limbé, usines), Yaoundé (administration, services), Kribi (port en eau profonde, gaz GNL, fer Mbalam).
\end{itemize}

\courssub{Les groupes semi-bantous (Grassfields) : l'originalité des Hauts Plateaux}

Occupent l'Ouest et le Nord-Ouest (régions de l'Ouest et du Nord-Ouest), sur les \textbf{Hauts Plateaux de l'Ouest} (alt. 1000-2500 m, volcaniques, sols ferralitiques riches, climat tempéré, fortes pluies). Densité rurale record (\textgreater 200 hab/km² par endroits).

\begin{itemize}
    \item \textbf{Groupes Bamiléké (Ouest) :} \textbf{Bafoussam}, \textbf{Bandjoun}, \textbf{Baham}, \textbf{Bamendjou}, \textbf{Bangangté}, \textbf{Dschang}, \textbf{Fongo-Tongo}, \textbf{Batcham}, etc. (\textasciitilde 100 groupements). Langues bamiléké (proches).
    \item \textbf{Groupes Bamoun (Noun) :} \textbf{Foumban} (capitale historique), sultanat ancien (XIVe s.), écriture \emph{Shümom} (inventée par le roi Njoya vers 1896), islamisés majoritaires (conquête peule avortée, puis conversion royale).
    \item \textbf{Groupes de l'Ouest/Nord-Ouest (Bamileke périphériques / Bamenda) :} \textbf{Bamenda} (Bamileke de l'Ouest), \textbf{Kom}, \textbf{Bafut}, \textbf{Nso}, \textbf{Bali}, \textbf{Wimbum} (Grassfields du Nord-Ouest). Chefferies \emph{Fon}.
\end{itemize}

\paragraph{Organisation : chefferies centralisées et dynamiques.}
\begin{itemize}
    \item \textbf{Le Fon / Chef :} Pouvoir sacré, politique, économique, judiciaire. Assisté de \textbf{notables} (conseil), sociétés secrètes (\emph{Kwifoyn}, \emph{Ngumba}, \emph{Maleu}), reines-mères (\emph{Mafo}).
    \item \textbf{Fonction économique :} Le chef redistribue les terres, organise le travail communautaire, gère les recettes du marché/chefferie. \textbf{Dynamisme entrepreneurial} : élites commerçantes, transporteurs, industriels (boissons, plastiques, agro-alimentaire) investissant au village (maisons modernes, routes, écoles, santé). \emph{« Développement par la base »}.
    \item \textbf{Intégration nationale :} Forte diaspora (Douala, Yaoundé, Nord, étranger). Réseaux de solidarité (\emph{tontines}, associations de développement villageois).
\end{itemize}

\paragraph{Agriculture intensive et pression foncière.}
\begin{itemize}
    \item \textbf{Intensification :} Fumure intensive (cendres, déchets, fumier), cultures associées (maïs, haricot, pomme de terre, taro, café arabica, arbres fruitiers), haies vives (\emph{Erythrina}, \emph{Leucaena}) anti-érosion/clôtures.
    \item \textbf{Pression foncière :} Morcellement extrême (héritage partageable), exode rural massif vers les villes (Bafoussam, Douala, Yaoundé) et les fronts pionniers (Est, Sud-Ouest, Adamaoua) pour le cacao/café/hévéa. Conflits fonciers diaspora/autochtones.
\end{itemize}

\courssub{Les Pygmées : les « premiers habitants » de la forêt}

\begin{definition}[Pygmées (Autochtones de la forêt)]
Populations chasseuses-cueilleuses nomades ou semi-nomades, de petite taille (moyenne \textasciitilde 1,50 m), vivant en symbiose avec la forêt dense humide du bassin du Congo. Au Cameroun : \textbf{Baka} (Sud-Est, arr. Moloundou, Salapoumbé, Yokadouma), \textbf{Bakola/Bagyeli} (Sud, arr. Bipindi, Lolodorf, Kribi), \textbf{Bedzan} (Centre-Est, Mbam-et-Kim, forêt de Nditam). Population estimée \textasciitilde \textbf{40 000 à 60 000} personnes (minorité vulnérable).
\end{definition}

\paragraph{Mode de vie et culture.}
\begin{itemize}
    \item \textbf{Économie :} Chasse (filets, arcs, pièges), cueillette (miel, fruits, champignons, chenilles, \emph{Gnetum}), pêche. Échanges troc avec villageois bantous (viande de brousse vs fer, sel, manioc, tissu, alcool).
    \item \textbf{Habitat :} Campements mobiles (\emph{mongulu} : huttes de feuillages en dôme), déplacements saisonniers selon ressources.
    \item \textbf{Organisation :} Bandes de 15-50 personnes, égalitaires, décisions consensuelles, pas de chef héréditaire. Chef de bande (\emph{Komba}) = chasseur expérimenté.
    \item \textbf{Spiritualité :} \emph{Jengi} (esprit de la forêt, initiation masculine), polyphonie vocale complexe (yodel), pharmacopée forestière pointue.
\end{itemize}

\paragraph{Menaces et droits.}
\begin{itemize}
    \item \textbf{Sédentarisation forcée :} Projets étatiques (années 60-70), pression des Bantous (terres), aires protégées (Parcs Nationaux : Boumba-Bek, Nki, Lobéké, Dja, Campo-Ma'an) interdisant chasse/cueillette traditionnelle.
    \item \textbf{Exploitation :} Travail forcé/servitude chez Bantous, non-paiement soins/école, non-reconnaissance citoyenneté (actes de naissance).
    \item \textbf{Cadre juridique :} Constitution 1996 (minorités), Déclaration ONU 2007 (Droits peuples autochtones), Plan National Développement Pygmées (PNDP). Progrès lents : scolarisation bilingue (Baka/Français), cartes d'identité, consultation préalable (FPIC) pour projets forestiers/miniers.
\end{itemize}

\begin{attention}[Ne pas confondre Pygmées et « Pygmoïdes »]
Le terme « Pygmoïde » (anthropologie physique ancienne) est \textbf{obsolète et péjoratif}. On dit \textbf{Pygmées} ou \textbf{Populations autochtones de la forêt} (terme officiel ONG/ONU). Au Bac, ne pas les assimiler à une « ethnie » unique : Baka, Bakola, Bedzan sont distincts.
\end{attention}

\coursec{L'inégale répartition de la population : les différences de densités}

\courssub{Une densité moyenne qui masque de forts contrastes}

Avec \textasciitilde 60 hab/km² (2023), le Cameroun semble moyennement peuplé. Mais cette moyenne cache une \textbf{répartition extrêmement inégale}, héritée de l'histoire (peuplement, colonisation), du milieu (climat, relief, sols, maladies) et de l'économie (villes, plantations, mines).

\begin{popfigure}
\begin{center}
\includegraphics[width=\linewidth,height=10cm,keepaspectratio]{N02/cmr_dens2000.jpg}
\end{center}
\legende{Carte de la densité de population du Cameroun en 2000 (Gridded Population of the World, CIESIN, Université Columbia ; légende en anglais, en personnes par km\textsuperscript{2}). Les plus fortes densités (brun, 250 à plus de 1~000 hab/km\textsuperscript{2}) forment des taches autour de Douala, de Yaoundé et des hautes terres de l'Ouest ; l'orange (25 à 249) couvre le Littoral, l'Ouest, le Nord-Ouest, le Centre et l'Extrême-Nord ; le plateau de l'Adamaoua, le Nord et le Sud sont peu peuplés (5 à 24) ; l'Est forestier est le \og vide \fg{} (1 à 4). \\ Source : Wikimedia Commons, SEDAC (CIESIN, Columbia University), CC BY 2.0.}
\end{popfigure}

\paragraph{Les trois grands types d'espaces démographiques.}

\begin{description}
    \item[\textbf{1. Les espaces faiblement peuplés (Densité < 10 hab/km²) : « Le vide relatif »}] \hfill \\
    \begin{itemize}
        \item \textbf{Région de l'Est (Kadey, Lom-et-Djérem, Boumba-et-Ngoko) :} Forêt dense, sols pauvres (ferralitiques lessivés), maladies (paludisme, onchocercose, trypanosomiase), isolement (routes mauvaises), activités extensives (chasse, cueillette, exploitation forestière peu employeuse). Pop. : Baka, Makaa, Njem, Kako.
        \item \textbf{Région du Sud (Vallée-du-Ntem, Mvila, Océan hors Kribi/Ebolowa) :} Forêt dense, relief accidenté, sols acides, enclavement.
        \item \textbf{Adamaoua Nord/Est (Mbéré, Djérem, Faro-et-Déo hors Ngaoundéré) :} Plateau soudanais, sols lessivés, tse-tse (trypanosomiase animale), pâturages extensifs, faible urbanisation.
        \item \textbf{Extrême-Nord : Logone-et-Chari Nord (bord du Lac Tchad) :} Zone sahélienne, sols sableux/salins, insécurité (Boko Haram), inondations/assèchement lac.
    \end{itemize}
    \item[\textbf{2. Les espaces moyennement peuplés (10 - 50 hab/km²) : « L'espace de production »}] \hfill \\
    \begin{itemize}
        \item \textbf{Plaines du Bénoué (Nord, dép. Bénoué, Faro, Mayo-Louti) :} Sol fertile (alluvions), coton (SODECOTON), élevage, villes moyennes (Garoua, Guider, Poli). Densité \textasciitilde 30-50 hab/km².
        \item \textbf{Adamaoua Centre (Vina, Ngaoundéré) :} Ville universitaire/administrative, élevage, cultures vivrières, chefferies lamidales. Densité \textasciitilde 20-40 hab/km².
        \item \textbf{Centre/Sud rural (hors axes routiers) :} Agriculture vivrière/rentière (cacao), densité \textasciitilde 15-30 hab/km².
        \item \textbf{Littoral rural (Sanaga-Maritime, Moungo, Nkam) :} Plantations (CDC, PHP, SOCAPALM), main-d'œuvre migrante (Bamiléké, Nordistes), densité \textasciitilde 30-60 hab/km².
    \end{itemize}
    \item[\textbf{3. Les espaces fortement peuplés (> 50 hab/km², souvent > 100-200) : « Les anomalies »}] \hfill \\
    \begin{itemize}
        \item \textbf{Hauts Plateaux de l'Ouest (Grassfields - Bamiléké, Bamoun, Bamenda) :} \textbf{Record national rural : 150 à 300+ hab/km²}. Facteurs : sols volcaniques fertiles, climat sain (pas de tse-tse), histoire (refuge, chefferies organisées), intensification agricole (fumure, haies), forte natalité, faible mortalité. \textbf{Anomalie de Malthus} : densité forte + croissance forte $\rightarrow$ exode massif.
        \item \textbf{Monts Mandara (Extrême-Nord, Mayo-Tsanaga, Mayo-Sava) :} \textasciitilde 100-150 hab/km². Facteurs : refuge historique, terrasses, fumure, cohésion clanique. (Voir ex. résolu).
        \item \textbf{Yaérés du Logone (Mayo-Danay, Logone-et-Chari Sud) :} \textasciitilde 50-100 hab/km². Facteurs : sols alluviaux fertiles (riz, sorgho), pêche, eau disponible, chefferies organisées (Massa, Kotoko, Toupouri).
        \item \textbf{Agglomérations urbaines (Douala, Yaoundé, Bafoussam, Bamenda, Garoua, Maroua, Ngaoundéré) :} Densités \textgreater 1000 hab/km² (communes d'arrondissement). Primatie céphalique (Douala/Yaoundé = 30\% pop. urbaine).
    \end{itemize}
\end{description}

\begin{aretenir}[Facteurs explicatifs de la répartition inégale]
\begin{itemize}
    \item \textbf{Historiques :} Refuges (Monts Mandara, Grassfields), axes de traite/colonisation (chemin de fer, route, fleuves), chefferies centralisées (Grassfields, Lamidats).
    \item \textbf{Physiques :} Sols (volcaniques riches vs ferralitiques pauvres), Climat (sain altitude vs chaud/humide forêt), Eau (fleuves, nappes), Maladies (tsé-tsé, paludisme).
    \item \textbf{Économiques :} Cultures de rente (coton, cacao, café, hévéa, banane), Mines (or, diamant, fer - potentiel), Villes (administration, industrie, port, université).
    \item \textbf{Politiques :} Aménagement (SODECOTON, SEMRY, UNVDA, PHP, CDC), Décentralisation, Insécurité (Boko Haram, crise anglophone).
\end{itemize}
\end{aretenir}

\coursec{Mobilité, brassage des populations et intégration nationale}

\courssub{Grande mobilité spatiale : migrations internes et externes}

Le Cameroun est un pays de \textbf{mobilité intense}. La population n'est pas figée dans les zones de peuplement traditionnel.

\begin{itemize}
    \item \textbf{Migrations Nord $\rightarrow$ Sud (Soudanais $\rightarrow$ Bantou) :} Flux historique (main-d'œuvre plantations CDC/PHP/SOCAPALM, commerce, éducation). \textbf{Peuls} (élevage, commerce bétail/viande), \textbf{Mbororo} (transhumance), \textbf{Nordistes} (fonctionnaires, militaires, étudiants). Création de quartiers « Nord » à Douala (New-Bell, Bépanda), Yaoundé (Mvog-Ada, Essos).
    \item \textbf{Migrations Ouest $\rightarrow$ Partout (Bamiléké/Bamoun $\rightarrow$ Sud, Est, Nord, Villes) :} Diaspora entrepreneuriale la plus visible. Commerce (marchés), transport, industrie, agriculture de rente (fronts pionniers Est/Sud-Ouest), professions libérales. Réseaux de solidarité (\emph{tontines}, associations).
    \item \textbf{Exode rural $\rightarrow$ Urbain :} Moteur principal de l'urbanisation (58\%). Jeunes, scolarisés, quête emploi/santé/éducation. Grossissement des bidonvilles (Douala, Yaoundé).
    \item \textbf{Migrations internationales :} \textbf{Immigration :} Nigérians (commerce, pêche, main-d'œuvre), Tchadiens, Centrafricains, Maliens (commerçants), Chinois (commerce, BTP), Français/occidentaux (expatriés). \textbf{Émigration :} Élites vers France, USA, Canada, Allemagne, Belgique (cerveaux), main-d'œuvre vers pays du Golfe, Gabon, Guinée Équatoriale.
    \item \textbf{Réfugiés / PDI :} \textasciitilde 470 000 réfugiés (majorité Centrafricains à l'Est, Nigérians au Nord - HCR 2023). \textasciitilde 1 million PDI (Extrême-Nord Boko Haram, Nord-Ouest/Sud-Ouest crise anglophone).
\end{itemize}

\courssub{Brassage ethnique et culturel : le creuset camerounais}

Le brassage est le processus de mélange des groupes par le contact prolongé, le mariage, l'école, la langue, le travail, l'espace urbain.

\begin{itemize}
    \item \textbf{Langues véhiculaires :} \textbf{Fulfulde} (Nord, langue de fait), \textbf{Ewondo} (Yaoundé/Centre), \textbf{Duala} (Littoral), \textbf{Pidgin English} (Nord-Ouest/Sud-Ouest), \textbf{Français/Anglais} (officiels, administration, école, médias). \textbf{Camfranglais} (jeunes urbains, mélange français/anglais/argot/vernaculaires).
    \item \textbf{Mariages interethniques :} Fréquents en ville (\textasciitilde 30-40\% couples mixtes à Douala/Yaoundé). Enfants bilingues/biculturels.
    \item \textbf{École et administration :} Internats, universités (Ngaoundéré, Maroua, Yaoundé I/II, Douala, Buea, Dschang, Bamenda), fonctionnaires affectés hors région d'origine. Brassage forcé mais efficace.
    \item \textbf{Fêtes et culture :} NGONDO (Sawa), NYEM-NYEM (Adamaoua), MEDUMBA (Bamiléké), Festival des Arts et Cultures (FESTAC), Coupe d'Afrique des Nations (CAN 2021) $\rightarrow$ sentiment national.
    \item \textbf{Religion :} Islam (Nord + poches Sud), Christianisme (Sud + poches Nord), syncrétismes. Dialogue interreligieux (Comité national).
\end{itemize}

\begin{exemplebox}[Le « Camerounais » type : portrait statistique d'un brassage]
Un cadre supérieur né à \textbf{Bafoussam} (Bamiléké), scolarisé à \textbf{Ngaoundéré} (université, vit avec des Peuls/Mboum), muté à \textbf{Bertoua} (Est, vit avec des Makaa/Baka), marié à une \textbf{Bassa} (Littoral), travaillant à \textbf{Douala} (quartier cosmopolite), parlant \textbf{Français, Anglais, Pidgin, Medumba, Fulfulde}. Ses enfants vont à l'école internationale, parlent Camfranglais. \textbf{C'est la norme, pas l'exception.}
\end{exemplebox}

\courssub{Tensions, conflits et enjeux d'intégration}

Le brassage ne va pas sans heurts. L'intégration nationale reste un \textbf{chantier permanent}.

\begin{itemize}
    \item \textbf{Conflits fonciers :} \textbf{Code foncier 1974 (modifié 2005)} : Terre = Domaine national (État). Conflit \textbf{Droit moderne (titre foncier) vs Droit coutumier (chef de terre / Fon / Lamido)}. Ex : Plantations agro-industrielles (Herakles Farms/SG-SOC Sud-Ouest, HEVECAM Sud) vs villages. Ex : Monts Mandara (Chef de terre vs Lamido/État).
    \item \textbf{Conflits agro-pastoraux :} Transhumance (Peuls/Mbororo) vs Agriculture (coton, vivrier). Couloirs de transhumance obstrués par champs. Dégradation ressources. Commissions mixtes (MINADER/MINEPIA) souvent inefficaces.
    \item \textbf{Crise anglophone (Nord-Ouest/Sud-Ouest, depuis 2016) :} Revendications identitaires (Common Law, éducation, marginalisation perçue) $\rightarrow$ conflit armé, sécessionnisme (Ambazonie), destructions, déplacements, rupture brassage.
    \item \textbf{Insécurité Extrême-Nord (Boko Haram / ISWAP) :} Stigmatisation ethnique (Kanouri, Choa, Peuls), recrutement forcé, méfiance intercommunautaire, PDI.
    \item \textbf{Tribalisme / Népotisme :} Accès emploi public, marchés, promotions perçus selon origine ethnique/régionale. « Quota » implicite. Frein au mérite et à l'unité.
\end{itemize}

\begin{methode}[Comment traiter un sujet sur « Population et Intégration nationale » au Bac ?]
\begin{enumerate}
    \item \textbf{Analyser le sujet :} Mots-clés (Répartition, Mobilité, Brassage, Intégration, Tensions, Politiques). Type (Commentaire carte, Dissertation, Croquis).
    \item \textbf{Structurer l'espace :} Carte mentale 3 zones (Nord/Soudanais, Ouest/Semi-Bantou, Sud-Est/Bantou-Pygmées) + Villes.
    \item \textbf{Mobiliser les données :} Chiffres clés (Pop. totale, taux croissance, % urbain, densités types, cheptel, réfugiés/PDI).
    \item \textbf{Articuler les échelles :} Local (village, chefferie) $\leftrightarrow$ Régional (Grand Nord, Grassfields) $\leftrightarrow$ National (Douala/Yaoundé, État) $\leftrightarrow$ International (CEMAC, CEEAC, Diaspora).
    \item \textbf{Nuancer :} Éviter déterminisme ethnique. Montrer dynamiques (brassage, métissage, urbanisation) vs inerties (foncier, chefferies, crises).
    \item \textbf{Conclusion prospective :} Intégration = processus continu (éducation, emploi, justice foncière, décentralisation effective, paix).
\end{enumerate}
\end{methode}

\coursec{TP 2 : Construction et lecture de la carte de répartition de la population au Cameroun}

\paragraph{Objectifs du TP.}
\begin{itemize}
    \item Construire une carte choroplèthe des densités de population par région (données RGPH 2005 / Projections INS 2020).
    \item Lire et interpréter la carte : identifier les contrastes, expliquer les facteurs, lier aux groupes humains et à l'intégration.
    \item Maîtriser la sémiologie graphique (choix des seuils, couleurs, légende ordonnée).
\end{itemize}

\paragraph{Document 1 : Données statistiques (Projections INS 2020, arrondis).}

\begin{center}
\begin{tabularx}{\linewidth}{|l|X|c|c|c|}\hline
\textbf{Région} & \textbf{Chef-lieu} & \textbf{Superficie (km²)} & \textbf{Population (2020)} & \textbf{Densité (hab/km²)} \\ \hline
Adamaoua & Ngaoundéré & 63 691 & 1 200 000 & 19 \\ \hline
Centre & Yaoundé & 68 953 & 4 500 000 & 65 \\ \hline
Est & Bertoua & 109 002 & 950 000 & 9 \\ \hline
Extrême-Nord & Maroua & 34 263 & 4 500 000 & 131 \\ \hline
Littoral & Douala & 20 248 & 3 800 000 & 188 \\ \hline
Nord & Garoua & 66 090 & 2 800 000 & 42 \\ \hline
Nord-Ouest & Bamenda & 17 300 & 2 100 000 & 121 \\ \hline
Ouest & Bafoussam & 13 892 & 2 800 000 & 202 \\ \hline
Sud & Ebolowa & 47 191 & 800 000 & 17 \\ \hline
Sud-Ouest & Buéa & 25 410 & 1 800 000 & 71 \\ \hline
\textbf{Total} & & \textbf{475 442} & \textbf{25 250 000} & \textbf{53 (moy.)} \\ \hline
\end{tabularx}
\end{center}
\textit{Source : INS, Annuaire statistique 2021 / Projections RGPH 2005. Note : Les chiffres sont des ordres de grandeur pour l'exercice.}

\paragraph{Consignes de construction (sur fond de carte muette des 10 régions).}

\begin{methode}[Étapes de la carte choroplèthe des densités]
\begin{enumerate}
    \item \textbf{Calculer les densités} (si non fournies) : $D = \frac{Pop}{Sup}$. Vérifier les totaux.
    \item \textbf{Choisir la discrétisation (seuils de classes) :} Méthode des \textbf{seuils naturels (Jenks)} ou \textbf{échelons géométriques} pour bien séparer les extrêmes. Proposition 5 classes :
        \begin{itemize}
            \item Classe 1 : \textbf{Très faible} $< 20$ hab/km² (Est, Sud, Adamaoua) $\rightarrow$ \textcolor{popgoldL!80}{Jaune très clair}.
            \item Classe 2 : \textbf{Faible} $20 - 50$ hab/km² (Nord, Adamaoua centre) $\rightarrow$ \textcolor{poporangeL!60}{Orange clair}.
            \item Classe 3 : \textbf{Moyenne} $50 - 100$ hab/km² (Centre hors Yaoundé, Sud-Ouest, Littoral rural) $\rightarrow$ \textcolor{poporange!60}{Orange moyen}.
            \item Classe 4 : \textbf{Forte} $100 - 150$ hab/km² (Extrême-Nord, Nord-Ouest) $\rightarrow$ \textcolor{poppink!60}{Rose/Orange foncé}.
            \item Classe 5 : \textbf{Très forte} $> 150$ hab/km² (Ouest, Littoral/Douala) $\rightarrow$ \textcolor{poppink}{Rouge/Rose vif}.
        \end{itemize}
    \item \textbf{Colorier chaque région} avec la couleur de sa classe. Ne pas colorier les chefs-lieux (points).
    \item \textbf{Construire la légende ordonnée} (du clair au foncé, ou faible à fort) : Surface colorée + Libellé classe + Effectif (nb régions).
    \item \textbf{Compléter les éléments cartographiques obligatoires :}
        \begin{itemize}
            \item \textbf{Titre :} « Répartition de la population au Cameroun par région (Densités 2020) ».
            \item \textbf{Échelle graphique} (ex: 1 cm = 100 km).
            \item \textbf{Orientation} (Flèche Nord).
            \item \textbf{Source} : INS, Projections 2020.
            \item \textbf{Auteur / Date}.
        \end{itemize}
    \item \textbf{Symboles ponctuels (optionnel mais recommandé) :} Proportionnels à la population urbaine des chefs-lieux (Douala, Yaoundé gros cercles ; Bertoua, Ebolowa petits).
\end{enumerate}
\end{methode}

\paragraph{Exemple de légende finale (à reproduire).}

\begin{center}
\begin{tabular}{|c|l|c|}\hline
\textbf{Symbole} & \textbf{Libellé} & \textbf{Nb Régions} \\ \hline
\textcolor{popgoldL!80}{\rule{8pt}{8pt}} & Très faible ($< 20$ hab/km²) & 3 (Est, Sud, Adamaoua) \\ \hline
\textcolor{poporangeL!60}{\rule{8pt}{8pt}} & Faible ($20 - 50$) & 1 (Nord) \\ \hline
\textcolor{poporange!60}{\rule{8pt}{8pt}} & Moyenne ($50 - 100$) & 3 (Centre*, Sud-Ouest, Littoral*) \\ \hline
\textcolor{poppink!60}{\rule{8pt}{8pt}} & Forte ($100 - 150$) & 2 (Extrême-Nord, Nord-Ouest) \\ \hline
\textcolor{poppink}{\rule{8pt}{8pt}} & Très forte ($> 150$) & 1 (Ouest) \\ \hline
{\large $\bullet$} \textcolor{popdark}{\rule{4pt}{4pt}} & Chef-lieu (Pop. urbaine) & Échelle : 1 mm = 100 000 hab \\ \hline
\end{tabular}
\end{center}
\emph{* Hors communes urbaines de Yaoundé et Douala qui faussent la densité régionale.}

\paragraph{Questions d'analyse et d'interprétation (Type Bac).}

\begin{exercice}[Analyse de la carte construite]
\begin{enumerate}
    \item \textbf{Description :} Identifier les 3 grands ensembles de densité (Faible, Moyenne, Forte/Très forte) et citer les régions concernées.
    \item \textbf{Explication :} Pour chaque ensemble, mobiliser \textbf{au moins 2 facteurs} (physiques, historiques, économiques, politiques) vus en cours.
    \item \textbf{Lien Groupes humains :} Comment la carte des densités recoupe-t-elle (ou pas) la carte des grands groupes humains (Soudanais, Bantous, Semi-Bantous, Pygmées) ? Donner 2 exemples de correspondance et 1 exemple de discordance.
    \item \textbf{Intégration nationale :} En quoi cette répartition inégale et cette mobilité (flux Nord-Sud, Ouest-Est, Rural-Urbain) constituent-elles à la fois un \textbf{atout} (brassage, marché national, main-d'œuvre) et un \textbf{défi} (conflits fonciers, pression urbaine, marginalisation périphéries) pour l'intégration nationale ?
    \item \textbf{Croquis rapide :} Sur un croquis de synthèse (A4 paysage), représenter : Zones de densité (3 couleurs), 2 flèches de migration majeure (Nord$\rightarrow$Sud, Ouest$\rightarrow$Est/Sud), 2 zones de tension (Extrême-Nord, Nord-Ouest/Sud-Ouest), 2 pôles d'attraction (Douala, Yaoundé).
\end{enumerate}
\end{exercice}

\begin{corrige}[Exercice TP 2 - Éléments de réponse]
\begin{enumerate}
    \item \textbf{Description :} \textbf{Faible} (Est 9, Sud 17, Adamaoua 19) : « Vide démographique » forestier/sahélien. \textbf{Moyenne}(Nord 42) : Plaine du Bénoué, coton, élevage. \textbf{Forte/Très forte} (Extrême-Nord 131, Nord-Ouest 121, Ouest 202, Littoral 188) : Monts Mandara/Yaérés, Grassfields, Métropoles.
    \item \textbf{Explication :}
        \begin{itemize}
            \item \textit{Faible :} Forêt dense/sols pauvres/maladies (Est, Sud) ; Plateau soudanais/tse-tse/extensif (Adamaoua).
            \item \textit{Moyenne :} Plaine alluviale Bénoué/coton/SODECOTON (Nord) ; Agriculture rente/vivrière + villes moyennes (Centre, Sud-Ouest, Littoral rural).
            \item \textit{Forte/Très forte :} Refuge historique/sols volcaniques/intensification/chefferies (Ouest, Nord-Ouest) ; Refuge Monts Mandara/terrasses (Extrême-Nord) ; Primatie urbaine/port/industrie (Littoral/Douala, Centre/Yaoundé).
        \end{itemize}
    \item \textbf{Lien Groupes humains :}
        \begin{itemize}
            \item \textit{Correspondance :} Haute densité Ouest/Nord-Ouest $\leftrightarrow$ Semi-Bantous (Bamiléké/Bamoun) ; Haute densité Monts Mandara $\leftrightarrow$ Paléo-soudanais (Mafa, Kapsiki).
            \item \textit{Discordance :} Faible densité Est/Sud $\leftrightarrow$ Bantous (Makaa, Baka) mais aussi Pygmées ; Forte densité Extrême-Nord $\leftrightarrow$ Mélange Paléo (Massa, Kotoko) et Néo (Peuls, Choa) inextricables.
        \end{itemize}
    \item \textbf{Intégration :}
        \begin{itemize}
            \item \textit{Atouts :} Brassage ethnique (mariages, langues véhiculaires Fulfulde/Pidgin/Français), marché national unifié (flux bétail Nord$\rightarrow$Sud, vivriers Sud$\rightarrow$Nord, main-d'œuvre Ouest$\rightarrow$Plantations), élites nationales multirégionales.
            \item \textit{Défis :} Conflits fonciers (titre moderne vs coutumier), pressions agro-pastorales (couloirs bouchés), crises sécuritaires (Boko Haram Extrême-Nord, Sécessionnisme NO/SO) fragmentant l'espace, bidonvillisation métropoles, tribalisme gestion publique.
        \end{itemize}
    \item \textbf{Croquis de synthèse :} Carte du Cameroun (contour 10 régions). Légende 3 couleurs densité. Flèches épaisses : Nord$\rightarrow$Douala/Yaoundé (bétail/commerce), Ouest$\rightarrow$Est/Sud-Ouest (cacao/café/hévéa). Symboles tension : Éclair rouge Extrême-Nord (Boko Haram), Éclair rouge NO/SO (Crise anglophone). Carrés rouges : Douala, Yaoundé (Pôles). Échelle, Nord, Titre, Source.
\end{enumerate}
\end{corrige}

\begin{aretenir}[Bilan de la leçon : Les grands groupes humains et l'intégration]
\begin{itemize}
    \item \textbf{Trois ensembles majeurs :} Soudanais (Nord, Paléo/Néo), Bantous/Semi-Bantous (Sud/Ouest), Pygmées (Forêt Sud-Est).
    \item \textbf{Organisations contrastées :} Segmentaires (Paléo) vs Centralisées/Lamidats (Néo) vs Chefferies/Fon (Semi-Bantous/Bantous) vs Bandes égalitaires (Pygmées).
    \item \textbf{Répartition inégale :} « Vide » forestier/sahélien vs « Plein » Hauts Plateaux/Monts Mandara/Yaérés vs Métropoles primatiques.
    \item \textbf{Mobilité/Brassage intenses :} Migrations Nord-Sud, Ouest-Tout, Rural-Urbain $\rightarrow$ Identité nationale plurielle (Camfranglais, élites mixtes).
    \item \textbf{Intégration :} Processus inachevé. Enjeux : Justice foncière, Gestion pastorale, Résolution crises (NO/SO, Extrême-Nord), Lutte contre tribalisme, Décentralisation effective.
\end{itemize}
\end{aretenir}

\coursec{Les groupes bantou et semi-bantou ; les pygmées}

Après avoir parcouru les vastes plaines soudanaises du Nord, franchissons la ligne de partage des eaux de l'Adamaoua. Nous basculons vers le Sud, dans l'univers vert et humide de la forêt équatoriale et des hautes herbes de l'Ouest. C'est ici que vivent les \cle{Bantou}, les \cle{Semi-bantou} et les \cle{Pygmées}. Si les Soudanais sont les « fils de la savane », les Bantous sont les « enfants de la forêt » et les Semi-bantous, ceux de la « transition ». Comprendre leur localisation, leur histoire et leur organisation, c'est saisir l'âme culturelle du Cameroun méridional.

\courssub{Localisation et origine : le grand Sud forestier et les migrations bantoues}

\begin{definition}[Les Bantous : peuplement et langues du Sud forestier]
Les \textbf{Bantous} désignent l'ensemble des populations occupant la zone forestière du Sud-Cameroun (régions du Centre, du Sud, de l'Est, du Littoral et une partie du Sud-Ouest). Ils parlent des langues appartenant à la grande famille \textbf{bantoue} (groupe Bantu au sens linguistique strict), issues d'une expansion historique majeure partie de la zone frontalière Cameroun-Nigéria (région du Grassfields actuel) il y a environ 3000 à 4000 ans.
\end{definition}

\begin{propriete}[Le foyer de dispersion bantou]
Le consensus scientifique (linguistique, archéologique, génétique) situe le \textbf{berceau bantou} dans la région des \cle{Grassfields} (hautes terres de l'Ouest et du Nord-Ouest actuels). De là, plusieurs vagues migratoires se sont échelonnées :
\begin{itemize}
    \item Vers l'Est et le Sud-Est (forêt dense) : ancêtres des Maka, Nzimé, Kako, Baka (contact).
    \item Vers le Sud-Ouest (littoral atlantique) : ancêtres des Sawa (Duala, Bakweri, Limba).
    \item Vers le Centre et le Sud (forêt à \textit{Triplochiton} / ayous) : ancêtres des Beti-Fang (Ewondo, Eton, Bulu, Fang), Bassa, Mvae.
\end{itemize}
Cette migration s'est faite par \textbf{petits groupes lignagers}, s'adaptant au milieu forestier (agriculture sur brûlis, pêche, chasse).
\end{propriete}

\begin{popfigure}
\begin{center}
\includegraphics[width=\linewidth,height=10cm,keepaspectratio]{N02/nigercongo.png}
\end{center}
\legende{Les langues nigéro-congolaises en Afrique (légende en anglais). On y lit l'aire bantoue (beige) qui part du sud-est du Cameroun et du golfe de Guinée et s'étend vers le sud et l'est, et, à la limite nord-ouest de cette aire, les langues bantoïdes (orange) : c'est la région d'origine supposée de l'expansion bantoue. Les limites sont des zones de transition, non des frontières rigides. \\ Source : Wikimedia Commons, Kimdime et Ulamm, CC BY-SA 3.0.}
\end{popfigure}

\begin{methode}[Légender une carte de répartition ethnolinguistique]
Pour réussir un croquis ou une carte sur les groupes humains au Bac :
\begin{enumerate}
    \item \textbf{Fond de carte} : Trace le contour du pays, les axes hydrographiques majeurs (Sanaga, Nyong, Wouri, Dja) et les 2-3 villes repères (Yaoundé, Douala, Bafoussam, Bertoua, Garoua, Maroua).
    \item \textbf{Grands ensembles d'abord} : Utilise 3-4 couleurs bien contrastées pour les \textbf{quatre macro-groupes} : Soudanais (Nord), Semi-bantous (Ouest/Nord-Ouest), Bantous (Sud/Est/Centre/Littoral), Pygmées (Sud-Est). Ne colorie pas peuple par peuple (illisible).
    \item \textbf{Exemples ciblés} : Place 2 à 3 noms de peuples \emph{représentatifs} par zone (ex: Beti, Bassa, Duala pour Bantous ; Bamiléké, Bamoun pour Semi-bantous ; Baka pour Pygmées).
    \item \textbf{Flèches historiques} : Si le sujet porte sur l'origine, trace des flèches de migration depuis le foyer (Grassfields) vers les zones d'installation actuelles.
    \item \textbf{Légende organisée} : Titre précis, hiérarchie des symboles (surfaces, points, flèches), échelle, orientation, source.
\end{enumerate}
\end{methode}

\begin{exemplebox}[Le couloir de peuplement du Sanaga]
La vallée du \textbf{Sanaga} a joué un rôle d'autoroute naturelle pour la migration bantoue. Les groupes \textbf{Bassa} et \textbf{Beti} ont progressé le long de ce fleuve et de ses affluents (Djérem, Lom, Pangar) avant de se disperser sur les interfluves. C'est pourquoi on retrouve une continuité linguistique et culturelle le long de cet axe, du plateau de l'Adamaoua jusqu'à l'estuaire du Wouri.
\end{exemplebox}

\courssub{Les peuples bantous : diversité et organisations socio-culturelles}

La zone bantoue n'est pas un bloc monolithique. Elle se divise en grands ensembles linguistiques et culturels que tu dois savoir situer et caractériser.

\begin{center}
\begin{tabularx}{\linewidth}{|l|X|X|}\hline
\rowcolor{popblueL}
\textbf{Ensemble} & \textbf{Peoples principaux (exemples)} & \textbf{Aire géographique type} \\ \hline
\textbf{Beti-Fang} & Ewondo (Yaoundé), Eton, Bulu (Sangmélima, Ebolowa), Fang (Sud, Est, Gabon, Guinée Eq.) & Centre, Sud, Est (forêt à \textit{Triplochiton}) \\ \hline
\textbf{Sawa (Côtiers)} & Duala (Wouri), Bakweri (pentes du Cameroun), Limba, Mungo, Isubu & Littoral, Sud-Ouest (côte \& montagne) \\ \hline
\textbf{Bassa} & Bassa, Bakoko, Yabassi & Littoral (Sanaga maritime), Centre (Nord-Ouest de Yaoundé) \\ \hline
\textbf{Est / Maka-Nzimé} & Maka, Nzimé (Njémé), Kako, Kwakum, Pol, Mbimo & Est (Bertoua, Batouri, Yokadouma), forêt dense humide \\ \hline
\textbf{Autres} & Mvae, Ngumba, Bafia (zone de transition Centre) & Zones de contact forêt/savane \\ \hline
\end{tabularx}
\end{center}

\begin{propriete}[Organisation socio-culturelle bantoue : le lignage et le village]
Malgré leur diversité, les sociétés bantoues partagent des structures communes :
\begin{itemize}
    \item \textbf{Sociétés lignagères (patrilinéaires le plus souvent)} : L'unité de base est le \textbf{lignage} (ensemble de descendants d'un ancêtre commun), regroupé en \textbf{clan} (ex: clans \textit{Etom}, \textit{Nnen} chez les Beti). La filiation, l'héritage (terre, titres) et les alliances matrimoniales s'organisent autour du lignage.
    \item \textbf{Le village (chefferie de village)} : Unité politique de base. Le \textbf{chef de village} (souvent le doyen du lignage fondateur ou élu par les notables) gère la terre, la justice coutumière, les rites agraires. Son pouvoir est \textbf{contrôlé par les sociétés secrètes} et le conseil des anciens.
    \item \textbf{Sociétés secrètes et rites d'initiation} : Véritables « poumons » de la démocratie traditionnelle et gardiennes de l'ordre social.
    \begin{itemize}
        \item \textbf{So'o (ou So)} chez les \textbf{Beti-Fang} : Société initiatique masculine puissante, régulateur politique, justice, lien avec les ancêtres.
        \item \textbf{Ngondo} chez les \textbf{Sawa (Duala, Bakweri)} : Culte de l'eau (Jengu), assemblée traditionnelle, fête annuelle sur les rives du Wouri (patrimoine UNESCO 2023).
        \item \textbf{Male} (Bassa), \textbf{Nkwa} (Maka) : Équivalents fonctionnels.
    \end{itemize}
    \item \textbf{Religion traditionnelle} : Culte des \textbf{ancêtres} (médiateurs auprès de Dieu créateur, \textit{Zambe} chez les Beti, \textit{Nyambe} chez les Bassa), croyance aux forces de la nature (esprits de l'eau, de la forêt), divination, médecine traditionnelle.
\end{itemize}
\end{propriete}

\begin{exoresolu}[Analyse d'un document : le Ngondo, facteur d'intégration]
\textbf{Document :} Extrait d'un article sur la fête du Ngondo à Douala (2023). « Le Ngondo rassemble les peuples Sawa autour du culte de Jengu, génie de l'eau. C'est un moment de cohésion, de règlement des conflits coutumiers et d'affirmation identitaire face à la modernité. Les chefs traditionnels y prêtent serment. »

\textbf{Question :} Montre que le Ngondo illustre la permanence et l'adaptation des institutions traditionnelles bantoues.

\textbf{Solution détaillée :}
\begin{enumerate}
    \item \textbf{Identification} : Le Ngondo est une institution \textbf{Sawa} (groupe bantou littoral). C'est une \textbf{société secrète / assemblée traditionnelle} (type \textit{So'o} chez les Beti).
    \item \textbf{Fonction d'intégration (permanence)} : « Rassemble les peuples Sawa » $\rightarrow$ unité inter-clans/ethnies (Duala, Bakweri, Mungo...). « Règlement des conflits coutumiers » $\rightarrow$ rôle judiciaire/politique du chef et des notables (structure villageoise). « Serment des chefs » $\rightarrow$ légitimation du pouvoir coutumier.
    \item \textbf{Adaptation (modernité)} : « Affirmation identitaire face à la modernité » $\rightarrow$ l'institution ne disparaît pas, elle se réinvente comme marqueur culturel (patrimoine UNESCO 2023) et espace de dialogue avec l'État (décentralisation).
    \item \textbf{Conclusion} : Le Ngondo prouve que les structures bantoues (lignage, chefferie, société secrète) ne sont pas des « survivances » mais des acteurs vivants de la gouvernance locale et de l'intégration nationale.
\end{enumerate}
\end{exoresolu}

\courssub{Les semi-bantous : les Grassfields, chefferies puissantes et dynamiques démographiques}

\begin{definition}[Les Semi-bantous : la zone de transition (Grassfields)]
Les \textbf{Semi-bantous} (ou \textbf{Bantoïdes de l'Ouest}) occupent les \textbf{hautes terres de l'Ouest} (région de l'Ouest) et du \textbf{Nord-Ouest} (région du Nord-Ouest). Cette zone, appelée \cle{Grassfields} (« champs d'herbes »), est une mosaïque de savanes d'altitude et de forêts de galerie. Leurs langues et cultures présentent un \textbf{mélange de traits bantous} (vocabulaire de base, agriculture) \textbf{et soudanais} (organisation politique centralisée, chefferies sacrées, architecture).
\end{definition}

\begin{popfigure}
\begin{center}
\includegraphics[width=\linewidth,height=11cm,keepaspectratio]{N02/bantu_phil.png}
\end{center}
\legende{Schéma de l'expansion bantoue en Afrique (modèle de Phillipson ; légende numérotée sans texte). En haut, les flèches rouges partent du foyer d'origine (point 1, aux confins du Cameroun et du Nigéria) puis se dirigent vers l'est (2a), vers le sud le long de la côte atlantique (2b, 5) et vers l'Afrique australe (7a, 7b). En bas, une autre représentation montre une dispersion en étoile depuis l'Afrique centrale (points 8 à 10). \\ Source : Wikimedia Commons, domaine public.}
\end{popfigure}

\begin{propriete}[Organisation semi-bantoue : les chefferies d'État (Fondoms et Sultanat)]
La spécificité majeure des Grassfields est l'existence de \textbf{chefferies complexes, hiérarchisées et centralisées}, souvent qualifiées de « royaumes » ou « chefferies d'État » :
\begin{itemize}
    \item \textbf{Le Fon (Bamiléké) / Le Sultan (Bamoun)} : Chef suprême, détenteur du pouvoir politique, judiciaire, religieux et foncier. Il est le garant de la fertilité (terre, femmes, pluie). Succession souvent patrilinéaire mais avec règles complexes (ex: choix par les notables \textit{Mkam Bvuh} chez les Bamiléké).
    \item \textbf{Hiérarchie administrative} : Le Fon/Sultan $\rightarrow$ Chefs de groupements/quartiers $\rightarrow$ Chefs de villages $\rightarrow$ Chefs de familles/lignages.
    \item \textbf{Sociétés secrètes structurantes} : 
    \begin{itemize}
        \item \textbf{Tukah (ou Kuosi)} : Société des éléphants/léopards, régulateur politique, police coutumière, masque à perles (art majeur).
        \item \textbf{Mbansi (ou Ngiri)} : Société des femmes / reine mère (\textit{Mafo}), contre-pouvoir essentiel.
        \item \textbf{Kamveu, Kemdje} : Sociétés de guerriers, de notables.
    \end{itemize}
    \item \textbf{Artisanat et écriture} : Maîtrise de la métallurgie (fer, laiton), sculpture sur bois, perles, teinture \textit{ndop}. Point d'orgue : \textbf{l'écriture Shümom} (Bamoun).
\end{itemize}
\end{propriete}

\begin{savaistu}[Le Shümom : une écriture africaine inventée au Cameroun]
Le \textbf{Sultan Ibrahim Njoya} (r. 1886--1933) de \textbf{Foumban} (peuple Bamoun) a créé au début du XX\textsuperscript{e} siècle une écriture originale, le \textbf{Shümom} (ou A-ka-u-ku). Partant de pictogrammes, il l'a simplifiée en un syllabaire d'environ 80 signes pour transcrire le bamoun. Il a fondé une école, une imprimerie, rédigé l'histoire des Bamoun, un manuel de médecine, une pharmacopée. Bien que réprimée par l'administration coloniale française, elle est aujourd'hui enseignée à l'Université de Dschang et au musée du Palais royal de Foumban (rénové en 2024). C'est l'une des rares écritures africaines indigènes encore vivantes.
\end{savaistu}

\begin{aretenir}[Densité record dans les Grassfields]
Les hautes terres de l'Ouest (pays Bamiléké : Bafoussam, Dschang, Mbouda, Bangangté) comptent parmi les \textbf{densités rurales les plus élevées d'Afrique}, dépassant souvent \num{300} à \num{500} hab/km² (locales > \num{1000} hab/km²). Cette surpopulation relative s'explique par : (1) sols volcaniques fertiles, (2) climat tempéré sain (paludisme faible), (3) organisation sociale favorisant l'intensification agricole (terrasses, fumure, cultures associées), (4) forte émigration vers les villes (Douala, Yaoundé) qui décongestionne le terroir.
\end{aretenir}

\begin{attention}[Ne pas confondre « Semi-bantou » et « Métis »]
\textbf{Erreur fréquente :} Dire que les Semi-bantous sont un « mélange biologique » entre Bantous et Soudanais.
\textbf{Correction :} C'est une \textbf{classification linguistique et culturelle}. Leurs langues (bantoïdes) sont apparentées au bantou mais ont conservé des traits archaïques ou emprunté au soudanais (notamment le système de classes nominaux simplifié). Leur organisation politique (chefferies centralisées) rappelle les modèles sahéliens/soudanais, mais leur base agricole (igname, maïs, taro, café) et leur habitat groupé sont typiques de la zone forestière de transition. Les Bamiléké ne sont pas des « Peuls noircis » ni des « Bantous islamisés » : ils sont les héritiers directs de la civilisation des Grassfields.
\end{attention}

\courssub{Les Pygmées : premiers habitants de la forêt, enjeux contemporains}

\begin{definition}[Les Pygmées (Autochtones de la forêt équatoriale)]
Les \textbf{Pygmées} (terme conventionnel, parfois remplacé par \textbf{« Peuples autochtones de la forêt »}) sont les \textbf{premiers occupants connus} de la forêt dense humide du bassin du Congo, bien avant l'arrivée des Bantous (il y a > 40 000 ans). Au Cameroun, on distingue trois groupes principaux :
\begin{itemize}
    \item \textbf{Baka} (Est, Sud-Est : Moloundou, Lomié, Yokadouma) -- le plus nombreux ($\sim$ \num{30000} à \num{40000}).
    \item \textbf{Bagyeli (ou Bakola)} (Sud : Bipindi, Kribi, Campo, forêt de Campo-Ma'an, Lolodorf) -- $\sim$ \num{4000} à \num{10000}.
    \item \textbf{Bakola} (parfois distingués des Bagyeli, zone de transition Est/Sud).
\end{itemize}
Ils ne sont \textbf{ni Bantous, ni Soudanais, ni Semi-bantous}. Ils constituent un \textbf{groupe humain à part}, génétiquement et culturellement distinct.
\end{definition}

\begin{attention}[Piège de l'examen : Classification des Pygmées]
\textbf{Ne classe jamais les Pygmées parmi les Bantous} sous prétexte qu'ils vivent au Sud. Ils ont leur propre groupe linguistique (langues oubangiennes pour les Baka, bantoues pour les Bagyeli/Bakola mais avec substrat propre), leur morphologie spécifique (petite taille adaptée à la forêt), et surtout leur \textbf{mode de vie chasseur-cueilleur nomade/séminomade} qui les oppose structurellement aux agriculteurs sédentaires bantous.
\end{attention}

\begin{propriete}[Mode de vie et organisation des Pygmées (Baka)]
\begin{itemize}
    \item \textbf{Économie} : \textbf{Chasse-cueillette} (arc, filets, pièges), pêche, récolte de miel, chenilles, champignons, ignames sauvages (\textit{Dioscorea}). Échanges traditionnels avec les villageois bantous (viande, miel, travail agricole contre fer, sel, manioc, tissu, alcool) : relation \textbf{« patron-clients »} souvent inégalitaire.
    \item \textbf{Habitat} : \textbf{Campements mobiles} (mongulu : huttes en feuillage en dôme), déplacés au rythme des ressources (saison du miel, des chenilles, de la chasse). Sédentarisation forcée croissante le long des pistes forestières.
    \item \textbf{Social} : Groupes de 15 à 50 personnes, égalitarisme marqué, pas de chef héréditaire (leadership situationnel : meilleur chasseur, ancien respecté). Décisions par consensus.
    \item \textbf{Spiritualité} : \textbf{Jengi (ou Djengi)} : Esprit majeur de la forêt, médiateur, célébré par des rites polyphoniques (polyphonie vocale yeyi, danse). Connaissance encyclopédique de la pharmacopée forestière.
\end{itemize}
\end{propriete}

\begin{exemplebox}[Marginalisation et droits : le cas du Parc de Lobéké]
La création d'aires protégées (Parc national de Lobéké, Boumba-Bek, Nki -- \textbf{Trinational de la Sangha}, patrimoine mondial UNESCO 2012) a souvent conduit à l'\textbf{expulsion des Baka} de leurs territoires ancestraux sans consentement libre, préalable et éclairé (CLPE). Ils deviennent « braconniers » sur leurs propres terres. Des ONG (WWF, GIZ, Okani) et l'État tentent désormais la \textbf{gestion participative} (zones d'usage communautaire, emploi d'écogardes Baka, partage des revenus du tourisme/taxes forestières). C'est un enjeu majeur d'\textbf{intégration nationale} et de justice environnementale.
\end{exemplebox}

\courssub{Religions et syncrétismes au Sud : christianisme, tradition, modernité}

Le paysage religieux du Sud-Cameroun est dominé par le christianisme, mais la tradition reste le socle culturel.

\begin{center}
\begin{tabularx}{\linewidth}{|l|X|X|}\hline
\rowcolor{popblueL}
\textbf{Composante} & \textbf{Caractéristiques} & \textbf{Exemples / Réalités} \\ \hline
\textbf{Christianisme (majoritaire)} & \textbf{\textasciitilde 70--80\%} au Sud. Catholique (arrivée 1890, Pallotins, Jésuites) et Protestant (Baptistes, Presbytériens, Luthériens, arrivés 1840s côte). Églises de réveil (pentecôtistes) en forte croissance. & Siège archidiocèse Yaoundé, Douala. Université catholique d'Afrique centrale. Écoles/hôpitaux confessionnels majeurs. \\ \hline
\textbf{Religion traditionnelle} & Culte des ancêtres, forces de la nature, divination, sociétés secrètes (So'o, Ngondo, Tukah). Pratiquée souvent \textbf{en parallèle} (double appartenance). & Rites funéraires (veillées, seconds enterrements), médecine traditionnelle, serments sur fétiches (Nkwa, Ngang). \\ \hline
\textbf{Islam (minorité au Sud)} & Présence ancienne via commerce (Haoussa, Peuls) et conversions récentes (élites, migrants nordistes). Minoritaire (< 5\% au Sud, mais visible en villes). & Mosquées à Yaoundé, Douala, Bafoussam. Chefs musulmans intégrés aux chefferies (ex: Lamido de Bafoussam). \\ \hline
\textbf{Syncrétismes \& Nouvelles religiosités} & \textbf{Églises prophétiques / spirituelles} (ex: Église du Christ au Cameroun - ECC, Kimbanguisme importé). \textbf{Mouvements de guérison/délivrance}. Réappropriation de rites traditionnels « christianisés ». & Culte des ancêtres « sanctifié » (messe de requiem), utilisation d'eau bénite / huile d'onction comme substituts aux pharmacopées. \\ \hline
\end{tabularx}
\end{center}

\begin{methode}[Analyser le fait religieux dans une zone d'étude (Bac)]
Si un sujet porte sur « Peuplement et religions » ou « Intégration nationale et diversité culturelle » :
\begin{enumerate}
    \item \textbf{Cartographie} : Superpose la carte des groupes humains et la carte des aires d'influence religieuse (Catholicisme au Centre/Sud/Est ; Protestantisme au Littoral/Ouest/Nord-Ouest ; Islam au Nord ; Traditionnel partout en sous-jacent).
    \item \textbf{Ne sépare pas} : Montre que l'appartenance ethnique ne détermine pas l'appartenance religieuse (un Bamiléké peut être catholique, musulman, protestant ou traditionaliste).
    \item \textbf{Facteur d'intégration / division} : L'Église (écoles, santé, médias) structure le territoire et crée du lien national (ex: Journées Mondiales de la Jeunesse, pèlerinages mariaux à Mvolyé, Marianberg). Mais les tensions existent (conflits fonciers Église/État, radicalismes, instrumentalisation politique de l'ethnie/religion).
\end{enumerate}
\end{methode}

\begin{exoresolu}[Commentaire de document statistique : Évolution religieuse (Recensements 1987, 2005, estimations 2020)]
\textbf{Document :} Tableau simplifié (pourcentages arrondis, population totale Cameroun).
\begin{center}
\begin{tabular}{|l|c|c|c|}\hline
\rowcolor{popblueL}
\textbf{Appartenance} & \textbf{1987} & \textbf{2005} & \textbf{Est. 2020} \\ \hline
Catholique & 38\% & 38\% & \textasciitilde 35\% \\ \hline
Protestant & 18\% & 22\% & \textasciitilde 25\% \\ \hline
Autres chrétiens (Réveil) & 3\% & 6\% & \textasciitilde 15\% \\ \hline
Musulman & 16\% & 21\% & \textasciitilde 25\% \\ \hline
Animiste / Traditionnel & 22\% & 6\% & \textasciitilde 5\% (déclaré) \\ \hline
Sans / Autre & 3\% & 7\% & \textasciitilde 5\% \\ \hline
\end{tabular}
\end{center}

\textbf{Question :} Analyse les grandes tendances de l'évolution religieuse au Cameroun entre 1987 et 2020.

\textbf{Solution détaillée :}
\begin{enumerate}
    \item \textbf{Données brutes} : Catholiques stables puis légère baisse (38\% $\rightarrow$ 35\%). Protestants classiques en hausse (18\% $\rightarrow$ 25\%). \textbf{Explosion des Églises de réveil} (3\% $\rightarrow$ 15\%). Islam en forte progression (16\% $\rightarrow$ 25\%). Animiste déclaré en chute libre (22\% $\rightarrow$ 5\%).
    \item \textbf{Interprétation géographique} : 
    \begin{itemize}
        \item \textbf{Progression de l'Islam} : Migration Nord $\rightarrow$ Sud (commerçants, étudiants, fonctionnaires), conversions en milieu urbain/bantou, démographie forte au Nord.
        \item \textbf{Essor du protestantisme évangélique} : Réponse aux crises urbaines (chômage, santé), prosélytisme médiatique, attraction pour la jeunesse, implantation forte dans les zones bantoues et semi-bantoues (Douala, Yaoundé, Bafoussam).
        \item \textbf{Chute de l'animisme « déclaré »} : Pas disparition de la tradition, mais \textbf{requalification identitaire} : on se déclare chrétien/musulman tout en pratiquant les rites ancestraux (double appartenance). Le recensement force un choix unique.
    \end{itemize}
    \item \textbf{Conclusion} : Le Cameroun reste un pays profondément religieux. La laïcité de l'État (Constitution) gère un pluralisme croissant. L'intégration nationale passe par la régulation du fait religieux (loi 2000 sur la liberté de culte, lutte contre l'extrémisme).
\end{enumerate}
\end{exoresolu}

\begin{aretenir}[Synthèse : Les trois piliers du peuplement méridional]
\begin{itemize}
    \item \textbf{Bantous (Sud forestier)} : Agriculteurs lignagers, chefferies villageoises, sociétés secrètes (So'o, Ngondo), christianisés majoritaires. Densité faible à moyenne.
    \item \textbf{Semi-bantous (Grassfields)} : Agriculteurs intensifs, \textbf{chefferies d'État centralisées} (Fon, Sultan), sociétés secrètes (Tukah), artisanat/écriture (Shümom). \textbf{Densités rurales records} (\num{300}--\num{500} hab/km²), forte diaspora.
    \item \textbf{Pygmées (Forêt profonde)} : Chasseurs-cueilleurs nomades, égalitaires, spirituels (Jengi), connaissances forestières uniques. \textbf{Marginalisés}, enjeux droits fonciers / aires protégées.
\end{itemize}
Ces trois mondes ne sont pas étanches : \textbf{brassage permanent} (mariages, migrations travail, commerce, école, politique) forge l'unité nationale dans la diversité.
\end{aretenir}

\coursec{L'inégale répartition de la population : les différences de densités}

\courssub{Transition : des groupes humains à leur répartition sur le territoire}

Après avoir identifié les \cle{grands groupes humains} qui composent la mosaïque camerounaise — \cle{Paléo-Soudanais}, \cle{Néo-Soudanais}, \cle{Bantous}, \cle{Semi-Bantous} et \cle{Pygmées} — nous devons maintenant comprendre \emph{comment} et \emph{pourquoi} ces populations se distribuent sur le territoire national. La carte des groupes ethniques ne se superpose pas à la carte des densités : une même région peut accueillir plusieurs groupes, et un même groupe peut être présent dans des zones de densité très contrastée. C'est cette \cle{répartition inégale} que nous allons analyser, car elle conditionne l'aménagement du territoire, la pression sur les ressources et les défis de l'\cle{intégration nationale}.

\begin{popfigure}
\begin{center}
\includegraphics[width=\linewidth,height=11cm,keepaspectratio]{N02/cmr_dens.png}
\end{center}
\legende{Carte de la densité de population des régions du Cameroun en quatre classes (légende en anglais, en hab/km\textsuperscript{2}). \textbf{Plus de 50} (bleu foncé) : Extrême-Nord, Nord-Ouest, Ouest et Littoral. \textbf{20 à 50} (bleu clair) : Sud-Ouest et Centre. \textbf{10 à 20} (vert) : Nord. \textbf{Moins de 10} (vert pâle) : Adamaoua, Est et Sud. Les contrastes opposent les hautes terres, le Littoral et l'Extrême-Nord au vide de l'Est et du Sud. \\ Source : Wikimedia Commons, Austiger, CC BY-SA 4.0.}
\end{popfigure}

\courssub{Rappel : effectifs totaux et densité moyenne nationale}

\begin{definition}[Densité de population]
La \cle{densité de population} est le rapport entre le nombre d'habitants d'un territoire et sa superficie. Elle s'exprime en habitants par kilomètre carré (\cle{hab/km²}).
\[
D = \frac{\text{Population totale}}{\text{Superficie (km²)}}
\]
C'est un \emph{indicateur de concentration} (intensif), à ne pas confondre avec l'effectif total (indicateur extensif).
\end{definition}

\begin{exemplebox}[Effectifs et densité moyenne du Cameroun (projections récentes)]
Selon les projections de l'\cle{INS} (Institut National de la Statistique) et du \cle{BUCREP} (Bureau Central des Recensements et des Études de Population), la population du Cameroun est estimée à environ \textbf{28 à 30 millions d'habitants} (projection 2023-2024). Pour une superficie de \SI{475 442}{km²}, la densité moyenne nationale se situe entre \textbf{55 et 60 hab/km²}.

\begin{center}
\begin{tabularx}{\linewidth}{|l|X|X|}\hline
\textbf{Indicateur} & \textbf{Valeur approximative} & \textbf{Observation} \\ \hline
Population totale & 28,5 -- 30 millions & Croissance \SI{\sim 2,6}{\%} / an \\ \hline
Superficie & \SI{475 442}{km²} & 53\textsuperscript{e} rang mondial \\ \hline
Densité moyenne & \textbf{58 hab/km²} & Masque d'énormes contrastes internes \\ \hline
\end{tabularx}
\end{center}
\end{exemplebox}

\courssub{Les zones densément peuplées : les « pleins » démographiques}

Trois ensembles majeurs concentrent la majorité de la population sur une fraction réduite du territoire.

\begin{enumerate}
\item \textbf{L'axe Littoral -- Ouest -- Nord-Ouest (le « triangle de forte densité »)} : C'est le cœur démographique du pays.
      \begin{itemize}
      \item \cle{Région du Littoral} : Avec \textbf{> 250 hab/km²} (Douala seule dépasse 3,5 millions d'habitants sur \SI{923}{km²}, soit \SI{\sim 3800}{hab/km²} en zone urbaine), c'est la région la plus dense. Port, industrie, services attirent les flux migratoires.
      \item \cle{Région de l'Ouest} (Pays Bamiléké) : \textbf{150 à 200 hab/km²}. Haute densité rurale historique liée à l'agriculture intensive sur sols volcaniques fertiles, organisation sociale forte (chefferies) et pression foncière séculaire.
      \item \cle{Région du Nord-Ouest} : \textbf{100 à 150 hab/km²}. Montagnes (Bamenda Highlands), sols fertiles, climat tempéré, agriculture de plantation (thé, café) et vivrière intensive.
      \end{itemize}
\item \textbf{La zone des Monts Mandara (Extrême-Nord, arrondissements de Mokolo, Koza, Tokombéré)} : Densités locales > 200 hab/km² sur les piémonts et versants (agriculture en terrasses, refuge historique), contrastant avec la plaine du Logone.
\item \textbf{Les métropoles : Yaoundé (Centre) et Douala (Littoral)} : \cle{Primatie urbaine}. Yaoundé (environ 4,5 millions d'habitants, \SI{\sim 2800}{hab/km²} en zone urbaine) et Douala concentrent à elles deux \textbf{près d'un quart de la population nationale} sur moins de 1 \% du territoire.
\end{enumerate}

\begin{aretenir}[Contraste saisissant]
Le \cle{Littoral} (superficie \SI{20 248}{km²}) et l'\cle{Est} (superficie \SI{109 002}{km²}, soit 5,4 fois plus grand) illustrent le paradoxe : le Littoral compte \textbf{plus de 4 millions d'habitants} (densité \SI{\sim 270}{hab/km²}), l'Est \textbf{environ 1 million} (densité \SI{\sim 9}{hab/km²}). Le rapport de densité dépasse \textbf{1 à 30}.
\end{aretenir}

\courssub{Les zones moyennement peuplées : zones de transition et d'axes}

Ces régions présentent des densités comprises entre \textbf{20 et 100 hab/km²}. La population s'y organise le long de \cle{axes structurants} (routes, fleuves, lignes de chemin de fer).

\begin{itemize}
\item \textbf{Région du Centre (hors Yaoundé)} : Densité \SI{\sim 35}{hab/km²}. Peuplement le long de l'axe Yaoundé--Ngaoundéré (route/rail) et du fleuve Sanaga. Agriculture de rente (cacao, café) et vivrière.
\item \textbf{Région du Sud-Ouest} : Densité \SI{\sim 65}{hab/km²}. Façade atlantique (Limbe, Buea), plantations agro-industrielles (CDC, PAMOL), sols volcaniques fertiles au pied du Mont Cameroun.
\item \textbf{Région de l'Adamaoua (axes)} : Densité moyenne \SI{\sim 25}{hab/km²}, mais très inégale. Concentration le long de l'axe Ngaoundéré--Garoua (route nationale N1, chemin de fer Transcamerounais) et autour des pôles urbains (Ngaoundéré, Tibati, Meiganga). Hors axes, la densité chute < 10 hab/km².
\item \textbf{Plaine du Logone (Extrême-Nord, départements du Logone-et-Chari, Mayo-Danay)} : Densité \SI{50--80}{hab/km²}. Alluvions fertiles, pêche, riziculture, chefferies Kotoko historiques. Zone de contact sahélien/soudanien.
\item \textbf{Région du Nord (bassin de la Bénoué, piémonts de l'Adamaoua)} : Densité \SI{\sim 40}{hab/km²}. Autour de Garoua (port fluvial, industrie cotonnière SODECOTON) et dans les piémonts (Monts Atlantika, Alantika) où les populations de refuge (Mafa, Kapsiki) maintiennent des densités élevées localement.
\end{itemize}

\courssub{Les zones faiblement peuplées : les « vides » démographiques}

Deux grands ensembles couvrent près de \textbf{40 \% du territoire} pour moins de \textbf{10 \% de la population} (densité < 20 hab/km², souvent < 10).

\begin{enumerate}
\item \textbf{Le Grand Sud forestier (Régions de l'Est et du Sud)} :
      \begin{itemize}
      \item \cle{Région de l'Est} : \SI{109 002}{km²}, \SI{\sim 9}{hab/km²}. Forêt dense humide, sols ferrallitiques lessivés (pauvres), maladies vectorielles (onchocercose, paludisme), isolement routier (peu d'axes bitumés), absence de pôles urbains forts (Bertoua, Batouri < 100 000 hab.).
      \item \cle{Région du Sud} : \SI{47 191}{km²}, \SI{\sim 13}{hab/km²}. Contexte forestier similaire, un peu plus ouvert sur la côte (Kribi, port en eau profonde, pipeline Tchad-Cameroun, hévéa SOCAPALM/HEVECAM) mais l'intérieur reste vide.
      \end{itemize}
\item \textbf{L'Adamaoua « hors axes » et le Nord sahélien (Extrême-Nord septentrional)} :
      \begin{itemize}
      \item Plateau de l'Adamaoua (hors corridor N1) : Pâturages extensifs (élevage transhumant Mbororo), sols ferrugineux tropicaux, climat tropical d'altitude, trypanosomose (mouche tsé-tsé) historiquement limitante.
      \item Extrême-Nord au nord du 11\textsuperscript{e} parallèle (département du Mayo-Sava, Mayo-Tsanaga hors Monts Mandara) : Zone sahélienne, pluviométrie < 700 mm, sols sableux, insécurité (Boko Haram depuis 2014 accentuant le vide et les déplacements).
      \end{itemize}
\end{enumerate}

\begin{popfigure}
\begin{tikzpicture}
\begin{axis}[
    ybar,
    bar width=7pt,
    width=\linewidth,
    height=9cm,
    enlarge x limits=0.15,
    ylabel={Densité (hab/km²)},
    ylabel style={font=\small},
    xlabel={Régions (ordre géographique Nord-Sud)},
    xlabel style={font=\small},
    xtick=data,
    xticklabels={Extr.-Nord, Nord, Adamaoua, N.-Ouest, Ouest, Littoral, Centre, S.-Ouest, Est, Sud},
    xticklabel style={rotate=45, anchor=east, font=\tiny},
    ymin=0,
    ymax=300,
    ymajorgrids=true,
    grid style={dashed, popmuted},
    legend style={at={(0.5,-0.38)}, anchor=north, legend columns=2, font=\footnotesize, cells={anchor=west}},
    axis line style={popdark},
    tick style={popdark},
    major grid style={popmuted},
]
\addplot[poppink!80, fill=poppink!60, draw=poppink!80] coordinates {(1,68) (2,42) (3,24) (4,130) (5,180) (6,270) (7,38) (8,65) (9,9) (10,13)};
\addlegendentry{Densité moyenne régionale (est. 2023)}
% Ligne densité moyenne nationale
\addplot[popblue, dashed, thick, forget plot] coordinates {(0.5,58) (10.5,58)};
\addlegendentry{Moyenne nationale (~58)}
\end{axis}
\end{tikzpicture}
\legende{Diagramme en barres : Densité de population moyenne par région du Cameroun (estimations récentes, hab/km²). La ligne pointillée bleue représente la densité moyenne nationale (~58 hab/km²). On observe le pic extrême du Littoral et de l'Ouest, et le creux de l'Est et du Sud. Source : Projections INS/BUCREP, calculs personnels.}
\end{popfigure}

\courssub{Facteurs explicatifs de la répartition inégale}

La répartition résulte de l'interaction entre facteurs \cle{naturels} (potentiel biophysique), \cle{historiques} (héritage des peuplements) et \cle{économiques} (attractivité actuelle).

\courssub{Facteurs naturels : le cadre biophysique}

\begin{itemize}
\item \textbf{Relief et altitude} : Les hautes terres de l'Ouest (1000--2500 m) et du Nord-Ouest offrent un \cle{climat tempéré} (paludisme moins virulent), des sols volcaniques fertiles (andosols) et une pluviométrie abondante (> 2000 mm). À l'inverse, les plaines basses de l'Extrême-Nord (sahélien) et les plateaux de l'Est/Sud (forêt dense) sont moins favorables à une forte densité rurale traditionnelle.
\item \textbf{Sols} : Sols volcaniques (Ouest, Mont Cameroun) et alluvions (Logone, Wouri, Bénoué) $\rightarrow$ forte potentialité agricole $\rightarrow$ forte densité. Sols ferrallitiques lessivés (Est, Sud) ou ferrugineux tropicaux lessivés (Adamaoua, Nord) $\rightarrow$ fragilité, besoin de jachère longue $\rightarrow$ faible densité.
\item \textbf{Climat et pluviométrie} : Gradient Nord-Sud. Soudanien (900--1500 mm) favorable aux cultures vivrières (sorgho, maïs, coton) $\rightarrow$ densités moyennes à fortes. Guinéen (1500--2000 mm) et Équatorial (2000 mm+) $\rightarrow$ forêt dense, sols lessivés, maladies $\rightarrow$ faible densité, sauf zones volcaniques (Ouest, SW).
\item \textbf{Hydrographie} : Fleuves (Sanaga, Bénoué, Logone, Wouri) = axes de communication, pêche, irrigation, villes-portuaires (Garoua, Douala, Kribi). Mais zones marécageuses (basse Sanaga, delta Wouri, Logone) = maladies, inondations $\rightarrow$ répulsion locale.
\item \textbf{Maladies endémiques} : \cle{Onchocercose} (cécité des rivières) a vidé les vallées de la Sanaga, Mbam, Bénoué jusqu'aux campagnes de traitement (années 1970-90, OMS/OCP). \cle{Paludisme} holoendémique en zone forestière et de plaine. \cle{Trypanosomose} (mouche tsé-tsé) a longtemps bloqué l'élevage et le peuplement dans l'Adamaoua et le Nord.
\end{itemize}

\courssub{Facteurs historiques et humains : l'héritage du peuplement}

\begin{itemize}
\item \textbf{Migrations et refuges} : Les \cle{Paléo-Soudanais} (Mafa, Kapsiki, etc.) se sont réfugiés dans les \cle{Monts Mandara} et \cle{Monts Atlantika} (zones de forte densité actuelle) pour échapper aux razzias esclavagistes (XIX\textsuperscript{e} s.) et à l'islamisation forcée. Les \cle{Bantous} ont migré du Sud vers le Centre/Est, mais la forêt a freiné l'expansion démographique.
\item \textbf{Traite négrière et colonisation} : La traite a dépeuplé les zones côtières (Douala, Kribi) et les axes fluviaux (Sanaga) aux XVII\textsuperscript{e}--XIX\textsuperscript{e} s. La colonisation allemande puis franco-britannique a créé des \cle{pôles artificiels} : ports (Douala, Kribi, Victoria/Limbé), capitales administratives (Yaoundé, Ngaoundéré, Garoua), chemin de fer (Transcamerounais Douala--Ngaoundéré), plantations (Sud-Ouest, Littoral). Cela a fixé les populations et créé des « pleins » nouveaux.
\item \textbf{Politiques de peuplement post-indépendance} : Opérations « \cle{Champs de la Révolution} » (années 1970) dans l'Adamaoua et l'Est (échec relatif). Déplacement forcé de populations pour barrages (Lagdo, Song Loulou, Memve'ele). Crise anglophone (depuis 2017) : déplacements internes massifs du Nord-Ouest/Sud-Ouest vers le Centre, Littoral, Ouest, modifiant les densités locales.
\item \textbf{Sécurité} : Insécurité dans l'Extrême-Nord (Boko Haram) $\rightarrow$ déplacements vers le Sud du département (Mokolo, Maroua) et vers le Nord/Adamaoua. Insécurité dans le N-O/S-O $\rightarrow$ surpeuplement relatif des villes d'accueil (Bafoussam, Douala, Yaoundé).
\end{itemize}

\courssub{Facteurs économiques : l'attractivité différentielle}

\begin{itemize}
\item \textbf{Industrialisation et emploi urbain} : Douala (capitale économique, \SI{\sim 50}{\%} de l'industrie formelle) et Yaoundé (capitale politique, administration, services) sont des \cle{aimants migratoires} majeurs. Bafoussam, Garoua, Maroua, Ngaoundéré, Bertoua, Ebolowa jouent le rôle de pôles régionaux secondaires.
\item \textbf{Infrastructures de transport} : L'axe \cle{Douala--Yaoundé--Ngaoundéré} (route bitumée, rail, pipeline, fibre optique) structure une « diagonale de forte densité ». L'Est et le Sud restent en marge (route bitumée Yaoundé--Bertoua--Batouri--Yokadouma récente mais trafic faible ; Kribi--Lolodorf bitumé pour le port/mines).
\item \textbf{Ressources minières et forestières} : Exploitation forestière (Est, Sud) et minière (bauxite Minim-Martap, fer Mbalam/Nabeba, or/diamant Est) créent des poches de densité (camps, villes minières) mais peu d'effet d'entraînement durable sur la densité rurale.
\item \textbf{Agriculture de rente} : Cacao/café (Centre, Sud, Sud-Ouest, Littoral), coton (Nord, Adamaoua), banane/plantations (Sud-Ouest, Littoral, PHP, CDC, Boh Plantations) fixent des populations salariées et paysannes.
\end{itemize}

\begin{methode}[Analyser une carte de répartition de la population (croquis ou choroplèthe)]
\begin{enumerate}
\item \textbf{Identifier les grandes classes de densité} (légende) : forte / moyenne / faible. Repérer les seuils (ex: >100, 20-100, <20 hab/km²).
\item \textbf{Dégager les grands ensembles spatiaux} : « triangle Ouest/Littoral/N-O », « diagonale Centre/Adamaoua/Nord », « vide forestier Est/Sud », « plein sahélien Logone/Mandara ».
\item \textbf{Expliquer CHAQUE contraste} en mobilisant \textbf{au moins un facteur naturel} (relief, sol, climat, maladie) \textbf{ET un facteur humain/historique} (migration, colonisation, ville, infrastructure, sécurité, politique).
\item \textbf{Nuancer} : distinguer densité rurale (agricole) et densité urbaine (artificielle, tertiaire). Noter les « poches » (Monts Mandara dans l'Extrême-Nord, piémonts Adamaoua dans le Nord).
\item \textbf{Conclure} sur les déséquilibres (surpeuplement rural Ouest vs sous-peuplement Est) et les enjeux d'aménagement (désenclavement Est, décongestionnement Ouest, développement pôles secondaires).
\end{enumerate}
\end{methode}

\courssub{Mobilité, brassage des populations et intégration nationale}

La forte \cle{mobilité spatiale} (migrations pendulaires, saisonnières, rurales-urbaines, interrégionales, forcées) est la conséquence directe des déséquilibres de densité et de développement. Elle génère un intense \cle{brassage des populations}.

\begin{itemize}
\item \textbf{Types de mobilité} :
      \begin{itemize}
      \item \emph{Migrations rurales-urbaines} : Exode massif des jeunes de l'Ouest, du Nord, de l'Adamaoua vers Douala, Yaoundé, Bafoussam.
      \item \emph{Migrations interrégionales « pionnières »} : Bamiléké (Ouest) vers l'Est, le Sud, l'Adamaoua, le Sud-Ouest (plantations, commerce, agriculture).
      \item \emph{Migrations forcées} : Déplacés internes (Extrême-Nord : Boko Haram ; N-O/S-O : crise anglophone) vers zones plus sûres.
      \item \emph{Mobilité pendulaire/seasonnière} : Transhumance Mbororo (Adamaoua/Nord/Extrême-Nord), commerce transfrontalier (CEMAC, Nigeria).
      \end{itemize}
\item \textbf{Conséquences : le brassage} :
      \begin{itemize}
      \item \emph{Positives} : Mariages interethniques, marchés intercommunautaires (ex: Marché Central Douala, Marché de Bafoussam), émergence de langues véhiculaires (\cle{fulfulde} au Nord, \cle{pidgin english} au SW/NW, \cle{camfranglais} urbain), enrichissement culturel.
      \item \emph{Tensions} : Conflits fonciers (autochtones vs allochtones, ex: Est, Adamaoua, Sud-Ouest), revendications identitaires, pression sur services urbains (eau, électricité, voirie, déchets), insécurité urbaine.
      \end{itemize}
\item \textbf{Enjeu d'intégration nationale} : Le brassage est le ciment de la nation \emph{si} il est encadré par des politiques d'aménagement équitables (PND 2020-2030, Schémas Directeurs d'Aménagement Régionaux). L'objectif : transformer la contrainte démographique en atout (main-d'œuvre, marché intérieur) par le \cle{désenclavement} (Est, Adamaoua), la \cle{décentralisation} effective (transferts de compétences/ressources aux CTD) et la \cle{création de pôles de croissance secondaires} (pôles agro-industriels, corridors de développement).
\end{itemize}

\courssub{Conséquences : déséquilibres régionaux et défis pour l'aménagement}

\begin{itemize}
\item \textbf{Surpeuplement rural à l'Ouest (Pays Bamiléké)} : Pression foncière extrême (morcellement parcellaire < 0,5 ha/exploitation), érosion des sols, exode rural massif vers les villes (Douala, Yaoundé, Bafoussam) et vers les « fronts pionniers » (Est, Sud, Adamaoua, plantations SW). Conflits fonciers autochtones/allochtones récurrents.
\item \textbf{Sous-exploitation de l'Est et du Sud} : « Vide » démographique = sous-valorisation des ressources (forêt, minerais, biodiversité, potentiel hydroélectrique \emph{Nachtigal, Lom Pangar, Memve'ele}). Difficulté à fournir des services publics (santé, éducation) à une population dispersée. Coût d'infrastructure au km² élevé.
\item \textbf{Primatie urbaine et macrocéphalie} : Douala + Yaoundé = \SI{\sim 25}{\%} de la population, \SI{\sim 60}{\%} du PIB urbain, \SI{\sim 80}{\%} de l'industrie moderne. Problèmes : habitat précaire (bidonvilles), chômage des diplômés, insécurité, gestion des déchets (HYSACAM saturé), congestion routière.
\item \textbf{Déséquilibres interrégionaux} : Indicateurs de développement (taux de scolarisation, accès à l'eau potable, électricité, route bitumée, revenu par tête) corrélés positivement avec les pôles urbains et négativement avec la faible densité (Est, Adamaoua hors axes, Extrême-Nord septentrional).
\item \textbf{Réponses de l'aménagement} : PND 2020-2030 (piliers : industrialisation, capital humain, gouvernance). Projets structurants : Autoroute Yaoundé-Douala, bitumage Est (Yokadouma), barrage de Nachtigal, port de Kribi (phase 2), corridor ferroviaire Douala-Ngaoundéré (réhabilitation), Zones Économiques Spéciales (ZES) à Kribi, Limbé, Ngaoundéré.
\end{itemize}

\begin{attention}[Pièges à éviter au Bac et en TP]
\begin{itemize}
\item \textbf{Confondre densité et effectif} : Dire « l'Est est peu peuplé » (vrai pour la densité) mais « l'Est a peu d'habitants » (faux en valeur absolue : \SI{\sim 1}{million}, plus que le Nord-Ouest \SI{\sim 0,9}{M}). \emph{Distingué toujours densité (hab/km²) et population totale (hab)}.
\item \textbf{Oublier l'urbain dans l'analyse rurale} : La densité du Littoral est tirée par Douala. Hors Douala, la densité rurale du Littoral est moyenne. Idem pour le Centre (Yaoundé). \emph{Préciser : « forte densité urbaine » vs « forte densité rurale »}.
\item \textbf{Réduire les facteurs au seul naturel} : Le déterminisme climatique/édaphique est insuffisant. L'Ouest est dense \emph{malgré} le relief accidenté, grâce à l'organisation sociale et l'histoire. L'Est est vide \emph{malgré} la forêt riche, à cause de l'histoire (traite, maladies, isolement colonial) et de l'économie (extractivisme sans peuplement).
\item \textbf{Légende de croquis incomplète} : En TP, une carte de densité sans seuils chiffrés dans la légende, sans flèche Nord, sans échelle, sans titre précis (« Répartition de la population au Cameroun en 2023 ») perd des points.
\item \textbf{Ne pas citer d'exemples précis} : « Dans le Nord » est vague. « Dans le bassin de la Bénoué (Garoua) » ou « Dans les Monts Mandara (Mokolo) » est attendu.
\end{itemize}
\end{attention}

\begin{savaistu}[Douala et Yaoundé : le duo infernal de la primatie]
Douala (capitale économique, port autonome, aéroport international, zone industrielle Bassa/Bonabéri) et Yaoundé (capitale politique, administrations centrales, ambassades, universités, hôpital de référence) concentrent \textbf{environ 8 à 9 millions d'habitants} (agglomérations 2023), soit \textbf{près de 30 \% de la population totale} du pays. Elles croissent à \SI{\sim 4--5}{\%} / an (naturel + migration nette). Le corridor Douala--Yaoundé (\SI{250}{km}, 3h de route) forme une \cle{région urbaine continue} en devenir (conurbation), concentrant infrastructures, richesse et pouvoir de décision, au détriment de l'équilibre territorial. Le projet de nouvelle capitale administrative (vers Nachtigal ou Mbandjock) reste au stade d'étude depuis les années 1980.
\end{savaistu}

\begin{exoresolu}[Commentaire de document : Carte de densité et facteurs]
\textbf{Énoncé :} À l'aide de la carte choroplèthe (Figure 1) et du diagramme en barres (Figure 2), expliquez les contrastes de densité de population entre la région de l'Ouest et la région de l'Est.

\textbf{Solution détaillée :}
\begin{enumerate}
\item \textbf{Constat chiffré (lecture des documents) :}
      \begin{itemize}
      \item Région de l'Ouest : Densité \SI{\sim 180}{hab/km²} (Figure 2, barre la plus haute après Littoral). Classe « Forte densité » (Figure 1, bleu foncé sur la carte de densité).
      \item Région de l'Est : Densité \SI{\sim 9}{hab/km²} (Figure 2, barre la plus basse avec le Sud). Classe « Faible densité » (Figure 1, couleur verte).
      \item Ratio : L'Ouest est \textbf{20 fois plus dense} que l'Est, pour une superficie 5,4 fois plus petite (\SI{109 002}{km²} vs \SI{20 248}{km²}... correction : Ouest \SI{13 892}{km²}, Est \SI{109 002}{km²} -> ratio ~7,8). L'Ouest a \SI{\sim 2,5}{M} d'hab., l'Est \SI{\sim 1}{M}.
      \end{itemize}
\item \textbf{Explication par les facteurs naturels :}
      \begin{itemize}
      \item \textbf{Ouest} : Relief montagneux (Bamiléké, 1000-2500 m) $\rightarrow$ climat tempéré (paludisme faible), sols volcaniques (andosols) très fertiles, pluviométrie élevée (> 2000 mm). \textbf{Favorable}.
      \item \textbf{Est} : Plateau pénéplané (500-800 m), forêt dense humide, sols ferrallitiques lessivés (pauvres, acides), pluviométrie forte mais lessivage, maladies vectorielles (onchocercose historique, paludisme). \textbf{Défavorable} à l'agriculture intensive et à la santé.
      \end{itemize}
\item \textbf{Explication par les facteurs historiques et humains :}
      \begin{itemize}
      \item \textbf{Ouest} : Peuple Bamiléké (Semi-Bantous), organisation en chefferies structurées (1\textsuperscript{er} et 2\textsuperscript{e} degré), agriculture intensive ancestrale (terrasses, fumure, haies vives), forte démographie endogène. Colonisation : caféiculture arabica (années 1920-30), routes précoces, éducation (missions).
      \item \textbf{Est} : Peuplement tardif et diffus (Bantous Makaa, Kako, Pygmées Baka/Bakola). Traite négrière (razzias depuis le Nord via Adamaoua par les Peuls). Colonisation : exploitation forestière/minière sans peuplement (concessions, camps temporaires, chemin de fer Douala-Ngaoundéré bypassant l'Est). Indépendance : isolement routier (piste jusqu'aux années 2000), absence de pôle urbain fort.
      \end{itemize}
\item \textbf{Explication par les facteurs économiques actuels :}
      \begin{itemize}
      \item \textbf{Ouest} : Saturation foncière $\rightarrow$ main-d'œuvre migrante vers les villes (Douala, Yaoundé) et plantations (Sud-Ouest). Marchés denses, commerce florissant, diaspora dynamique (investissements au pays).
      \item \textbf{Est} : Économie de rente (bois : SFID, SEFAC ; or/diamant : artisanal + Capam ; futur fer Mbalam/Nabeba, bauxite Minim-Martap) = enclaves, peu d'emplois locaux durables, faible marché local. Port de Kribi (Sud) et route Yaoundé-Bertoua bitumée (2010s) : début d'attractivité mais effet limité pour l'instant.
      \end{itemize}
\item \textbf{Synthèse :} Le contraste Ouest/Est est le \textbf{paradigme du déséquilibre camerounais} : surpeuplement rural historique sur terres fertiles mais étroites vs sous-peuplement récent sur terres vastes mais marginalisées par l'histoire, la géographie physique et l'économie extractive.
\end{enumerate}
\end{exoresolu}

\begin{exercice}[Entraînement : Croquis de répartition (TP 2)]
Réalisez un croquis simplifié de la répartition de la population au Cameroun sur un contour de carte fourni (ou dessiné).
\textbf{Consignes :}
\begin{enumerate}
\item Représentez les trois classes de densité (forte, moyenne, faible) par trois motifs/couleurs distincts (hachures croisées, pointillés, tramage léger).
\item Localisez et nommez : les 4 métropoles (Douala, Yaoundé, Bafoussam, Garoua/Maroua), les Monts Mandara, le Mont Cameroun, le barrage de Lagdo.
\item Tracez l'axe structurant principal (Douala--Yaoundé--Ngaoundéré) et l'axe transversal (Bafoussam--Bamenda--Mamfé).
\item Légendez soigneusement (titre, sources, échelle, orientation, signature).
\item Rédigez un commentaire de 10 lignes expliquant le contraste Est / Ouest en mobilisant 2 facteurs naturels et 2 facteurs humains.
\end{enumerate}
\end{exercice}

\begin{corrige}[Exercice n°1 : Croquis de répartition]
\textbf{Éléments attendus sur le croquis :}
\begin{itemize}
\item \textbf{Zone Forte (Rose/Hachures croisées)} : Régions Littoral, Ouest, Nord-Ouest + Monts Mandara (pointe Nord-Ouest Extrême-Nord), piémonts Atlantika (Nord), Mont Cameroun (Sud-Ouest).
\item \textbf{Zone Moyenne (Orange/Pointillés)} : Régions Centre, Sud-Ouest, Adamaoua (bande le long N1/rail), Nord (bassin Bénoué + piémonts), Extrême-Nord (plaine Logone au Sud).
\item \textbf{Zone Faible (Vert/Tramage léger)} : Régions Est, Sud, Adamaoua (Est/Nord hors axes), Extrême-Nord (bande Nord sahélienne au-delà du 11\textsuperscript{e} parallèle).
\item \textbf{Symboles ponctuels} : $\bullet$ Douala, Yaoundé, Bafoussam, Garoua, Maroua, Ngaoundéré, Bertoua, Ebolowa, Kribi, Limbé. $\triangle$ Mont Cameroun, Monts Mandara. $\sim$ Barrage Lagdo (sur Bénoué).
\item \textbf{Lignes} : Flèches épaisses pour axe Douala-Yaoundé-Ngaoundéré (rail/route), flèches moyennes pour transversal Ouest-NO (Bafoussam-Bamenda).
\item \textbf{Légende} : Organisée par thèmes (Densité, Villes, Relief, Infrastructures). Titre : « Répartition de la population au Cameroun (estimation 2023) ». Source : INS/BUCREP. Auteur / Date.
\end{itemize}
\textbf{Commentaire type (10 lignes) :}
« L'Ouest (180 hab/km²) et l'Est (9 hab/km²) illustrent le contraste majeur. \textbf{Naturellement}, l'Ouest bénéficie de sols volcaniques fertiles (andosols) et d'un climat tempéré d'altitude limitant le paludisme, tandis que l'Est subit des sols ferrallitiques pauvres, la forêt dense humide et l'onchocercose historique. \textbf{Humainement}, l'Ouest est un foyer de peuplement ancien (Bamiléké, chefferies structurées, agriculture intensive sur terrasses) et de colonisation caféière précoce, alors que l'Est a subi la traite négrière, une colonisation extractive (forêt/mine) sans peuplement et un isolement routier prolongé. Aujourd'hui, l'exode rural ouestien alimente les métropoles et les fronts pionniers, pendant que l'Est peine à attirer hors enclaves minières/forestières malgré le bitumage récent de la N10. »
\end{corrige}

\coursec{Mobilité, brassage des populations et intégration nationale}

Nous venons de voir que la population camerounaise est très inégalement répartie : l'Ouest, le Littoral et le Centre concentrent les densités les plus fortes, tandis que l'Est, le Nord et l'Adamaoua restent faiblement peuplés. Mais cette carte n'est pas figée : elle bouge en permanence. Chaque jour, des milliers de Camerounais quittent leur village pour la ville, leur région pour une autre, leurs champs pour une plantation ou un chantier. Cette dynamique de \cle{mobilité} produit un phénomène majeur pour l'intégration nationale : le \cle{brassage des populations}. C'est l'objet de cette section.

\begin{definition}[Le brassage des populations]
Le \cle{brassage des populations} est le \textbf{mélange de groupes humains différents sur un même espace}, résultant des migrations. Il peut être spatial (des ethnies différentes cohabitent dans un quartier ou un village), culturel (adoption de langues, de cuisines, de musiques d'autres groupes) ou biologique (mariages mixtes). Il ne supprime pas les identités : il les fait se rencontrer et évoluer.
\end{definition}

\courssub{Une grande mobilité interne, ancienne et multiforme}

\textbf{Intuition.} Prends le car qui relie Bafoussam à Douala un vendredi soir : il est plein de commerçants, d'étudiants, de fonctionnaires en mission. Prends maintenant la route de Foumban à Bafoussam en saison sèche : tu croiseras des troupeaux de zébus conduits par des éleveurs \cle{Bororo} en \cle{transhumance}. Ces deux scènes illustrent la même réalité : les Camerounais sont un peuple en mouvement.

\begin{attention}[Les migrations ne sont pas un phénomène récent]
Erreur fréquente en dissertation : présenter la mobilité comme une conséquence de la seule modernité. Or les mobilités sont \textbf{anciennes} au Cameroun : la transhumance pastorale des Bororo existe depuis des siècles ; le \cle{commerce caravanier} (kola, sel, esclaves puis produits de traite) reliait déjà le Nord soudanien à la côte avant la colonisation ; les migrations de conquête (peuplement fulbé de l'Adamaoua au XIX\textsuperscript{e} siècle) ont remodelé la carte humaine. La colonisation et l'indépendance ont \emph{accéléré} et \emph{réorienté} ces mouvements, elles ne les ont pas créés.
\end{attention}

On distingue plusieurs formes de migrations internes :

\begin{itemize}
\item Les \cle{migrations saisonnières} : transhumance des éleveurs Bororo (de l'Adamaoua vers le Nord-Ouest et l'Ouest en saison sèche), départs temporaires des jeunes du Nord vers les villes du Sud pendant la saison sèche (phénomène parfois appelé « exode de saison sèche »).
\item L'\cle{exode rural} : départ définitif des jeunes des campagnes vers les villes, à la recherche d'emplois, d'études ou de distractions. Il alimente la croissance rapide de Douala et Yaoundé et vieillit les campagnes.
\item Les \cle{migrations de travail} : historiquement vers les grandes plantations du littoral (CDC à Tiko, Mondoni, Idenau ; PHP à Penja ; SOCAPALM à Dibombari, Mbongo, Eseka, Kienké ; SABC à Douala), les chantiers du port de Douala, aujourd'hui aussi vers les chantiers de construction, les mines (ex. : cobalt de Nkamouna, or de Bétaré-Oya) et les services urbains.
\end{itemize}

\courssub{Les grands courants migratoires internes}

L'observation des cartes de densité et des enquêtes de recensement (RGPH 2005, 2010) permet de dégager trois grands courants :

\begin{enumerate}
\item \textbf{Du Nord vers le Sud} : les \cle{Fulbé} (Peuls, groupe néo-soudanais) et les commerçants \cle{Haoussa} descendent de l'Extrême-Nord, du Nord et de l'Adamaoua vers les villes du Centre, du Littoral et de l'Ouest, où ils dominent certains commerces (bétail, tissus, change, quincaillerie). Les quartiers comme Briqueterie à Yaoundé ou New-Bell à Douala en témoignent.
\item \textbf{De l'Ouest vers le Littoral et le Centre} : les \cle{Bamiléké}, issus des hautes terres très densément peuplées de l'Ouest, migrent massivement vers Douala (commerce, transport, emplois portuaires) et Yaoundé, mais aussi vers les campagnes du Mbam, de la Menoua et du Noun comme \cle{colons agricoles}.
\item \textbf{Du Sud vers les villes} : les populations du Centre, du Sud et de l'Est (Ewondo, Bulu, Bassa, Bafia...) alimentent l'exode rural vers Yaoundé, Douala, Édéa, Kribi, attirées par l'administration, les études et les industries.
\end{enumerate}

\begin{popfigure}
\begin{center}
\includegraphics[width=\linewidth,height=10cm,keepaspectratio]{N02/cmr_regions.jpg}
\end{center}
\legende{Fond de carte pour les principaux courants migratoires internes au Cameroun : grands ensembles régionaux (nord, Adamaoua, hautes terres de l'Ouest, littoral, plateau sud), villes principales (Douala, Yaoundé, Bamenda, Bafoussam, Maroua, Garoua, Ngaoundéré), fleuves et routes (noms en anglais). Les flux à y lire : de l'Extrême-Nord et des hautes terres de l'Ouest vers Douala et Yaoundé, des campagnes vers les chefs-lieux, vers les fronts pionniers du Centre et de l'Est. \\ Source du fond de carte : Wikimedia Commons, Peter Fitzgerald, CC BY 3.0.}
\end{popfigure}

\courssub{Le brassage dans les villes}

Les villes sont les grands creusets du brassage. À \cle{Douala} et \cle{Yaoundé}, aucun groupe ethnique ne vit seul : les quartiers sont \cle{multiethniques}. New-Bell, Deïdo, Akwa ou Bonabéri à Douala ; Briqueterie, Mvog-Ada, Nkolndongo ou Essos à Yaoundé accueillent des populations venues de toutes les régions.

\begin{exemplebox}[Des quartiers-mosaïques]
Dans certains quartiers de Douala, on compte des ressortissants de \textbf{plus de 100 groupes ethniques différents} (ordre de grandeur couramment cité, le Cameroun comptant environ 250 groupes ethniques). Un même immeuble peut abriter un commerçant bamiléké, un fonctionnaire ewondo, un boucher haoussa, un docker bassa et une couturière bafia.
\end{exemplebox}

Ce brassage urbain prend plusieurs formes :

\begin{itemize}
\item Les \cle{mariages mixtes}, de plus en plus fréquents en ville, qui créent des familles et des enfants à double ou triple appartenance culturelle.
\item Les \cle{langues véhiculaires} qui permettent la communication entre groupes : le \cle{français} et l'\cle{anglais} (langues officielles), le \cle{pidgin} (au Littoral, au Sud-Ouest et au Nord-Ouest), le \cle{fulfuldé} (au Nord), l'\cle{ewondo} (au Centre) et le « \cle{camfranglais} » (parler mixte des jeunes citadins).
\item Les associations de développement et tontines, souvent ethniques à l'origine, qui s'ouvrent progressivement aux voisins d'autres origines.
\end{itemize}

\courssub{Le brassage dans les campagnes}

Le brassage ne concerne pas que les villes. Dans les campagnes, deux figures dominent :

\begin{itemize}
\item Le \cle{colon agricole} : migrant venu d'une région surpeuplée qui défriche et cultive une terre ailleurs. L'exemple classique est celui des \cle{Bamiléké} installés dans la \cle{Menoua} (autour de Dschang), dans le \cle{Mbam} (Centre) et dans le Noun (autour de Foumban), où ils ont développé des cultures vivrières et de rente (maraîchage, maïs, café). Leur réussite économique est parfois source d'admiration, parfois de tensions foncières avec les populations autochtones.
\item L'\cle{éleveur transhumant} : les \cle{Bororo} (Peuls pasteurs) conduisent leurs zébus de l'Adamaoua vers les pâturages du Nord-Ouest et de l'Ouest en saison sèche. Cette transhumance, vieille de plusieurs siècles, met en contact éleveurs et agriculteurs sédentaires. Les populations \cle{Pygmées} (Baka, Bakola, Bagyéli) de l'Est et du Sud, bien que moins mobiles, sont aussi concernées par les politiques de sédentarisation et de scolarisation qui favorisent leur brassage progressif avec les populations bantoues voisines.
\end{itemize}

\courssub{Atouts et défis du brassage pour l'intégration nationale}

\textbf{Les atouts.} Le brassage est d'abord une richesse :

\begin{itemize}
\item \textbf{Échanges économiques} : les migrations mettent en relation des espaces complémentaires (la viande du Nord, les vivres de l'Ouest, le poisson du Littoral) ; les migrants créent des réseaux commerciaux nationaux et envoient des transferts d'argent vers leurs villages d'origine.
\item \textbf{Enrichissement culturel} : musiques (mangambeu, bikutsi, makossa devenus nationaux), cuisines, langues et savoir-faire circulent et se mélangent.
\item \textbf{Sentiment national} : la cohabitation quotidienne, l'école, l'armée, le service civique national et les mariages mixtes construisent une identité camerounaise au-delà des appartenances ethniques.
\end{itemize}

\textbf{Les limites et tensions.} Le brassage n'est pas toujours harmonieux :

\begin{itemize}
\item Les \cle{conflits agriculteurs-éleveurs}, notamment au \cle{Nord-Ouest}, où les troupeaux en transhumance détruisent parfois les champs, provoquant des affrontements meurtriers.
\item La \cle{concurrence foncière} entre autochtones et colons agricoles (Menoua, Mbam, Noun), qui peut dégénérer en conflits violents.
\item Les \cle{stéréotypes ethniques} et le \cle{régionalisme} (« tel groupe domine tel secteur »), qui alimentent la méfiance, les discriminations à l'embauche et parfois les discours de haine, amplifiés par les réseaux sociaux.
\end{itemize}

\textbf{Les réponses de l'État.} Face à ces défis, l'État camerounais promeut plusieurs politiques : la proclamation de l'\cle{unité nationale} (fête nationale du 20~mai, hymne, drapeau), le \cle{bilinguisme} français-anglais comme facteur d'unité, la \cle{décentralisation} (régions, communes) pour rapprocher la décision des populations, la recherche de l'\cle{équilibre régional} dans les nominations aux postes publics (principe constitutionnel), et la promotion du « \cle{vivre ensemble} » comme valeur cardinale.

\begin{popfigure}
\begin{center}
\begin{tikzpicture}[node distance=0.35cm, every node/.style={font=\small}]
\node[draw=popdark, fill=popblueL, rounded corners, thick, minimum width=4.6cm, minimum height=1.1cm, align=center] (centre) {\textbf{Brassage des populations}\\ moteur de l'intégration nationale};
\node[draw=popgreen, fill=popgreenL, rounded corners, thick, left=2.2cm of centre, minimum width=4.6cm, align=center] (atouts) {\textbf{Atouts}\\ échanges économiques\\ enrichissement culturel\\ sentiment national};
\node[draw=poporange, fill=poporangeL, rounded corners, thick, right=2.2cm of centre, minimum width=4.6cm, align=center] (defis) {\textbf{Défis}\\ conflits agriculteurs-éleveurs\\ concurrence foncière\\ stéréotypes, régionalisme};
\node[draw=poppurple, fill=poppurpleL, rounded corners, thick, below=1.1cm of centre, minimum width=7.5cm, align=center] (etat) {\textbf{Réponses de l'État}\\ unité nationale, bilinguisme, décentralisation,\\ équilibre régional, « vivre ensemble »};
\draw[->, thick, popgreen] (atouts.east) -- (centre.west);
\draw[->, thick, poporange] (defis.west) -- (centre.east);
\draw[->, thick, poppurple] (centre.south) -- (etat.north);
\draw[->, thick, poppurple, dashed] (etat.west) .. controls (-6.2,-2.6) .. (atouts.south west);
\draw[->, thick, poppurple, dashed] (etat.east) .. controls (6.2,-2.6) .. (defis.south east);
\end{tikzpicture}
\end{center}
\legende{Schéma bilan : le brassage des populations, entre atouts et défis, et les réponses de l'État pour en faire un levier d'intégration nationale.}
\end{popfigure}

\begin{methode}[Structurer une dissertation sur l'intégration nationale]
Pour un sujet du type « Le brassage des populations au Cameroun : atout ou menace pour l'intégration nationale ? », structure ton devoir en trois temps :
\begin{enumerate}
\item \textbf{Diversité (constat)} : le Cameroun compte environ 250 groupes ethniques répartis en grands ensembles (soudanais : paléo-soudanais et néo-soudanais ; bantou ; semi-bantou ; pygmées), inégalement répartis.
\item \textbf{Brassage (dynamique)} : migrations anciennes et actuelles, grands courants (Nord$\rightarrow$Sud, Ouest$\rightarrow$Littoral/Centre, Campagnes$\rightarrow$Villes), mélange dans les villes et les campagnes.
\item \textbf{Unité (enjeu et politiques)} : atouts (économie, culture, sentiment national), tensions (conflits, foncier, régionalisme), et réponses de l'État (unité nationale, bilinguisme, équilibre régional, vivre ensemble).
\end{enumerate}
Conclus en montrant que le brassage est un \emph{processus} : il peut renforcer l'unité s'il est accompagné par des politiques justes et inclusives.
\end{methode}

\begin{exoresolu}[Analyse de document : la transhumance au Nord-Ouest]
\textbf{Énoncé.} Document : « Chaque saison sèche, des milliers de zébus conduits par des éleveurs bororo quittent l'Adamaoua pour les pâturages du Nord-Ouest. Des conflits éclatent régulièrement avec les agriculteurs dont les champs sont endommagés. Des comités de dialogue "agropasteurs" ont été créés pour délimiter les couloirs de transhumance. » Questions : 1) Identifie les deux groupes en présence et la cause du conflit. 2) Montre que ce cas illustre à la fois le brassage et ses limites. 3) Propose une solution.
\tcblower
\textbf{Solution détaillée.}
1) Les deux groupes sont les \textbf{éleveurs bororo} (Peuls transhumants, groupe néo-soudanais) et les \textbf{agriculteurs sédentaires} du Nord-Ouest (groupes semi-bantou). La cause du conflit est la \textbf{concurrence pour l'espace} : les troupeaux détruisent les cultures, et les agriculteurs empiètent parfois sur les pâturages et couloirs de passage.
2) Ce cas illustre le \textbf{brassage} car la transhumance met en contact durable deux groupes humains différents sur un même espace, avec des échanges économiques réels (lait, fumure des champs, viande). Mais il en montre aussi les \textbf{limites} : la cohabitation peut dégénérer en violences lorsque les règles d'usage de l'espace ne sont pas claires ou respectées.
3) Solution : la \textbf{délimitation concertée des couloirs de transhumance} et des zones de pâturage par des comités mixtes (éleveurs, agriculteurs, autorités administratives et traditionnelles), assortie d'un calendrier de passage et de mécanismes d'indemnisation des dégâts. On peut y ajouter la sensibilisation au « vivre ensemble » et l'application équitable de la loi foncière (loi n°74-1 du 6 juillet 1974).
\end{exoresolu}

\begin{savaistu}[Le « vivre ensemble », un concept au cœur du Bac]
L'expression « vivre ensemble » est devenue un mot d'ordre officiel au Cameroun, régulièrement rappelé par les pouvoirs publics face aux tensions identitaires. Au Baccalauréat, le conflit agriculteurs-éleveurs du Nord-Ouest et la politique d'équilibre régional sont des \textbf{exemples fréquemment attendus} dans les dissertations sur l'intégration nationale : retiens-les, avec leurs mécanismes et leurs réponses.
\end{savaistu}

\begin{aretenir}[L'essentiel de la section]
\begin{itemize}
\item La mobilité interne est \textbf{ancienne} (transhumance, commerce caravanier) et prend des formes variées : migrations saisonnières, exode rural, migrations de travail.
\item Trois grands courants : Nord $\rightarrow$ Sud (Fulbé, Haoussa), Ouest $\rightarrow$ Littoral/Centre (Bamiléké), campagnes $\rightarrow$ villes.
\item Le \textbf{brassage} se réalise dans les villes (quartiers multiethniques, mariages mixtes, langues véhiculaires) et dans les campagnes (colons agricoles, éleveurs transhumants).
\item Il est un \textbf{atout} (économie, culture, sentiment national) mais aussi un \textbf{défi} (conflits, foncier, stéréotypes).
\item L'État répond par l'\textbf{unité nationale}, le \textbf{bilinguisme}, la \textbf{décentralisation}, l'\textbf{équilibre régional} et la promotion du \textbf{vivre ensemble}.
\end{itemize}
\end{aretenir}

\begin{exercice}[Pour t'entraîner]
1) Définis le brassage des populations et donne deux exemples camerounais (un urbain, un rural). 2) Sur un croquis du Cameroun, trace et légende les trois grands courants migratoires internes. 3) Rédige un paragraphe argumenté répondant à la question : « Le brassage des populations renforce-t-il l'unité nationale camerounaise ? »
\end{exercice}

\begin{corrige}[Exercice]
1) Le brassage des populations est le mélange de groupes humains différents sur un même espace, résultant des migrations. Exemple urbain : les quartiers multiethniques de Douala (plus de 100 groupes ethniques représentés dans certains quartiers). Exemple rural : les colons agricoles bamiléké installés dans la Menoua ou le Mbam.
2) Le croquis doit comporter : le contour simplifié du Cameroun, les villes repères (Douala, Yaoundé, Bafoussam, Garoua, Maroua), une flèche épaisse de l'Ouest vers le Littoral, une flèche moyenne du Nord vers le Sud, des flèches fines des campagnes du Sud vers les villes, une flèche du nord et une légende numérotée.
3) Paragraphe attendu : oui, à condition d'être encadré. Le brassage crée des solidarités économiques (réseaux commerciaux, complémentarités régionales), culturelles (langues véhiculaires, musiques nationales) et matrimoniales (mariages mixtes) qui dépassent les appartenances ethniques. Mais il peut aussi nourrir des tensions (conflits agriculteurs-éleveurs, concurrence foncière, régionalisme). C'est pourquoi les politiques d'unité nationale, de bilinguisme, d'équilibre régional et de « vivre ensemble » sont indispensables pour transformer la diversité brassée en force nationale.
\end{corrige}

\coursec{TP 2 : Construction et lecture de la carte de répartition de la population}

Dans les sections précédentes, nous avons vu que les Camerounais ne restent pas figés dans leurs terroirs d'origine : migrations vers les villes, brassage des groupes humains, peuplement des fronts pionniers. Mais une question demeure : comment \emph{voir} ces réalités d'un seul coup d'œil ? Comment montrer, sur une seule carte, que l'Ouest étouffe sous la densité humaine pendant que l'Est reste presque vide ? C'est exactement le travail du géographe — et ce sera le tien dans ce TP : construire toi-même la carte de la répartition de la population camerounaise, puis la lire comme un document de diagnostic.

\courssub{Objectifs et matériel du TP}

\begin{definition}[Carte choroplèthe]
Une \cle{carte choroplèthe} est une carte thématique où des unités territoriales (régions, départements, communes) sont colorées selon la valeur d'un phénomène mesuré (ici la \cle{densité de population}), en utilisant une gamme de couleurs dégradées : plus la valeur est forte, plus la couleur est foncée.
\end{definition}

\begin{aretenir}[Objectifs du TP 2]
À l'issue de ce TP, tu dois être capable de :
\begin{enumerate}
\item calculer une \cle{densité de population} à partir de données brutes ;
\item définir des \cle{classes de densité} et choisir une gamme de couleurs cohérente ;
\item construire et légender une carte choroplèthe du Cameroun ;
\item lire et interpréter cette carte (décrire, expliquer, conclure) ;
\item prolonger l'analyse par un diagramme en barres et un croquis des courants migratoires.
\end{enumerate}
\end{aretenir}

\textbf{Matériel nécessaire.} Un fond de carte du Cameroun vierge avec les limites des dix régions ; le tableau statistique des populations et superficies par région (fourni ci-dessous) ; une calculatrice ; trois crayons de couleur d'une même gamme (par exemple jaune clair, orange, rouge foncé) ; un crayon noir pour la légende et les toponymes.

\courssub{Étape 1 : Calculer les densités par région}

La densité de population est le rapport entre le nombre d'habitants et la superficie du territoire qu'ils occupent :
\[ \text{Densité (hab/km}^2\text{)} = \frac{\text{Population (habitants)}}{\text{Superficie (km}^2\text{)}} \]

Voici le tableau statistique de travail (données issues des projections officielles du BUCREP à partir du RGPH 2005 ; les chiffres sont des ordres de grandeur arrondis, suffisants pour un travail de classe).

\begin{center}
\begin{tabularx}{\linewidth}{|l|X|X|X|X|}
\hline
\rowcolor{popblueL} \textbf{Région} & \textbf{Population (millions)} & \textbf{Superficie (km$^2$)} & \textbf{Densité (hab/km$^2$)} & \textbf{Classe} \\
\hline
Adamaoua & 1,2 & \num{63701} & $\approx 19$ & Faible \\
\hline
Centre (avec Yaoundé) & 4,1 & \num{68953} & $\approx 59$ & Moyenne \\
\hline
Est & 0,8 & \num{109002} & $\approx 7$ & Faible \\
\hline
Extrême-Nord & 3,9 & \num{34246} & $\approx 114$ & Forte \\
\hline
Littoral (avec Douala) & 3,3 & \num{20239} & $\approx 163$ & Forte \\
\hline
Nord & 2,4 & \num{65576} & $\approx 37$ & Moyenne \\
\hline
Nord-Ouest & 1,9 & \num{17812} & $\approx 107$ & Forte \\
\hline
Ouest & 1,9 & \num{13872} & $\approx 137$ & Forte \\
\hline
Sud & 0,7 & \num{47110} & $\approx 15$ & Faible \\
\hline
Sud-Ouest & 1,5 & \num{25410} & $\approx 59$ & Moyenne \\
\hline
\textbf{Cameroun} & \textbf{21,7} & \textbf{\num{475440}} & \textbf{$\approx 46$ (RGPH 2005) / $\approx 55$ (estimation récente)} & --- \\
\hline
\end{tabularx}
\end{center}

\begin{exoresolu}[Calculer une densité : le cas de l'Ouest]
Énoncé. La région de l'Ouest compte environ 1,9 million d'habitants sur \num{13872} km$^2$. Calcule sa densité et compare-la à celle de l'Est (0,8 million d'habitants sur \num{109002} km$^2$).
\tcblower
Solution détaillée.
\begin{align*}
\text{Densité de l'Ouest} &= \frac{\num{1900000}}{\num{13872}} \approx 137 \text{ hab/km}^2 \\
\text{Densité de l'Est} &= \frac{\num{800000}}{\num{109002}} \approx 7 \text{ hab/km}^2
\end{align*}
Conclusion : la densité de l'Ouest est environ \textbf{vingt fois} supérieure à celle de l'Est. Sur une même superficie, l'Ouest concentre vingt fois plus d'hommes : c'est un contraste majeur de peuplement que la carte devra rendre visible.
\end{exoresolu}

\begin{exemplebox}[La moyenne nationale masque les contrastes]
La densité moyenne du Cameroun était d'environ \cle{46 hab/km$^2$} au RGPH 2005 ; les estimations récentes la situent autour de \cle{55 à 60 hab/km$^2$}. Mais cette moyenne est un \emph{mirage statistique} : aucune région ne vit réellement à cette densité. Le Littoral dépasse 160 hab/km$^2$ tandis que l'Est tombe sous 10 hab/km$^2$. Retiens ce réflexe de géographe : \textbf{une moyenne nationale doit toujours être décomposée} pour révéler les disparités territoriales.
\end{exemplebox}

\courssub{Étape 2 : Définir les classes et choisir les couleurs}

Une fois les densités calculées, il faut les \emph{classer}. Nous adoptons ici un découpage en trois classes, simple et lisible :

\begin{center}
\begin{tabularx}{\linewidth}{|l|X|X|}
\hline
\rowcolor{popblueL} \textbf{Classe} & \textbf{Densité (hab/km$^2$)} & \textbf{Couleur proposée} \\
\hline
Zone densément peuplée & plus de 100 & rouge foncé (ou brun) \\
\hline
Zone moyennement peuplée & 20 à 100 & orange (ou rose) \\
\hline
Zone faiblement peuplée & moins de 20 & jaune clair \\
\hline
\end{tabularx}
\end{center}

\begin{methode}[Règles d'or d'une légende de carte choroplèthe]
Une légende de carte choroplèthe doit toujours comporter :
\begin{enumerate}
\item un \cle{titre} précis (phénomène, espace, date) : « Densités de population par région au Cameroun » ;
\item des \cle{classes ordonnées}, de la plus faible à la plus forte, avec leurs bornes chiffrées et leur unité (hab/km$^2$) ;
\item des \cle{couleurs dégradées} d'une même gamme, du clair au foncé, l'intensité croissant avec la valeur ;
\item les \cle{sources} des données (ex. : RGPH 2005, BUCREP) et, sur la carte, l'\cle{orientation} (flèche du nord) et une \cle{échelle} indicative.
\end{enumerate}
\end{methode}

\begin{attention}[Erreur fréquente : les couleurs sans logique ordinale]
Ne colorie jamais les régions avec des couleurs choisies au hasard (vert ici, bleu là, rouge ailleurs) : le lecteur ne peut plus comparer. En choroplèthe, \textbf{l'intensité de la couleur doit croître avec la valeur} : clair = faible densité, foncé = forte densité. Autre piège : oublier les bornes chiffrées des classes dans la légende — une couleur sans chiffre ne signifie rien.
\end{attention}

\courssub{Étape 3 : Colorier la carte et rédiger la légende complète}

Reporte chaque région dans sa classe : forte densité (Littoral, Ouest, Extrême-Nord, Nord-Ouest) ; densité moyenne (Centre, Sud-Ouest, Nord) ; faible densité (Adamaoua, Sud, Est). Colorie ensuite chaque région avec la couleur de sa classe, écris les noms des régions et des principales villes, trace la flèche du nord et rédige la légende. Voici le modèle attendu :

\begin{popfigure}
\begin{center}
\includegraphics[width=\linewidth,height=11cm,keepaspectratio]{N02/cmr_dens.png}
\end{center}
\legende{Exemple de carte choroplèthe de densités par région au Cameroun (quatre classes, légende en anglais). 1 : zones densément peuplées, plus de 50 hab/km\textsuperscript{2} (Littoral, Ouest, Nord-Ouest, Extrême-Nord) ; 2 : zones moyennement peuplées, 10 à 50 (Centre, Sud-Ouest, Nord) ; 3 : zones faiblement peuplées, moins de 10 (Adamaoua, Sud, Est). \\ Source : Wikimedia Commons, Austiger, CC BY-SA 4.0.}
\end{popfigure}

\courssub{Étape 4 : Lecture et interprétation de la carte}

Construire une carte ne suffit pas : il faut la faire \emph{parler}.

\begin{methode}[Lire une carte de densités en trois étapes]
\begin{enumerate}
\item \textbf{Décrire (où ?)} : localise les trois types de zones avec un vocabulaire d'orientation précis (littoral, ouest, nord, sud-est...). Cite les régions et les chiffres.
\item \textbf{Expliquer (pourquoi ?)} : relie chaque fait à ses causes physiques (relief, climat, sols, maladies) et humaines (histoire du peuplement, activités économiques, villes, axes de transport, migrations).
\item \textbf{Conclure (conséquences ?)} : dégage les effets sur l'aménagement du territoire, les besoins des populations et les flux migratoires.
\end{enumerate}
\end{methode}

\textbf{Application.} La description montre un Cameroun à trois vitesses : un \emph{arc de forte densité} allant du littoral (Douala) aux hautes terres de l'Ouest et du Nord-Ouest, prolongé par la cuvette de l'Extrême-Nord autour de Maroua ; une \emph{ceinture intermédiaire} (Centre avec Yaoundé, Nord, Sud-Ouest) ; un \emph{vaste sud-est presque vide} (Est, Sud, Adamaoua). L'explication croise les facteurs : sols volcaniques fertiles et climat frais des hautes terres bamiléké et bamoun ; port et industries de Douala ; agriculture de sécurité et vieux peuplement du bassin du Logone ; à l'inverse, forêt équatoriale dense, sols lessivés, enclavement et faible accessibilité dans l'Est et le Sud.

\courssub{Étape 5 : Rédiger un court commentaire géographique}

Le commentaire suit la logique \emph{constat — explications — conséquences}. Modèle de rédaction attendu :

\begin{exemplebox}[Exemple de commentaire rédigé]
\textbf{Constat.} La répartition de la population camerounaise est profondément inégale : avec une moyenne nationale d'environ 55 hab/km$^2$ (estimation récente), le Littoral atteint 163 hab/km$^2$ et l'Ouest 137 hab/km$^2$, tandis que l'Est tombe à 7 hab/km$^2$ et le Sud à 15 hab/km$^2$.
\textbf{Explications.} Les fortes densités s'expliquent par la fertilité des sols volcaniques des hautes terres de l'Ouest, l'attraction du port et des industries de Douala, et l'ancienneté du peuplement de l'Extrême-Nord ; les faibles densités tiennent à la forêt dense, aux sols pauvres et à l'enclavement du Sud-Est.
\textbf{Conséquences.} Ces contrastes alimentent de puissantes migrations vers les villes et les zones denses, posent le problème de la pression foncière à l'Ouest et celui de la mise en valeur des régions vides : l'État doit y concentrer routes, écoles et centres de santé pour équilibrer le territoire.
\end{exemplebox}

\begin{savaistu}[La carte de densité, un outil de planification]
La carte que tu viens de construire n'est pas un exercice scolaire sans lendemain : c'est un véritable \cle{instrument de décision}. Les services de l'État (MINEPAT, BUCREP) s'appuient sur les cartes de densité pour décider \emph{où} implanter un lycée, un hôpital de district, un marché ou une route nationale. Une région dense exige plus d'infrastructures sociales ; une région vide exige d'abord des axes de désenclavement. Le géographe éclaire ainsi l'action de l'aménageur.
\end{savaistu}

\courssub{Prolongement : diagramme en barres et croquis des migrations}

Pour comparer d'un seul regard les dix régions, traduis le tableau en diagramme en barres :

\begin{popfigure}
\begin{center}
\begin{tikzpicture}
\begin{axis}[
    ybar, bar width=9pt,
    width=0.95\linewidth, height=6.2cm,
    ymin=0, ymax=185,
    ylabel={Densité (hab/km$^2$)},
    symbolic x coords={Littoral, Ouest, E.-Nord, N.-Ouest, Centre, S.-Ouest, Nord, Adamaoua, Sud, Est},
    xtick=data, x tick label style={rotate=35, anchor=east, font=\scriptsize},
    ymajorgrids, grid style={popline!40},
    every axis plot/.append style={fill=popblue, draw=popdark}
]
\addplot coordinates {(Littoral,163) (Ouest,137) (E.-Nord,114) (N.-Ouest,107) (Centre,59) (S.-Ouest,59) (Nord,37) (Adamaoua,19) (Sud,15) (Est,7)};
\draw[popink, thick, dashed] (axis cs:Littoral,55) -- (axis cs:Est,55) node[pos=0.15, above, font=\scriptsize, popink] {moyenne nationale $\approx$ 55};
\end{axis}
\end{tikzpicture}
\end{center}
\legende{Diagramme en barres des densités régionales (ordres de grandeur, projections issues du RGPH 2005). La ligne pointillée matérialise la moyenne nationale (estimation récente) : cinq régions sont au-dessus, cinq en dessous — preuve visuelle du déséquilibre du peuplement.}
\end{popfigure}

Achève le TP par un croquis des \cle{courants migratoires} : sur un second fond de carte, trace des flèches d'épaisseur proportionnelle à l'importance des flux — flèches épaisses des régions denses rurales (Ouest, Nord-Ouest, Extrême-Nord) vers Douala et Yaoundé, flèches moyennes vers les fronts pionniers agricoles (périphérie de Douala, Mbam-et-Kim, plaine des Mbo), flèches fines vers les villes régionales. Ce croquis relie le constat statistique du TP à la dynamique humaine étudiée dans la section précédente : les hommes se déplacent des zones saturées vers les pôles d'emploi et les terres disponibles, tissant ainsi, par le brassage des populations, le vivre-ensemble national.

\begin{popfigure}
\begin{center}
\includegraphics[width=\linewidth,height=10cm,keepaspectratio]{N02/cmr_regions.jpg}
\end{center}
\legende{Fond de carte pour le croquis des principaux courants migratoires internes au Cameroun (noms en anglais). Le croquis attendu y trace : en flèches épaisses, l'exode des hautes terres de l'Ouest et du Nord-Ouest (Bamenda, Bafoussam) et de l'Extrême-Nord vers Douala et Yaoundé ; en flèches moyennes, la progression vers les fronts pionniers agricoles (Mbam, Mbo) ; en flèches fines, les déplacements vers les chefs-lieux régionaux. \\ Source du fond de carte : Wikimedia Commons, Peter Fitzgerald, CC BY 3.0.}
\end{popfigure}

\begin{aretenir}[L'essentiel du TP 2]
\begin{itemize}
\item Densité = population $\div$ superficie ; la moyenne nationale ($\approx 55$ hab/km$^2$, estimation récente) masque des contrastes extrêmes (163 au Littoral, 7 dans l'Est).
\item Une carte choroplèthe exige des classes ordonnées, des couleurs dégradées du clair au foncé, un titre, des sources, le nord et une échelle.
\item La lecture d'une carte suit toujours trois temps : décrire, expliquer, conclure.
\item Trois zones de peuplement : l'arc dense Ouest--Littoral--Extrême-Nord, la ceinture moyenne, le Sud-Est quasi vide.
\item Les migrations internes (exode rural, fronts pionniers) sont la réponse dynamique à ces déséquilibres de peuplement.
\end{itemize}
\end{aretenir}

\begin{exercice}[Pour t'entraîner]
À partir du tableau statistique du TP :
\begin{enumerate}
\item Calcule la densité de la région du Nord.
\item Classe les dix régions dans les trois catégories (forte, moyenne, faible).
\item Rédige en dix lignes un commentaire de la carte construite, en suivant le plan constat--explications--conséquences.
\end{enumerate}
\end{exercice}

\begin{corrige}[Exercice du TP 2]
\begin{enumerate}
\item Densité du Nord : $\frac{\num{2400000}}{\num{65576}} \approx 37$ hab/km$^2$ (classe moyenne).
\item Forte densité : Littoral (163), Ouest (137), Extrême-Nord (114), Nord-Ouest (107) ; Moyenne densité : Centre (59), Sud-Ouest (59), Nord (37) ; Faible densité : Adamaoua (19), Sud (15), Est (7).
\item Le commentaire doit : localiser les trois zones avec chiffres à l'appui ; invoquer au moins deux facteurs physiques (sols volcaniques, forêt dense, climat) et deux facteurs humains (port de Douala, enclavement, migrations, histoire du peuplement) ; conclure sur les conséquences pour l'aménagement (pression foncière à l'Ouest, désenclavement du Sud-Est, implantation des infrastructures, intégration nationale).
\end{enumerate}
\end{corrige}

\newpage

\coursec{Activité d'intégration}

\coursec{Leçon 2 : Les grands groupes humains}

\courssub{Introduction : définitions et cadre d'étude}
\begin{definition}[Groupe humain, peuple, ethnie]
Un \textbf{groupe humain} est un ensemble d'individus partageant une communauté de langue, de culture, d'histoire et souvent d'espace. Au Cameroun, on recense environ \textbf{250 groupes ethniques} (ou peuples) regroupés en \textbf{quatre grandes familles linguistiques et culturelles} : les \cle{Soudanais} (au Nord), les \cle{Bantou} (au Centre, Sud, Est, Littoral, Ouest, Sud-Ouest), les \cle{Semi-Bantou} (aux Hauts Plateaux de l'Ouest et du Nord-Ouest) et les \cle{Pygmées} (forêts de l'Est et du Sud).
\end{definition}

\begin{savaistu}[Le peuplement du Cameroun : une mosaïque ancienne]
Le territoire camerounais est un carrefour de migrations séculaires. Les données archéologiques (Shum Laka, Obobogo) attestent une présence humaine depuis plus de 30 000 ans. Les expansions \cle{bantoues} (vers 3000-2000 av. J.-C. depuis la région du lac Tchad/Nigeria actuel) et les poussées \cle{soudanaises} (venant du Nord-Est) ont structuré la carte ethno-linguistique actuelle. La colonisation (allemande, puis française/britannique) a figé des frontières administratives coupant parfois des ensembles ethniques unitaires (ex. : les Mafa entre Cameroun et Nigeria, les Bassa entre Cameroun et RCA).
\end{savaistu}

\courssub{Les groupes soudanais : paléo-soudanais et néo-soudanais}
\begin{aretenir}[Localisation et caractéristiques des groupes soudanais]
\begin{itemize}
    \item \textbf{Aire géographique :} Régions de l'\textbf{Extrême-Nord}, du \textbf{Nord} et de l'\textbf{Adamaoua} (bassin du Lac Tchad, plaine du Diamaré, piémonts des Mandara, plateau de l'Adamaoua, bassin de la Bénoué).
    \item \textbf{Paléo-soudanais (peuplements anciens) :} Installés de longue date dans les montagnes (Mandara, Alantika) et les plaines (Diamaré, Logone). \textbf{Exemples :} \cle{Mafa}, \cle{Mofu}, \cle{Kapsiki}, \cle{Toupouri}, \cle{Massa}, \cle{Moundang}, \cle{Musey}. Organisation en \textbf{chefferies traditionnelles} ou sociétés segmentaires. Agriculture pluviale intensive (mil, sorgho, maïs) sur terrasses (Mandara) ou en décrue (Logone). Islamisation ancienne (XI\textsuperscript{e}-XV\textsuperscript{e} s.) pour certains, christianisation récente pour d'autres.
    \item \textbf{Néo-soudanais (peuplements plus récents, XVIII\textsuperscript{e}-XIX\textsuperscript{e} s.) :} Principalement les \cle{Peuls} (ou \cle{Fulbé}, \cle{Mbororo}) venus du Fouta-Djalal (Guinée) et du Macina (Mali). \textbf{Deux modes de vie :} \emph{Peuls sédentaires} (intégrés aux chefferies, agriculture/élevage, islamisés, langue fulfulde) et \cle{Peuls nomades (Mbororo)} (pastoralisme extensif, transhumance Nord-Sud, habitat mobile, forte identité culturelle). Présence forte dans l'Adamaoua, le Nord, l'Extrême-Nord et de plus en plus au Centre/Est.
\end{itemize}
\end{aretenir}

\begin{exemplebox}[Organisation socio-politique : Le Lamidat]
Hérité du jihad d'Ousman dan Fodio (début XIX\textsuperscript{e} s.), le \textbf{lamidat} est une structure politico-religieuse (émirat) dirigée par un \textbf{Lamido} (chef suprême, autorité spirituelle et temporelle). Il s'appuie sur une administration centrale (Waziri, Sarkin Yaki, etc.) et des chefs de cantons/villages (Lawans). Exemples majeurs : \textbf{Rey-Bouba}, \textbf{Tibati}, \textbf{Ngaoundéré}, \textbf{Maroua}, \textbf{Mokolo}. Cette structure a facilité l'administration coloniale (indirect rule) et persiste comme relais de l'État moderne et gardien de la coutume.
\end{exemplebox}

\courssub{Les groupes bantou et semi-bantou ; les pygmées}
\begin{aretenir}[Répartition et diversité des groupes bantou et semi-bantou]
\begin{itemize}
    \item \textbf{Bantou (famille Niger-Congo, branche Bantoïde) :} Occupent le \textbf{Centre}, le \textbf{Sud}, l'\textbf{Est}, le \textbf{Littoral}, le \textbf{Sud-Ouest} et une partie de l'\textbf{Ouest}.
        \begin{itemize}
            \item \textit{Zone forêt (Est/Sud/Centre) :} \cle{Beti-Fang} (Ewondo, Boulou, Fang, Eton), \cle{Bassa}, \cle{Bakoko}, \cle{Makaa-Njem}, \cle{Kako}, \cle{Konabembe}. Organisation en \textbf{lignages patrilinéaires}, villages de case groupées, sociétés secrètes (\cle{Ngil}, \cle{So}, \cle{Mvet}).
            \item \textit{Zone côtière et montagneuse (Littoral/Sud-Ouest/Ouest) :} \cle{Duala}, \cle{Bakweri}, \cle{Balong}, \cle{Bassa} (Littoral), \cle{Bamileke} (Ouest - voir Semi-Bantou), \cle{Bamoun} (Ouest - voir Semi-Bantou). Royaumes côtiers (Duala) ou chefferies de haut plateau.
        \end{itemize}
    \item \textbf{Semi-Bantou (ou Bantoïdes de l'Ouest / Grassfields) :} \textbf{Hauts Plateaux de l'Ouest (Bamiléké) et du Nord-Ouest (Bamoun, Tikar, Widikum, Ngemba, etc.)}. \textbf{Particularité :} Langues à classes nominales réduites, organisation en \textbf{chefferies complexes (fonctions, sociétés secrètes, arts de la cour)}. Très forte densité démographique, agriculture intensive (terrasses, fumure), dynamisme entrepreneurial et diaspora nationale/internationale.
    \item \textbf{Pygmées (Premiers habitants de la forêt) :} \cle{Baka} (Est, Sud-Est), \cle{Bakola/Bagyéli} (Sud, Sud-Ouest, littoral). \textbf{Mode de vie :} Chasse-cueillette, pêche, échange avec les villageois bantous (produits forestiers vs fer, sel, manioc). Habitat semi-nomade (campements \emph{mongulu}). \textbf{Menaces :} Déforestation, aires protégées (Dja, Lobéké, Boumba-Bek), sédentarisation forcée, marginalisation sociale. Droits reconnus par la Constitution (2008) et la Déclaration ONU (2007), application encore lacunaire.
\end{itemize}
\end{aretenir}

\begin{savaistu}[Religions et syncrétisme]
Le paysage religieux reflète l'histoire : \textbf{Islam} (majoritaire au Nord, Adamaoua, chez les Peuls ; $\approx$ 25-30 \% pop.), \textbf{Christianisme} (catholique/protestant, majoritaire au Sud, Centre, Ouest, Littoral, Nord-Ouest, Sud-Ouest ; $\approx$ 40-45 \% pop.), \textbf{Religions traditionnelles} (ancêtres, forces de la nature, rites agraires ; $\approx$ 20-25 \% pop., souvent en syncrétisme). La \textbf{laïcité} de l'État (Loi 1905, Constitution 1996) garantit la liberté de culte. Le \cle{vivre ensemble} interreligieux est une réalité quotidienne (fêtes partagées, mariages mixtes) mais des tensions existent (radicalisation Boko Haram au Nord, instrumentalisation politique).
\end{savaistu}

\courssub{L'inégale répartition de la population : les différences de densités}
\begin{definition}[Densité de population]
Rapport entre la population totale d'un territoire et sa superficie. Formule : $D = \frac{P}{S}$ (exprimée en \textbf{hab/km²}). C'est un \textbf{moyen grossier} (masque les vides et les concentrations locales) mais un \textbf{indicateur de synthèse} indispensable pour la comparaison régionale.
\end{definition}

\begin{aretenir}[Facteurs de la répartition inégale au Cameroun]
\begin{itemize}
    \item \textbf{Physiques :} \textbf{Altitude/Sols :} Hauts Plateaux de l'Ouest (volcaniques, fertiles, climat tempéré) $\to$ fortes densités. Plaine du Diamaré/Logone (alluvions) $\to$ fortes densités. Plateau de l'Adamaoua (sols lessivés, ferralitiques, altitude 1000-1500 m) $\to$ faibles densités. Forêt dense (Est/Sud) : sols acides, lessivés, maladies (paludisme, onchocercose, trypanosomiase) $\to$ faibles densités. \textbf{Climat :} Soudanais (saison sèche marquée, cultures pluviales) plus favorable que l'équatorial (pluies constantes, contrainte forestière).
    \item \textbf{Historiques :} \textbf{Traite négrière} (XVII\textsuperscript{e}-XIX\textsuperscript{e} s.) : vidage démographique de l'Adamaoua, de l'Est et du Sud (razzias). \textbf{Colonisation :} Axe de pénétration \textbf{Douala-Yaoundé-Ngaoundéré} (chemin de fer Transcamerounais, RN1) $\to$ concentration le long de l'axe. Plantations agro-industrielles (Sud-Ouest/Littoral : banane, palme, hévéa, cacao) $\to$ attraction de main-d'œuvre (Bamiléké, Nordistes, Nigérians).
    \item \textbf{Économiques actuels :} \textbf{Métropolisation :} \cle{Douala} (capitale économique, port, industrie, $\approx$ 4 M hab. aire urbaine) et \cle{Yaoundé} (capitale politique, administration, services, $\approx$ 3,5 M hab.) = pôles d'aspiration massive. \textbf{Agro-industrie :} Bananeraies (Njombé, Mungo), Palmeraies (Socapalm : Dibombari, Eseka, Kienké), Hévéa (Hévécam), Coton (Sodecoton : Nord/Adamaoua). \textbf{Mines (projets) :} Fer de \cle{Mbalam} (Est), Bauxite de \cle{Minim-Martap} (Adamaoua), Cobalt/Nickel de \cle{Nkamouna} (Est) $\to$ futurs pôles de croissance.
\end{itemize}
\end{aretenir}

\courssub{Mobilité, brassage des populations et intégration nationale}
\begin{propriete}[Types de mobilités spatiales au Cameroun]
\begin{itemize}
    \item \textbf{Migrations internes (les plus intenses) :}
        \begin{itemize}
            \item \textit{Rural $\to$ Urbain :} Exode massif vers Douala, Yaoundé, Bafoussam, Bamenda, Maroua, Garoua, Ngaoundéré, Bertoua, Ebolowa, Kumba, Limbe. Taux d'urbanisation $\approx$ 58 \% (2023).
            \item \textit{Rural $\to$ Rural :} \textbf{Est-Ouest} (Bamiléké $\to$ plantations SW/Littoral, forêts Est/Sud pour cacao/bois) ; \textbf{Nord-Sud} (Peuls/Mbororo $\to$ Adamaoua, Est, Centre pour pâturages) ; \textbf{Sud-Nord} (Beti, Bassa $\to$ Nord/Adamaoua pour commerce, administration).
            \item \textit{Urbain $\to$ Urbain :} Circulation entre métropoles et villes secondaires.
        \end{itemize}
    \item \textbf{Migrations internationales :} \textbf{Immigration :} Nigérians (commerce, pêche, éducation), Centrafricains, Tchadiens, réfugiés (Centrafrique au Est/Adamaoua, Nigeria au Extrême-Nord). \textbf{Émigration :} Diaspora qualifiée (France, USA, Allemagne, Canada, pays du Golfe) $\to$ transferts de fonds ($\approx$ 300-400 Mds FCFA/an, > aide publique au développement).
    \item \textbf{Mobilité circulaire / Transhumance :} \cle{Mbororo} (Nord $\leftrightarrow$ Adamaoua/Centrafrique/Est), \cle{Arabes Choa} (Logone $\leftrightarrow$ Yaérés). Conflits récurrents avec agriculteurs sédentaires (extension cultures coton/maïs).
\end{itemize}
\end{propriete}

\begin{methode}[Analyser le brassage comme facteur d'intégration (ou de tension)]
\begin{enumerate}
    \item Identifier les \textbf{flux principaux} (origines $\to$ destinations, volumes, périodes).
    \item Analyser les \textbf{interactions} : marchés intercommunautaires (ex. : Maroua, Ngaoundéré, Douala, Yaoundé), mariages mixtes, apprentissage des langues (fulfulde, ewondo, bamiléké, pidgin, français, anglais), associations de développement (GIC, tontines).
    \item Évaluer les \textbf{tensions} : conflits fonciers (droit coutumier vs droit moderne, autochtones vs allogènes), conflits agro-pastoraux, replis identitaires, crise anglophone (Nord-Ouest/Sud-Ouest).
    \item Conclure sur le \textbf{bilan} : le brassage est un \textbf{fait majeur} (peu de zones ethniquement pures), source de richesse culturelle et de main-d'œuvre mobile, mais nécessite une \textbf{gouvernance foncière et politique équitable} (décentralisation, code foncier, commission nationale bilinguisme/multiculturalisme) pour devenir un ciment de l'unité nationale.
\end{enumerate}
\end{methode}

\courssub{TP 2 : Construction et lecture de la carte de répartition de la population}

\courssub{Objectifs de l'activité}
\begin{itemize}
    \item \textbf{Calculer} les densités de population des dix régions du Cameroun à partir de données démographiques récentes (projections 2023).
    \item \textbf{Construire} une carte choroplèthe à trois classes de densité (faible, moyenne, forte) en respectant la sémiologie graphique (choix des seuils, nuancier, légende ordonnée).
    \item \textbf{Analyser} les contrastes spatiaux mis en évidence : localisation des zones denses (hauts plateaux, plaines du Nord, métropoles), zones moyennes (piémonts, façade atlantique) et zones creuses (forêt de l'Est, plateau de l'Adamaoua, massif du Sud).
    \item \textbf{Expliquer} ces inégalités par la combinaison de facteurs physiques (relief, climat, sols), historiques (peuplement ancien, colonisation, axes de pénétration) et économiques (urbanisation, activités minières/agro-industrielles).
    \item \textbf{Proposer} deux mesures concrètes d'aménagement du territoire favorisant un meilleur équilibre régional et le renforcement de l'intégration nationale (ex. : désenclavement de l'Est, pôles de développement secondaire, gestion concertée des ressources pastorales).
\end{itemize}

\courssub{Situation}
Le Ministère de l'Économie, de la Planification et de l'Aménagement du Territoire (MINEPAT) lance, dans le cadre de la \cle{Stratégie Nationale de Développement 2020-2030 (SND30)}, un concours national à destination des lycées : \og Ma carte pour un Cameroun équilibré \fg. L'objectif est de sensibiliser la jeunesse aux défis de l'aménagement du territoire pour consolider l'\cle{intégration nationale}.

Votre équipe reçoit un \textbf{dossier documentaire} composé de :
\begin{enumerate}
    \item Un \textbf{tableau statistique} (source : BUCREP, projections 2023) donnant la population totale et la superficie de chacune des dix régions.
    \item Une \textbf{carte de base} (fond de carte au 1/5 000 000) représentant les limites régionales, les chefs-lieux, le réseau hydrographique principal (Sanaga, Bénoué, Logone) et les grands axes routiers (RN1, RN3, RN6, RN10).
    \item Deux \textbf{photographies légendées} : \textit{Photo A} : « Monts Mandara, département du Mayo-Sava (Extrême-Nord) : habitat dense, terrasses culturales, forte pression foncière » ; \textit{Photo B} : « Forêt dense humide du Haut-Nyong (Est) : village de chasseurs-cueilleurs Baka, piste cyclable unique, densité \textless 5 hab/km² ».
    \item Deux \textbf{extraits de presse} : \textit{Extrait 1} (Le Jour, mars 2024) : « Conflits agriculteurs-éleveurs dans le Mayo-Rey : la transhumance bloquée par l'extension des cultures de coton » ; \textit{Extrait 2} (Cameroon Tribune, janvier 2024) : « Projet autoroutier Yaoundé-Ngaoundéré : un levier pour l'intégration de l'Adamaoua et de l'Est ».
\end{enumerate}

\begin{center}
\begin{tabularx}{\linewidth}{|l|X|X|X|}\hline
\rowcolor{popblueL}
\textbf{Région} & \textbf{Superficie (km²)} & \textbf{Population 2023 (hab)} & \textbf{Densité (hab/km²)} \\ \hline
Adamaoua & 63 691 & 1 180 000 & \textit{À calculer} \\ \hline
Centre (y.c. Yaoundé) & 68 926 & 4 520 000 & \textit{À calculer} \\ \hline
Est & 109 002 & 920 000 & \textit{À calculer} \\ \hline
Extrême-Nord & 34 263 & 4 480 000 & \textit{À calculer} \\ \hline
Littoral (y.c. Douala) & 20 248 & 3 810 000 & \textit{À calculer} \\ \hline
Nord & 66 090 & 2 780 000 & \textit{À calculer} \\ \hline
Nord-Ouest & 17 300 & 2 200 000 & \textit{À calculer} \\ \hline
Ouest & 13 892 & 2 790 000 & \textit{À calculer} \\ \hline
Sud & 47 110 & 800 000 & \textit{À calculer} \\ \hline
Sud-Ouest & 25 380 & 2 010 000 & \textit{À calculer} \\ \hline
\textbf{TOTAL} & \textbf{475 602} & \textbf{25 490 000} & \textbf{~53,6} \\ \hline
\end{tabularx}
\end{center}
\legende{Tableau 1 : Données démographiques par région (Projections BUCREP 2023, arrondis). La densité nationale moyenne est de 53,6 hab/km².}

\courssub{Consignes}
\begin{enumerate}
    \item \textbf{Compléter} le tableau 1 en calculant la densité de population de chaque région (arrondi à l'entier). Identifier la densité nationale moyenne.
    \item \textbf{Définir} trois classes de densité (faible, moyenne, forte) en vous appuyant sur la moyenne nationale (53,6 hab/km²) et la distribution des valeurs régionales. Justifier le choix de vos seuils.
    \item \textbf{Réaliser} la carte choroplèthe sur le fond de carte fourni : colorier chaque région selon sa classe, tracer la légende complète (titre, seuils, surfaces colorées, source, échelle, orientation).
    \item \textbf{Rédiger} un commentaire structuré (15-20 lignes) comprenant :
    \begin{itemize}
        \item La localisation précise des trois types d'espaces (régions concernées, position relative).
        \item L'explication des contrastes (facteurs physiques, historiques, économiques).
        \item Deux propositions d'aménagement concrètes pour réduire les déséquilibres et favoriser le \cle{vivre ensemble}.
    \end{itemize}
    \item \textbf{Présenter} oralement votre carte et vos propositions en 3 minutes (restitution collective).
\end{enumerate}

\courssub{Ressources à mobiliser}
\begin{definition}[Densité de population]
Rapport entre la population totale d'un territoire et sa superficie. Formule : $D = \frac{P}{S}$ (hab/km²).
\end{definition}

\begin{methode}[Construction d'une carte choroplèthe (méthode des seuils statistiques)]
\begin{enumerate}
    \item Calculer l'indicateur (densité) pour chaque unité spatiale (région).
    \item Ordonner les valeurs et observer la distribution (écart-type, médiane, moyenne).
    \item Choisir un nombre de classes réduit (3 à 5) et des seuils \textbf{signifiants} (moyenne nationale, ruptures naturelles). Éviter les classes à effectifs nuls.
    \item Appliquer un \textbf{nuancier ordonné} : teintes claires pour les valeurs faibles, teintes sombres/saturées pour les valeurs fortes (ex. : jaune $\to$ orange $\to$ rouge ; ou vert clair $\to$ jaune $\to$ brun).
    \item Légender rigoureusement : titre informatif, classes avec bornes incluses/exclues [min - max[ , source, date, auteur.
\end{enumerate}
\end{methode}

\begin{aretenir}[Facteurs de la répartition inégale au Cameroun]
\begin{itemize}
    \item \textbf{Physiques :} Altitude (monts Cameroun, Mandara, Adamaoua $\to$ sols volcaniques fertiles vs sols lessivés), Climat (soudanais favorable aux cultures pluviales vs équatorial contraignant), Hydrologie (vallées alluviales).
    \item \textbf{Historiques :} Ancienneté du peuplement (Paléo-soudanais aux Mandara, Bantou à l'Ouest), Traite négrière (vidage de l'Est/Adamaoua), Colonisation (axe Douala-Yaoundé-Ngaoundéré, plantations SW/Littoral).
    \item \textbf{Économiques :} Métropolisation (Douala, Yaoundé $\to$ effet d'aspiration), Agro-industrie (bananeraies Njombé, palmeraies Sanaga, coton Nord), Mines (bauxite Minim-Martap, fer Mbalam $\to$ pôles futurs).
\end{itemize}
\end{aretenir}

\courssub{Démarche et corrigé commenté}
\begin{exoresolu}[Activité d'intégration : Carte de la répartition de la population]
\textbf{1. Complétion du tableau (Calculs)}

On applique $D = P / S$ pour chaque région (arrondi à l'unité) :
\begin{align*}
\text{Adamaoua} &: 1\,180\,000 / 63\,691 \approx \textbf{19 hab/km²} \\
\text{Centre} &: 4\,520\,000 / 68\,926 \approx \textbf{66 hab/km²} \\
\text{Est} &: 920\,000 / 109\,002 \approx \textbf{8 hab/km²} \\
\text{Extrême-Nord} &: 4\,480\,000 / 34\,263 \approx \textbf{131 hab/km²} \\
\text{Littoral} &: 3\,810\,000 / 20\,248 \approx \textbf{188 hab/km²} \\
\text{Nord} &: 2\,780\,000 / 66\,090 \approx \textbf{42 hab/km²} \\
\text{Nord-Ouest} &: 2\,200\,000 / 17\,300 \approx \textbf{127 hab/km²} \\
\text{Ouest} &: 2\,790\,000 / 13\,892 \approx \textbf{201 hab/km²} \\
\text{Sud} &: 800\,000 / 47\,110 \approx \textbf{17 hab/km²} \\
\text{Sud-Ouest} &: 2\,010\,000 / 25\,380 \approx \textbf{79 hab/km²}
\end{align*}
\textit{Vérification :} Total Pop / Total Surf = 25,49 M / 475 602 $\approx$ 53,6 hab/km² (cohérence).

\textbf{2. Choix des seuils et classes (3 classes)}

Valeurs ordonnées : 8 (Est) ; 17 (Sud) ; 19 (Adamaoua) ; 42 (Nord) ; 66 (Centre) ; 79 (Sud-Ouest) ; 127 (Nord-Ouest) ; 131 (Extrême-Nord) ; 188 (Littoral) ; 201 (Ouest).
Médiane $\approx$ (66+79)/2 = 72,5. Moyenne nationale = 53,6.
On observe une rupture nette entre le Nord (42) et le Centre (66), et une autre entre le Sud-Ouest (79) et le Nord-Ouest (127).
\emph{Choix pédagogique pertinent :} Seuil bas $\approx$ moyenne nationale (54) arrondi à \textbf{50} pour séparer « Faible » et « Moyenne ». Seuil haut $\approx$ rupture forte avant les très fortes densités (127) $\to$ \textbf{100} pour séparer « Moyenne » et « Forte ».
\begin{itemize}
    \item \textbf{Classe 1 (Faible) :} $[0 - 50[$ hab/km² $\to$ Est, Sud, Adamaoua, Nord. (4 régions)
    \item \textbf{Classe 2 (Moyenne) :} $[50 - 100[$ hab/km² $\to$ Centre, Sud-Ouest. (2 régions)
    \item \textbf{Classe 3 (Forte) :} $[100 - +\infty[$ hab/km² $\to$ Extrême-Nord, Nord-Ouest, Ouest, Littoral. (4 régions)
\end{itemize}
\emph{Justification :} Ce découpage isole clairement le « Cameroun vide » (Est/Sud/Adamaoua + Nord sahélien peu dense), le « Cameroun intermédiaire » (axe central/atlantique) et le « Cameroun plein » (hauts plateaux de l'Ouest, plaine du Logone, métropole doubalaise).

\textbf{3. Carte choroplèthe (Réalisation graphique)}

\begin{popfigure}
\begin{center}
\includegraphics[width=\linewidth,height=10.5cm,keepaspectratio]{N02/cmr_dens2000.jpg}
\end{center}
\legende{Carte de la densité de population du Cameroun en 2000 (Gridded Population of the World, CIESIN, Université Columbia ; légende en anglais, en personnes par km\textsuperscript{2}). On y lit : les plus fortes densités (brun foncé, 250 à plus de 1~000 hab/km\textsuperscript{2}) en taches isolées autour de Douala sur la côte, de Yaoundé et des hautes terres de l'Ouest (Bamenda, Bafoussam) ; une densité moyenne (orange, 25 à 249 hab/km\textsuperscript{2}) sur le Littoral, l'Ouest, le Nord-Ouest, le Centre et l'Extrême-Nord ; de faibles densités (5 à 24 hab/km\textsuperscript{2}) sur le plateau de l'Adamaoua, le Nord et le Sud ; et le \og vide \fg{} de l'Est forestier (1 à 4 hab/km\textsuperscript{2}). Elle sert de modèle de lecture pour la carte choroplèthe à construire (classes de densité faible, moyenne, forte). \\ Source : Wikimedia Commons, SEDAC (CIESIN, Columbia University), CC BY 2.0.}
\end{popfigure}

\textbf{4. Commentaire type (Modèle de rédaction)}

\emph{« La carte choroplèthe révèle un Cameroun à trois vitesses démographiques.}

\emph{\textbf{1. Les zones faiblement peuplées (< 50 hab/km²) }occupent plus de 60 \% du territoire national : l'Est (8 hab/km²), le Sud (17 hab/km²), l'Adamaoua (19 hab/km²) et le Nord (42 hab/km²). Ces espaces correspondent à la forêt dense humide (contraintes climatiques, sols lessivés, maladies vectorielles), au plateau de l'Adamaoua (altitude, climat frais, sols pauvres) et à la zone sahélienne (aridité, insécurité Boko Haram). Ce sont des « vides » relatifs, traversés par des axes de pénétration récents (RN10, RN15).}

\emph{\textbf{2. Les zones moyennement peuplées (50-100 hab/km²) }concernent la région du Centre (66 hab/km², hors Yaoundé) et le Sud-Ouest (79 hab/km²). Le Centre bénéficie de l'effet d'entraînement de la capitale politique Yaoundé (3,5 M d'hab. aire urbaine) et de l'axe routier/ferroviaire transversal. Le Sud-Ouest, malgré la crise anglophone (depuis 2016), conserve une densité soutenue par l'agro-industrie (CDC, Pamol) et l'urbanisation côtière (Limbe, Buea).}

\emph{\textbf{3. Les zones fortement peuplées (> 100 hab/km²) }se concentrent sur 15 \% du territoire : l'Ouest (201 hab/km²), le Nord-Ouest (127 hab/km²), le Littoral (188 hab/km²) et l'Extrême-Nord (131 hab/km²). L'Ouest et le Nord-Ouest (Hauts Plateaux Bamileke/Bamoun) allient sols volcaniques fertiles, climat tempéré, ancienneté du peuplement (organisation chefferies) et forte dynamique entrepreneuriale. Le Littoral est porté par Douala (capitale économique, port, 4 M hab. aire urbaine). L'Extrême-Nord (plaine du Diamaré, monts Mandara) combine fertilité alluviale, paléo-peuplement dense (Paléo-soudanais) et pression foncière extrême (Photo A).}

\emph{\textbf{Explication des contrastes :} }L'opposition Est/Ouest structure le pays : à l'Est, le « vide » forestier hérité de la traite et de l'enclavement ; à l'Ouest, le « plein » des hautes terres et des métropoles littorales. Le Nord présente une dichotomie : plaine dense (Extrême-Nord) vs plateau creux (Nord/Adamaoua). L'urbanisation (Yaoundé, Douala) et les axes de transport (Transcamerounais) agissent comme des aimants démographiques.

\emph{\textbf{Propositions pour l'intégration nationale :}}
\emph{(1) \textbf{Achèvement du corridor routier/ferroviaire Yaoundé-Ngaoundéré-Bertoua (RN10/RN15)} pour désenclaver l'Est et l'Adamaoua, valoriser les ressources minières (fer Mbalam, bauxite Minim-Martap) et créer des pôles urbains secondaires (Bertoua, Batouri, Meiganga) capables d'absorber la croissance démographique.}
\emph{(2) \textbf{Mise en place d'une « Agence des Hauts Plateaux et du Sahel »} pour gérer concurremment les conflits fonciers agriculteurs-éleveurs (Extrême-Nord, Nord, Adamaoua) via la délimitation de couloirs de transhumance sécurisés, la promotion de l'élevage sédentaire amélioré et la restauration des sols (lutte contre l'érosion dans l'Ouest), garantissant la paix sociale et le \cle{vivre ensemble}. »}

\textbf{5. Restitution orale}
L'équipe présente la carte (vidéo-projecteur ou affiche), souligne le choix des seuils (moyenne nationale + ruptures), lit le commentaire et défend les deux propositions face aux questions de la classe (ex. : coût du ferroviaire, faisabilité couloirs transhumance).
\tcblower
\textbf{Corrigé commenté par le professeur}
\begin{itemize}
    \item \textbf{Calculs :} Vérifier l'arrondi (ex. Adamaoua 18,5 $\to$ 19). La densité du Nord (42) la place en classe 1, ce qui est discutable (proche du seuil 50) mais justifié par la rupture naturelle avec le Centre (66). Accepter une classe « Nord » à part si argumenté (transition).
    \item \textbf{Seuils :} L'utilisation de la moyenne nationale (53,6) comme seuil bas est une excellente référence. Le seuil haut à 100 isole bien les 4 régions « moteurs ». Une alternative : seuils quartiles (33 / 100) mais moins signifiants géographiquement.
    \item \textbf{Carte :} Vérifier : Légende \textbf{obligatoire} (Titre, 3 cases colorées avec bornes exactes, Source, Auteur, Date, Échelle, Nord). Respect du nuancier ordonné (clair $\to$ foncé). Pas de hachures, aplats de couleur pleins. Noms des 10 chefs-lieux positionnés.
    \item \textbf{Commentaire :} Exiger la structure en 3 paragraphes (Localisation, Explication, Propositions). Sanctionner l'absence de chiffres (densités, \% territoire) ou de noms de régions/villes. Les propositions doivent être \textbf{spatiales} (localisées) et \textbf{opérationnelles} (acteurs, leviers), pas des vœux pieux (« il faut développer l'Est » $\to$ insuffisant).
    \item \textbf{Intégration :} Lier explicitement « désenclavement » $\to$ « brassage populations » $\to$ « intégration nationale ». Citer les groupes humains concernés (Baka à l'Est, Mbororo à l'Adamaoua, Paléo-soudanais au Nord, Bantou à l'Ouest).
\end{itemize}
\end{exoresolu}

\courssub{Bilan de l'activité}
Cette activité d'intégration a permis de mettre en évidence que :
\begin{itemize}
    \item La \textbf{répartition de la population camerounaise est extrêmement contrastée} (coefficient de variation inter-régional > 100 \%), structurée par une diagonale « vide » (Est-Sud-Adamaoua-Nord) et un croissant « plein » (Extrême-Nord $\to$ Ouest $\to$ Littoral $\to$ Sud-Ouest).
    \item Ces contrastes ne sont pas figés : ils résultent d'une \textbf{histoire longue} (peuplements anciens, traite, colonisation) et de \textbf{dynamiques actuelles} (métropolisation, conflits, projets d'infrastructures).
    \item La \textbf{carte choroplèthe} est un outil puissant de diagnostic territorial, à condition de maîtriser le choix des seuils (moyenne, ruptures naturelles) et la sémiologie (nuancier ordonné).
    \item L'\textbf{intégration nationale} ne se décrète pas : elle se construit par l'aménagement (désenclavement de l'Est, corridor Nord-Sud) et la \textbf{gouvernance foncière/pastorale} (prévention des conflits agriculteurs-éleveurs), conditions sine qua non du \cle{vivre ensemble} entre les 250 groupes ethniques (Soudanais, Bantou, Semi-Bantou, Pygmées) qui composent la nation.
\end{itemize}

\begin{attention}[Pièges à éviter pour le Bac]
\begin{itemize}
    \item Confondre \textbf{densité absolue} (hab/km²) et \textbf{densité relative} (hab/ha cultivable) — le programme exige la densité absolue.
    \item Oublier d'indiquer \textbf{l'échelle} et \textbf{l'orientation} sur la carte (perte de 1 à 2 points).
    \item Utiliser un \textbf{nuancier non ordonné} (ex. : rouge pour faible, vert pour fort) ou des hachures au lieu d'aplats.
    \item Rédiger des propositions \textbf{hors-sol} (ex. « construire des hôpitaux partout ») sans les localiser ni les lier aux flux de population.
    \item Négliger le \textbf{Nord} dans l'analyse : ce n'est pas un bloc homogène (Extrême-Nord dense vs Nord/Adamaoua creux).
\end{itemize}
\end{attention}

\newpage

\coursec{Méthodes et exercices résolus}

\coursec{Les grands groupes humains}

\courssub{Observer, localiser et décrire les groupes humains du Cameroun}

\begin{methode}[Décrire et localiser un groupe humain à partir d'une carte]
Pour traiter une question du type « Localisez et décrivez les grands groupes humains du Cameroun », suis toujours la même démarche :
\begin{enumerate}
\item \textbf{Nommer les ensembles} : identifier les trois grands groupes (soudanais, bantou et semi-bantou, pygmées) et les sous-groupes (paléo-soudanais, néo-soudanais).
\item \textbf{Localiser par rapport à des repères fixes} : utiliser les régions administratives (Extrême-Nord, Nord, Adamaoua, Centre, Sud, Est, Littoral, Ouest, Nord-Ouest, Sud-Ouest), les villes (Maroua, Garoua, Ngaoundéré, Yaoundé, Ebolowa, Bertoua, Douala, Bafoussam, Bamenda, Buéa) et les lignes de partage (le fleuve Sanaga, le plateau de l'Adamaoua).
\item \textbf{Décrire la répartition} : du nord vers le sud, en précisant qui occupe quelle zone (soudanais au nord de l'Adamaoua, bantou et semi-bantou au sud, pygmées dans la forêt du Sud-Est).
\item \textbf{Caractériser chaque groupe} : une ou deux traits socio-culturels (habitat, activité, organisation sociale, religion dominante).
\item \textbf{Conclure} par une phrase de synthèse reliant répartition et intégration nationale.
\end{enumerate}
\textbf{Réflexe} : jamais de liste sans localisation ; jamais de localisation sans repère (région, ville ou fleuve).
\end{methode}

\begin{exoresolu}[Localisation des grands groupes humains]
\textbf{Énoncé.} À l'aide de vos connaissances, localisez et décrivez les grands groupes humains du Cameroun. (10 points)
\tcblower
\textbf{Solution détaillée.}
\begin{enumerate}
\item \textbf{Nomenclature.} La population camerounaise se rattache à trois grands ensembles humains : les \cle{soudanais} (subdivisés en paléo-soudanais et néo-soudanais), les \cle{bantou} et \cle{semi-bantou}, et les \cle{pygmées}.
\item \textbf{Localisation.}
\begin{itemize}
\item Les \textbf{soudanais} occupent le nord du pays, au nord du plateau de l'Adamaoua : régions de l'Extrême-Nord (Maroua), du Nord (Garoua) et de l'Adamaoua (Ngaoundéré).
\item Les \textbf{bantou} occupent le sud forestier : régions du Centre (Yaoundé), du Sud (Ebolowa), de l'Est (Bertoua), du Littoral (Douala) et du Sud-Ouest (Buéa).
\item Les \textbf{semi-bantou} (ou peuples des Grassfields) occupent les hautes terres de l'Ouest (Bafoussam, Dschang) et du Nord-Ouest (Bamenda).
\item Les \textbf{pygmées} (Baka, Bagyeli, Medzan) vivent dans la forêt dense du Sud-Est, autour de Yokadouma, Lomié, Moloundou et Campo.
\end{itemize}
\item \textbf{Caractérisation.}
\begin{itemize}
\item \textbf{Soudanais :} peuples de savane, agriculteurs et éleveurs. Les néo-soudanais (Peuls, Arabes Choa, Kanouri) sont largement islamisés et organisés en lamibés. Les paléo-soudanais (Kirdi : Mafa, Mofu, Kapsiki, Matakam des monts Mandara) ont conservé des religions traditionnelles et des chefferies de montagne.
\item \textbf{Bantou et semi-bantou :} peuples de la forêt et des hautes terres, pratiquant l'agriculture itinérante sur brûlis (forêt) ou l'agriculture intensive sur sols volcaniques (Ouest), de religion chrétienne ou animiste, organisés en chefferies et sociétés secrètes.
\item \textbf{Pygmées :} chasseurs-cueilleurs semi-nomades de la forêt équatoriale, organisation égalitaire en campements.
\end{itemize}
\item \textbf{Conclusion.} Cette répartition, héritée des grandes migrations (bantoues, soudanaises), fait du Cameroun une mosaïque de plus de 250 groupes ethniques, dont le brassage constitue un enjeu majeur de l'intégration nationale.
\end{enumerate}
\end{exoresolu}

\courssub{Calculer et interpréter une densité de population}

\begin{methode}[Calculer une densité et classer les espaces]
\begin{enumerate}
\item \textbf{Formule} : densité $D = \dfrac{\text{population (habitants)}}{\text{superficie (km}^2\text{)}}$, exprimée en hab/km$^2$.
\item \textbf{Vérifier les unités} avant le calcul (population en habitants, superficie en km$^2$).
\item \textbf{Comparer} les résultats entre eux et à la densité moyenne nationale (environ 55 à 60 hab/km$^2$ en 2023, ordre de grandeur à actualiser selon les dernières projections de l'INS).
\item \textbf{Classer} les espaces en trois catégories : densément peuplés (fortes densités, supérieures à 100 hab/km$^2$ localement), moyennement peuplés (autour de la moyenne nationale, 10 à 100 hab/km$^2$), faiblement peuplés (moins de 10 hab/km$^2$).
\item \textbf{Interpréter} : relier la densité aux facteurs explicatifs (relief, climat, sols, histoire, activités économiques, urbanisation).
\end{enumerate}
\textbf{Réflexe} : une densité seule ne dit rien ; c'est la \emph{comparaison} qui crée l'information géographique.
\end{methode}

\begin{exoresolu}[Calcul et classement des densités régionales]
\textbf{Énoncé.} On donne (données simplifiées, ordres de grandeur) : Littoral : 3,3 millions d'habitants pour \SI{20239}{\kilo\meter\squared} ; Ouest : 1,9 million d'habitants pour \SI{13872}{\kilo\meter\squared} ; Est : 0,8 million d'habitants pour \SI{109002}{\kilo\meter\squared}. 1) Calculez la densité de chaque région. 2) Classez-les et proposez une explication. (10 points)
\tcblower
\textbf{Solution détaillée.}
\begin{enumerate}
\item \textbf{Calculs.}
\begin{align*}
D_{\text{Littoral}} &= \frac{3\,300\,000}{20\,239} \approx 163 \text{ hab/km}^2 \\
D_{\text{Ouest}} &= \frac{1\,900\,000}{13\,872} \approx 137 \text{ hab/km}^2 \\
D_{\text{Est}} &= \frac{800\,000}{109\,002} \approx 7 \text{ hab/km}^2
\end{align*}
\item \textbf{Classement.} Le Littoral et l'Ouest sont \textbf{densément peuplés} (densités très supérieures à la moyenne nationale d'environ 55--60 hab/km$^2$) ; l'Est est \textbf{faiblement peuplé} (densité très inférieure à 10 hab/km$^2$).
\item \textbf{Explication.} Le Littoral concentre les activités (port de Douala, industries, emplois tertiaires) qui attirent les migrants. L'Ouest doit ses fortes densités à la fertilité des sols volcaniques des hautes terres et à une ancienneté du peuplement. L'Est, couvert de forêt dense équatoriale difficile à défricher, enclavé et pauvre en infrastructures, reste un vide humain.
\item \textbf{Conclusion.} La population camerounaise est donc inégalement répartie : quelques régions concentrent l'essentiel des hommes pendant que de vastes espaces restent quasi vides.
\end{enumerate}
\end{exoresolu}

\courssub{Calculer une évolution démographique et une part en pourcentage}

\begin{methode}[Taux de variation et part relative]
\begin{enumerate}
\item \textbf{Taux de variation} entre deux dates : $t = \dfrac{V_f - V_i}{V_i} \times 100$ (en \%). Un $t$ positif traduit une augmentation.
\item \textbf{Part d'un élément dans un ensemble} : $p = \dfrac{\text{effectif de l'élément}}{\text{effectif total}} \times 100$ (en \%).
\item \textbf{Toujours} écrire la formule littérale, puis l'application numérique, puis le résultat arrondi avec son unité (\%).
\item \textbf{Interpréter} le résultat en une phrase : « la population a été multipliée par... », « telle région rassemble près de... \% des Camerounais ».
\end{enumerate}
\end{methode}

\begin{exoresolu}[Croissance démographique et concentration urbaine]
\textbf{Énoncé.} La population du Cameroun est passée d'environ 10,5 millions d'habitants en 1987 à environ 28 millions en 2023 (ordres de grandeur). Douala et Yaoundé comptent ensemble environ 8 millions d'habitants en 2023. 1) Calculez le taux d'accroissement de la population entre 1987 et 2023. 2) Calculez la part des deux métropoles dans la population nationale en 2023. 3) Commentez. (10 points)
\tcblower
\textbf{Solution détaillée.}
\begin{enumerate}
\item \textbf{Taux d'accroissement.}
\[ t = \frac{28 - 10{,}5}{10{,}5} \times 100 = \frac{17{,}5}{10{,}5} \times 100 \approx 166{,}7\,\% \]
La population a donc augmenté d'environ 167 \% en 36 ans : elle a été multipliée par environ $1 + 1{,}667 \approx 2{,}7$.
\item \textbf{Part des deux métropoles.}
\[ p = \frac{8}{28} \times 100 \approx 28{,}6\,\% \]
Douala et Yaoundé rassemblent donc près de 29 \% des Camerounais, soit presque 3 habitants sur 10.
\item \textbf{Commentaire.} La très forte croissance démographique s'accompagne d'une concentration urbaine spectaculaire : les migrations intérieures, alimentées par l'exode rural, gonflent les deux métropoles, pôles d'emplois et de services. Ce déséquilibre accentue l'inégale répartition de la population et pose des défis (logement, chômage, insécurité foncière, gestion des déchets).
\item \textbf{Conclusion.} Croissance rapide et concentration urbaine sont les deux traits majeurs de la dynamique démographique camerounaise.
\end{enumerate}
\end{exoresolu}

\courssub{Construire et légender une carte de répartition de la population (TP 2)}

\begin{methode}[Construire une carte de densités en trois classes]
\begin{enumerate}
\item \textbf{Tracer le fond de carte} : contour simplifié du Cameroun, limites des 10 régions, flèche du nord, titre précis, échelle indicative.
\item \textbf{Choisir trois classes de densité} : fortement peuplé ($>$ 100 hab/km$^2$), moyennement peuplé (10 à 100 hab/km$^2$), faiblement peuplé ($<$ 10 hab/km$^2$).
\item \textbf{Associer une gamme de couleurs} : une seule couleur dégradée (par exemple orange clair $\rightarrow$ orange $\rightarrow$ orange foncé) ; la teinte la plus forte pour la densité la plus forte.
\item \textbf{Reporter les données} région par région à partir du tableau statistique fourni.
\item \textbf{Rédiger la légende} : chaque couleur correspond à une classe chiffrée ; ajouter les villes principales (points nommés) et les axes de migration (flèches) si demandé.
\item \textbf{Lire la carte} : dégager les contrastes (Nord-Ouest/Ouest/Littoral/Extrême-Nord peuplés ; Est et Sud-Est vides) et les expliquer.
\end{enumerate}
\textbf{Réflexe} : une carte sans titre, sans nord ou sans légende chiffrée est une carte non achevée : elle perd des points.
\end{methode}

\begin{exoresolu}[TP 2 : croquis de la répartition de la population]
\textbf{Énoncé.} À partir du tableau des densités régionales fourni, construisez un croquis du Cameroun montrant les zones densément, moyennement et faiblement peuplées, puis dégagez les grandes leçons de la répartition. (10 points)
\tcblower
\textbf{Solution détaillée.}
\begin{enumerate}
\item \textbf{Classement préalable (données indicatives).}
\begin{itemize}
\item \textbf{Densément peuplées ($>$ 100 hab/km$^2$) :} Ouest, Nord-Ouest, Littoral, Extrême-Nord.
\item \textbf{Moyennement peuplées (10 à 100 hab/km$^2$) :} Centre, Sud-Ouest, Nord, Adamaoua, Sud.
\item \textbf{Faiblement peuplée ($<$ 10 hab/km$^2$) :} Est.
\end{itemize}
\item \textbf{Construction du croquis} (modèle schématique) :
\end{enumerate}
\begin{popfigure}
\begin{center}
\includegraphics[width=\linewidth,height=10cm,keepaspectratio]{N02/cmr_dens2000.jpg}
\end{center}
\legende{Fond de carte pour le croquis de la répartition de la population au Cameroun : densités en 2000 en six classes (légende en anglais). Le croquis attendu en retient trois ensembles : \textbf{1} : les taches densément peuplées (Ouest, Nord-Ouest, Littoral autour de Douala, Extrême-Nord) ; \textbf{2} : les zones intermédiaires (Centre, Sud, Nord, Adamaoua, Sud-Ouest) ; \textbf{3} : le vide humain de l'Est forestier, en beige très clair. \\ Source : Wikimedia Commons, SEDAC (CIESIN, Columbia University), CC BY 2.0.}
\end{popfigure}
\begin{enumerate}
\setcounter{enumi}{2}
\item \textbf{Lecture.} Trois leçons se dégagent : (a) les hommes se concentrent sur les hautes terres de l'Ouest (sol volcanique, ancienneté), le Littoral industrialisé (Douala) et les plaines de l'Extrême-Nord (bassin cotonnier, Maroua) ; (b) le Centre s'est densifié autour de Yaoundé (fonction administrative) ; (c) l'Est et le Sud-Est forestiers restent presque vides (enclavement, forêt dense, sols fragiles).
\item \textbf{Conclusion.} La répartition de la population camerounaise est fortement contrastée : moins de la moitié du territoire concentre l'essentiel des habitants, ce qui nourrit de puissants courants migratoires internes.
\end{enumerate}
\end{exoresolu}

\courssub{Analyser les mobilités et le brassage des populations}

\begin{methode}[Commenter un document sur les migrations intérieures]
\begin{enumerate}
\item \textbf{Identifier la nature du document} (carte de flux, tableau, texte, graphique) et sa source.
\item \textbf{Dégager les directions des flux} : départs (campagnes pauvres, Nord surpeuplé/sécheresse, Est enclavé, Ouest saturé) et arrivées (Douala, Yaoundé, zones de plantations du Sud-Ouest/Littoral, frontières agricoles de l'Adamaoua/Est).
\item \textbf{Expliquer les causes} : facteurs de répulsion (pauvreté, sécheresse, manque de terres, insécurité) et facteurs d'attraction (emplois, services, sols fertiles, éducation).
\item \textbf{Montrer les conséquences} : brassage ethnique et culturel dans les villes, urbanisation rapide, métissage, mais aussi tensions foncières, pression sur les services urbains, bidonvilles.
\item \textbf{Relier à l'intégration nationale} : le brassage fabrique une identité camerounaise commune (mariages mixtes, langues véhiculaires : français, anglais, pidgin, fulfuldé, école, administration, service militaire).
\end{enumerate}
\end{methode}

\begin{exoresolu}[Commentaire : les migrations vers Douala et Yaoundé]
\textbf{Énoncé.} « Chaque année, des dizaines de milliers de jeunes quittent les régions du Nord et de l'Ouest pour s'installer à Douala et à Yaoundé. » À partir de ce constat et de vos connaissances, montrez que la mobilité des populations est un facteur de brassage et d'intégration nationale au Cameroun. (10 points)
\tcblower
\textbf{Solution détaillée.}
\begin{enumerate}
\item \textbf{Constat.} Le Cameroun connaît une grande mobilité interne : exode rural massif, migrations saisonnières de travail (plantations), migrations pendulaires et installations définitives vers les métropoles.
\item \textbf{Directions et causes.} Les flux partent des zones rurales surpeuplées ou fragilisées (hautes terres de l'Ouest, savanes du Nord touchées par la sécheresse et l'insécurité) vers les pôles d'attraction : Douala (port, industries, commerce informel) et Yaoundé (administration, services, universités). Les facteurs de répulsion (chômage rural, pauvreté, pression foncière, insécurité) se combinent aux facteurs d'attraction (emplois, éducation, santé, anonymat urbain).
\item \textbf{Brassage.} Dans les villes, toutes les ethnies se côtoient quotidiennement : un quartier de Douala ou Yaoundé mêle Bamiléké, Beti, Duala, Bassa, Peuls, Arabes Choa, anglophones du Nord-Ouest/Sud-Ouest... Ce voisinage de fait produit des mariages mixtes, la diffusion des langues véhiculaires (français, anglais, pidgin, fulfuldé) et des pratiques culturelles partagées (cuisine, musique, habillement).
\item \textbf{Intégration nationale.} L'école (programmes communs), l'administration (fonctionnaires mutés), l'armée (service national, brassage des recrues) et le marché du travail brassent les jeunes générations et construisent une citoyenneté commune au-delà des appartenances ethniques ou régionales d'origine.
\item \textbf{Limites et défis.} Ce brassage crée aussi des tensions : concurrence pour l'accès à la terre et à l'emploi, replis identitaires ou communautarismes, prolifération de bidonvilles, insécurité urbaine.
\item \textbf{Conclusion.} La mobilité est donc un double processus : elle mélange les populations et renforce l'unité nationale par le bas, tout en posant des défis d'aménagement (logement, emploi, cohésion sociale) que les politiques publiques doivent accompagner.
\end{enumerate}
\end{exoresolu}

\courssub{Distinguer paléo-soudanais, néo-soudanais et bantou dans un commentaire}

\begin{methode}[Construire un tableau comparatif de groupes humains]
\begin{enumerate}
\item \textbf{Choisir des critères de comparaison} fixes : localisation, habitat, activités économiques, organisation sociale, religion.
\item \textbf{Remplir le tableau} avec des traits précis et localisés (exemples de peuples, de régions, de villes).
\item \textbf{Dégager ensuite} en rédaction une ou deux idées générales (opposition savane/forêt, sédentarité/nomadisme, islamisation/christianisation).
\item \textbf{Ne jamais essentialiser} : préciser que ces caractères sont des traits dominants historiques, bousculés aujourd'hui par l'urbanisation, l'école et le brassage des populations.
\end{enumerate}
\end{methode}

\begin{exoresolu}[Tableau comparatif des grands groupes humains]
\textbf{Énoncé.} Présentez dans un tableau les différences entre paléo-soudanais, néo-soudanais, bantou et semi-bantou, pygmées (localisation, activités, organisation socio-religieuse), puis dégagez l'idée principale. (10 points)
\tcblower
\textbf{Solution détaillée.}
\begin{center}
\begin{tabularx}{\linewidth}{|l|X|X|X|}
\hline
\rowcolor{popblueL} \textbf{Groupe} & \textbf{Localisation} & \textbf{Activités principales} & \textbf{Organisation socio-religieuse} \\
\hline
Paléo-soudanais & Monts Mandara et plaines de l'Extrême-Nord (Mafa, Mofu, Kapsiki, Matakam, Kotoko*) & Agriculture sur terrasses, élevage, pêche (Kotoko) & Chefferies de montagne, religions traditionnelles, résistance historique à l'islamisation \\
\hline
Néo-soudanais & Plaines du Nord (Bénoué) et de l'Adamaoua (Peuls, Arabes Choa, Kanouri, Moundang) & Élevage bovin transhumant, agriculture, commerce & Islam dominant, lamibés (chefferies islamisées), société hiérarchisée (nobles, artisans, captifs) \\
\hline
Bantou et semi-bantou & Sud forestier (Beti, Fang, Bassa, Duala, Bakoko) et hautes terres de l'Ouest/Nord-Ouest (Bamiléké, Bamoun, Tikar, groupes des Grassfields) & Agriculture sur brûlis (forêt), agriculture intensive volcans (Ouest), plantations, commerce & Chefferies traditionnelles, sociétés secrètes, christianisme (catholique, protestant) et religions traditionnelles \\
\hline
Pygmées & Forêt dense du Sud-Est (Baka, Bagyeli, Medzan) & Chasse, cueillette, pêche, semi-nomadisme & Groupes de campements égalitaires, religions animistes liées à l'esprit de la forêt (Jengi) \\
\hline
\end{tabularx}
\end{center}
\textbf{Note :} *Les Kotoko, héritiers de la civilisation Sao, sont souvent classés parmi les paléo-soudanais mais sont largement islamisés aujourd'hui.
\textbf{Idée principale.} L'opposition fondamentale structurelle est celle entre les peuples de la \emph{savane} (soudanais, au nord de l'Adamaoua, marqués par l'élevage, le mil/sorgho et l'islam pour les néo-soudanais) et les peuples de la \emph{forêt} (bantou et pygmées, au sud, marqués par l'agriculture sur brûlis, les tubercules/plantains, le christianisme ou l'animisme). Les semi-bantou des hautes terres de l'Ouest forment une zone de transition originale (volcans, forte densité, chefferies dynamiques). Mais l'urbanisation et les migrations estompent aujourd'hui ces contrastes : les groupes se mélangent dans les villes, ce qui nourrit le brassage national.
\textbf{Conclusion.} La diversité des groupes humains, héritée de l'histoire des migrations (bantoues venues du Nord-Est, soudanaises venues du Nord), fait du Cameroun « l'Afrique en miniature » et constitue à la fois une richesse culturelle exceptionnelle et un défi permanent d'intégration nationale.
\end{exoresolu}

\begin{attention}[Les 5 erreurs qui coûtent des points au Bac]
\begin{enumerate}
\item \textbf{Confondre les groupes et leurs localisations} : placer les pygmées dans l'Ouest ou les soudanais au Sud est une erreur éliminatoire. Retiens l'axe nord-sud : soudanais au nord de l'Adamaoua, bantou et semi-bantou au sud, pygmées dans la forêt du Sud-Est.
\item \textbf{Oublier la distinction paléo-soudanais / néo-soudanais} : les premiers (Kirdi des monts Mandara : Mafa, Mofu, Kapsiki) sont restés attachés aux religions traditionnelles, les seconds (Peuls, Arabes Choa) sont islamisés et organisés en lamibés. Ne pas confondre.
\item \textbf{Se tromper de formule de densité} : c'est population divisée par superficie, jamais l'inverse. Vérifie la cohérence du résultat : une densité de 0,14 hab/km$^2$ pour le Littoral signale une inversion.
\item \textbf{Oublier les éléments obligatoires de la carte (TP)} : titre précis, flèche du nord, légende \textbf{chiffrée} (classes de densité avec leurs valeurs numériques), échelle indicative. Une carte sans légende chiffrée perd la moitié des points du TP.
\item \textbf{Essentialiser les groupes ethniques} : écrire « les Bamiléké sont des commerçants » ou « les Peuls sont des éleveurs » sans nuance est une faute d'analyse géographique. Précise toujours qu'il s'agit de traits dominants historiques, transformés par l'urbanisation, l'école, la diversification des activités et le brassage des populations.
\end{enumerate}
\end{attention}

\newpage

\coursec{Exercices}

\courssub{Je vérifie mes connaissances}

\begin{exercice}[QCM : les grands groupes humains]
Pour chaque question, une seule réponse est exacte. Entoure la bonne lettre.

\textbf{1.} Les peuples dits \og soudanais \fg{} occupent principalement :
\begin{itemize}
\item[a)] le littoral et le Sud forestier ;
\item[b)] les régions du Nord (Nord, Extrême-Nord, Adamaoua) ;
\item[c)] uniquement les hautes terres de l'Ouest.
\end{itemize}

\textbf{2.} Les Kotoko, installés autour du lac Tchad et du Logone, appartiennent au groupe :
\begin{itemize}
\item[a)] néo-soudanais ; \quad b)] paléo-soudanais ; \quad c)] bantou.
\end{itemize}

\textbf{3.} Les pygmées (Baka, Bagyeli) vivent traditionnellement :
\begin{itemize}
\item[a)] dans la plaine du Bénoué ;
\item[b)] dans la forêt du Sud et de l'Est du Cameroun ;
\item[c)] dans les monts Mandara.
\end{itemize}

\textbf{4.} La région du Cameroun dont la densité de population dépasse couramment 100 habitants au km\textsuperscript{2} est :
\begin{itemize}
\item[a)] l'Est ; \quad b)] l'Adamaoua ; \quad c)] l'Ouest.
\end{itemize}
\end{exercice}

\begin{exercice}[Vrai ou faux — à justifier]
Indique si chaque affirmation est vraie ou fausse, puis justifie ta réponse en une ou deux phrases.
\begin{enumerate}
\item Les néo-soudanais sont organisés en chefferies centralisées (lamidats) marquées par l'islamisation.
\item Les peuples paléo-soudanais des monts Mandara ont fui historiquement les plaines pour échapper aux razzias.
\item La région de l'Est est l'une des plus densément peuplées du Cameroun.
\item Les migrations intérieures ne concernent que les jeunes ruraux à la recherche d'emplois industriels.
\end{enumerate}
\end{exercice}

\begin{exercice}[Définitions]
Définis en deux ou trois lignes chacun des termes suivants, en donnant pour chacun un exemple camerounais : \cle{groupe ethnique} ; \cle{paléo-soudanais} ; \cle{semi-bantou} ; \cle{densité de population} ; \cle{brassage des populations}.
\end{exercice}

\begin{exercice}[Calculs de densités]
Le tableau ci-dessous donne des ordres de grandeur pour trois régions du Cameroun.

\begin{center}
\begin{tabularx}{\linewidth}{|l|X|X|X|}
\hline
\rowcolor{popblueL} Région & Superficie (km\textsuperscript{2}) & Population (habitants) & Densité (hab/km\textsuperscript{2}) \\
\hline
Ouest & \num{14000} & \num{2000000} & \ldots \\
\hline
Adamaoua & \num{64000} & \num{1200000} & \ldots \\
\hline
Littoral & \num{20000} & \num{3300000} & \ldots \\
\hline
\end{tabularx}
\end{center}

\begin{enumerate}
\item Calcule la densité de population de chaque région (arrondis à l'unité).
\item Classe ces trois régions de la plus densément peuplée à la moins densément peuplée.
\item Calcule la part de la population du Littoral dans la population totale des trois régions (en \%, arrondi au dixième).
\end{enumerate}
\end{exercice}

\courssub{Je m'entraîne}

\begin{exercice}[Paléo-soudanais et néo-soudanais : question à réponse courte]
Distingue les \cle{paléo-soudanais} des \cle{néo-soudanais} en donnant pour chaque groupe :
\begin{enumerate}
\item deux exemples de peuples du Cameroun ;
\item leur espace de vie principal ;
\item un trait caractéristique de leur organisation socio-politique.
\end{enumerate}
Ta réponse ne doit pas dépasser une demi-page.
\end{exercice}

\begin{exercice}[Classer les peuples]
Voici une liste de peuples du Cameroun : Bamiléké, Kotoko, Douala, Mafa, Peul (Foulbé), Baka, Toupouri, Bétis, Tikar, Massa, Ewondo, Kirdi (montagnards).
\begin{enumerate}
\item Range chaque peuple dans la colonne qui convient : soudanais (paléo ou néo), bantou, semi-bantou, pygmées.
\item Pour chaque grand groupe, indique la ou les régions du Cameroun où il est majoritairement implanté.
\item Identifie deux peuples de cette liste que l'on retrouve aujourd'hui en nombre important à Yaoundé ou à Douala, et explique pourquoi.
\end{enumerate}
\end{exercice}

\begin{exercice}[Lecture d'un tableau de densités]
Le tableau ci-dessous donne les densités approximatives des dix régions du Cameroun (ordres de grandeur).

\begin{center}
\begin{tabularx}{\linewidth}{|l|X|l|X|}
\hline
\rowcolor{popblueL} Région & Densité (hab/km\textsuperscript{2}) & Région & Densité (hab/km\textsuperscript{2}) \\
\hline
Ouest & environ 140 & Centre & environ 50 \\
\hline
Nord-Ouest & environ 90 & Nord & environ 45 \\
\hline
Littoral & environ 160 & Extrême-Nord & environ 80 \\
\hline
Sud-Ouest & environ 45 & Adamaoua & environ 18 \\
\hline
Sud & environ 25 & Est & environ 10 \\
\hline
\end{tabularx}
\end{center}

\begin{enumerate}
\item Classe les régions en trois groupes : densément peuplées (plus de 75 hab/km\textsuperscript{2}), moyennement peuplées (25 à 75 hab/km\textsuperscript{2}), faiblement peuplées (moins de 25 hab/km\textsuperscript{2}).
\item Représente ces dix densités par un diagramme en barres ordonné (régions sur l'axe horizontal, densités sur l'axe vertical). N'oublie ni le titre ni les unités.
\item Propose une explication géographique (naturelle ou historique) de la forte densité de l'Ouest et de la faible densité de l'Est.
\end{enumerate}
\end{exercice}

\begin{exercice}[Croquis des migrations intérieures]
Sur un fond de carte schématique du Cameroun (contour simplifié, dix régions délimitées, nord fléché) :
\begin{enumerate}
\item place et nomme les principales zones de départ (monts Mandara, hautes terres de l'Ouest et du Nord-Ouest) et les principales zones d'accueil (Douala, Yaoundé, plantations du littoral, Adamaoua, périphéries urbaines) ;
\item trace par des flèches d'épaisseur proportionnelle les grands courants migratoires intérieurs ;
\item rédige une légende numérotée et un titre ;
\item en cinq lignes, explique comment ces migrations favorisent le brassage des populations et l'intégration nationale.
\end{enumerate}
\end{exercice}

\courssub{Je me prépare au Bac}

\begin{exercice}[Dissertation type Bac]
\textbf{Sujet :} \og Les grands groupes humains du Cameroun : une diversité au service de l'intégration nationale ? \fg

Tu rédigeras une dissertation géographique complète comprenant :
\begin{enumerate}
\item une introduction avec accroche, définition des termes clés, présentation du cadre spatial et temporel, problématique et annonce du plan ;
\item un développement en deux ou trois parties clairement articulées (on pourra par exemple montrer d'abord la réalité de la mosaïque humaine et de l'inégale répartition, puis les dynamiques de brassage et leurs limites) ;
\item une conclusion qui répond explicitement à la problématique et propose une ouverture.
\end{enumerate}
\emph{Barème indicatif :} introduction 4 pts ; développement 12 pts ; conclusion 2 pts ; présentation et langue 2 pts.
\end{exercice}

\begin{exercice}[Analyse de documents : peuplement historique et répartition actuelle]
\textbf{Document 1 (carte décrite).} Une carte des grands groupes humains du Cameroun montre : au nord, les peuples soudanais (paléo-soudanais des monts Mandara — Mafa, Mofu, Kapsiki — et néo-soudanais des plaines — Peuls, Kotoko, Arabes Choa) ; à l'ouest et au nord-ouest, les peuples semi-bantou des Grassfields (Bamiléké, Bamoun, Tikar) ; au centre, au sud et à l'est, les peuples bantou (Bétis, Ewondo, Douala, Bassa, Bulu) ; en forêt du Sud-Est, les pygmées (Baka, Bagyeli).

\textbf{Document 2 (tableau).} Densités régionales approximatives (hab/km\textsuperscript{2}) : Littoral 160 ; Ouest 140 ; Nord-Ouest 90 ; Extrême-Nord 80 ; Centre 50 ; Sud-Ouest 45 ; Nord 45 ; Sud 25 ; Adamaoua 18 ; Est 10.

\begin{enumerate}
\item À partir du document 1, décris la localisation des quatre grands groupes humains du Cameroun (4 pts).
\item À partir du document 2, dégage les trois grandes classes de densité et cite les régions correspondantes (3 pts).
\item Montre, en croisant les deux documents, que les fortes densités actuelles correspondent souvent à d'anciens foyers de peuplement (monts Mandara, hautes terres de l'Ouest) ; propose deux explications historiques et deux explications naturelles (6 pts).
\item Explique pourquoi le Littoral et le Centre connaissent des densités élevées alors qu'ils ne sont pas des foyers de peuplement ancien comparables aux précédents (4 pts).
\item Conclus en montrant en quoi la connaissance du peuplement historique éclaire l'organisation actuelle de l'espace national (3 pts).
\end{enumerate}
\end{exercice}

\begin{exercice}[Construction cartographique : la carte des densités (TP 2)]
On te fournit le tableau suivant (ordres de grandeur) :

\begin{center}
\begin{tabularx}{\linewidth}{|l|X|X|X|}
\hline
\rowcolor{popblueL} Région & Densité (hab/km\textsuperscript{2}) & Région & Densité (hab/km\textsuperscript{2}) \\
\hline
Extrême-Nord & 80 & Nord-Ouest & 90 \\
\hline
Nord & 45 & Ouest & 140 \\
\hline
Adamaoua & 18 & Littoral & 160 \\
\hline
Centre & 50 & Sud-Ouest & 45 \\
\hline
Est & 10 & Sud & 25 \\
\hline
\end{tabularx}
\end{center}

\begin{enumerate}
\item Choisis un découpage en trois classes (densément, moyennement, faiblement peuplé) et justifie tes seuils (3 pts).
\item Construis la carte : contour simplifié du Cameroun, dix régions délimitées, chaque classe représentée par une teinte ou un tramage distinct, flèche du nord, échelle indicative (8 pts).
\item Rédige une légende complète et numérotée, ainsi qu'un titre précis (3 pts).
\item Lis ta carte : localise les trois ensembles de peuplement (Nord soudanais, hautes terres de l'Ouest, axe Douala--Yaoundé ; diagonale vide Adamaoua--Est--Sud) et propose une explication pour chacun (4 pts).
\item Indique deux limites de ta carte (choix des seuils, échelle régionale masquant les contrastes locaux) (2 pts).
\end{enumerate}
\end{exercice}

\newpage

\coursec{Corrigés des exercices}

\begin{corrige}[Exercice 1 — QCM : les grands groupes humains]
\textbf{1. Réponse b).} Les peuples soudanais occupent les régions septentrionales du Cameroun (Extrême-Nord, Nord, Adamaoua), dans le domaine des savanes soudaniennes. Le littoral et le Sud forestier sont le domaine des peuples bantou ; les hautes terres de l'Ouest, celui des semi-bantou.

\textbf{2. Réponse b).} Les Kotoko, descendants des anciens États du Sao installés autour du lac Tchad et du Logone, sont des \cle{paléo-soudanais} : ils occupent le pays depuis longtemps, avant l'islamisation. Les néo-soudanais (Peuls, Arabes Choa) sont arrivés plus tardivement, avec l'islam.

\textbf{3. Réponse b).} Les pygmées (Baka, Bagyeli, Bakola) vivent traditionnellement dans la forêt dense du Sud et de l'Est du Cameroun, où ils pratiquent la chasse, la cueillette et la pêche.

\textbf{4. Réponse c).} La région de l'Ouest dépasse couramment 100 hab/km\textsuperscript{2} (environ 140 hab/km\textsuperscript{2}), grâce à ses sols volcaniques fertiles et à un peuplement ancien. L'Est (environ 10 hab/km\textsuperscript{2}) et l'Adamaoua (environ 18 hab/km\textsuperscript{2}) sont au contraire faiblement peuplés.
\end{corrige}

\begin{corrige}[Exercice 2 — Vrai ou faux]
\begin{enumerate}
\item \textbf{Vrai.} Les néo-soudanais (Peuls notamment) ont constitué à partir du XIX\textsuperscript{e} siècle des chefferies centralisées, les \cle{lamidats} (Ngaoundéré, Garoua, Maroua, Banyo), organisés hiérarchiquement autour du lamido et fortement marqués par l'islam.
\item \textbf{Vrai.} Les paléo-soudanais des monts Mandara (Mafa, Mofu, Kapsiki, etc.) se sont réfugiés dans les reliefs pour échapper aux razzias esclavagistes des cavaliers peuls et aux conquêtes des lamidats ; la montagne leur a servi de forteresse naturelle.
\item \textbf{Faux.} L'Est est la région la moins densément peuplée du Cameroun (environ 10 hab/km\textsuperscript{2}) : forêt dense difficile à défricher, enclavement, maladies (paludisme, trypanosomiase) et éloignement des axes majeurs expliquent ce vide relatif.
\item \textbf{Faux.} Les migrations intérieures concernent toutes les catégories : jeunes ruraux vers les villes, mais aussi éleveurs peuls vers l'Adamaoua et le Sud, planteurs bamiléké vers le littoral et le Sud-Ouest, commerçants, fonctionnaires mutés, etc. Les motifs sont économiques, scolaires, administratifs et matrimoniaux.
\end{enumerate}
\end{corrige}

\begin{corrige}[Exercice 3 — Définitions]
\begin{itemize}
\item \textbf{Groupe ethnique :} ensemble humain partageant une langue, des coutumes, des institutions et souvent une origine commune revendiquée. Exemple : les Bamiléké des hautes terres de l'Ouest.
\item \textbf{Paléo-soudanais :} peuples anciennement installés dans les savanes et les montagnes du Nord-Cameroun, avant l'islamisation ; beaucoup se sont réfugiés dans les monts Mandara. Exemple : les Mafa, les Kotoko.
\item \textbf{Semi-bantou :} peuples des Grassfields (Ouest, Nord-Ouest) dont les langues et les organisations (chefferies hiérarchisées, sociétés secrètes) présentent des traits intermédiaires entre les mondes soudanais et bantou. Exemple : les Bamiléké, les Bamoun, les Tikar.
\item \textbf{Densité de population :} nombre moyen d'habitants par unité de surface, obtenu en divisant la population par la superficie ($D = P/S$, en hab/km\textsuperscript{2}). Exemple : environ 140 hab/km\textsuperscript{2} dans l'Ouest.
\item \textbf{Brassage des populations :} mélange des peuples résultant des migrations intérieures, des mariages mixtes et de la cohabitation dans les villes et les zones pionnières. Exemple : la présence de toutes les ethnies du pays à Douala et à Yaoundé.
\end{itemize}
\end{corrige}

\begin{corrige}[Exercice 4 — Calculs de densités]
\textbf{1. Calcul des densités} ($D = P/S$) :
\begin{align*}
D_{\text{Ouest}} &= \frac{\num{2000000}}{\num{14000}} \approx 142,9 \approx 143 \text{ hab/km}^2 \\
D_{\text{Adamaoua}} &= \frac{\num{1200000}}{\num{64000}} = 18,75 \approx 19 \text{ hab/km}^2 \\
D_{\text{Littoral}} &= \frac{\num{3300000}}{\num{20000}} = 165 \text{ hab/km}^2
\end{align*}

\textbf{2. Classement} (de la plus dense à la moins dense) : Littoral (165) $>$ Ouest (143) $>$ Adamaoua (19).

\textbf{3. Part du Littoral} dans la population totale des trois régions :
\[ P_{\text{totale}} = \num{2000000} + \num{1200000} + \num{3300000} = \num{6500000} \text{ habitants} \]
\[ \text{Part} = \frac{\num{3300000}}{\num{6500000}} \times 100 \approx 50,8\,\% \]
Le Littoral regroupe à lui seul plus de la moitié de la population des trois régions, alors qu'il n'occupe qu'environ 20\,\% de leur superficie totale : c'est un foyer de concentration démographique lié au port de Douala et aux activités industrielles et commerciales.

\textbf{Point de vigilance :} ne jamais additionner des densités ; on additionne des populations et des superficies, puis on divise.
\end{corrige}

\begin{corrige}[Exercice 5 — Paléo-soudanais et néo-soudanais]
\textbf{Paléo-soudanais.}
\begin{itemize}
\item Peuples : Mafa et Kotoko (on peut citer aussi Mofu, Kapsiki, Massa, Toupouri).
\item Espace de vie : monts Mandara (pour les montagnards réfugiés) et plaine du Logone--lac Tchad (Kotoko, Massa).
\item Organisation : sociétés villageoises décentralisées, sans État centralisé ; autorité des chefs de terre et des notables ; religions traditionnelles (islam peu ou pas implanté chez les montagnards).
\end{itemize}
\textbf{Néo-soudanais.}
\begin{itemize}
\item Peuples : Peuls (Foulbé) et Arabes Choa.
\item Espace de vie : plaines de la Bénoué, du Logone et du Diamaré (Nord, Extrême-Nord, Adamaoua).
\item Organisation : chefferies centralisées et hiérarchisées, les \cle{lamidats} (Garoua, Maroua, Ngaoundéré), dirigées par un lamido ; forte islamisation.
\end{itemize}
\textbf{Idée directrice :} la distinction repose sur l'ancienneté de l'installation (avant/après l'islamisation du XIX\textsuperscript{e} siècle) et sur le type d'organisation politique (villages autonomes contre États centralisés).
\end{corrige}

\begin{corrige}[Exercice 6 — Classer les peuples]
\textbf{1. Classement.}
\begin{center}
\begin{tabularx}{\linewidth}{|l|X|}
\hline
\rowcolor{popblueL} Groupe & Peuples de la liste \\
\hline
Paléo-soudanais & Kotoko, Mafa, Toupouri, Massa, Kirdi (montagnards) \\
\hline
Néo-soudanais & Peul (Foulbé) \\
\hline
Semi-bantou & Bamiléké, Tikar \\
\hline
Bantou & Douala, Bétis, Ewondo \\
\hline
Pygmées & Baka \\
\hline
\end{tabularx}
\end{center}

\textbf{2. Localisation.}
\begin{itemize}
\item Paléo-soudanais : Extrême-Nord (monts Mandara, plaine du Logone) et Nord (plaine de la Bénoué pour les Toupouri).
\item Néo-soudanais : Nord, Extrême-Nord et Adamaoua (plaines et lamidats).
\item Semi-bantou : Ouest et Nord-Ouest (Grassfields, hautes terres).
\item Bantou : Centre, Sud, Littoral, Sud-Ouest et Est (domaine forestier et côtier).
\item Pygmées : forêts du Sud et de l'Est.
\end{itemize}

\textbf{3. Peuples présents en nombre à Yaoundé ou Douala.} Les Bamiléké et les Ewondo (mais aussi les Bétis et les Douala) sont très nombreux dans les deux métropoles. Les Bamiléké y ont migré massivement depuis les hautes terres surpeuplées de l'Ouest pour le commerce, l'artisanat et les emplois ; les Ewondo et les Bétis, originaires du Centre, sont présents autour de Yaoundé, capitale administrative. C'est l'illustration du \cle{brassage des populations} par les migrations intérieures.
\end{corrige}

\begin{corrige}[Exercice 7 — Lecture d'un tableau de densités]
\textbf{1. Classement en trois groupes.}
\begin{itemize}
\item \textbf{Densément peuplées (plus de 75 hab/km\textsuperscript{2}) :} Littoral (160), Ouest (140), Nord-Ouest (90), Extrême-Nord (80).
\item \textbf{Moyennement peuplées (25 à 75 hab/km\textsuperscript{2}) :} Centre (50), Nord (45), Sud-Ouest (45), Sud (25).
\item \textbf{Faiblement peuplées (moins de 25 hab/km\textsuperscript{2}) :} Adamaoua (18), Est (10).
\end{itemize}

\textbf{2. Diagramme en barres ordonné.}

\begin{popfigure}
\begin{center}
\begin{tikzpicture}
\begin{axis}[
  width=0.95\linewidth, height=7cm,
  ybar, bar width=9pt,
  ymin=0, ymax=180,
  ylabel={Densité (hab/km\textsuperscript{2})},
  symbolic x coords={Littoral,Ouest,Nord-Ouest,Extrême-Nord,Centre,Nord,Sud-Ouest,Sud,Adamaoua,Est},
  xtick=data,
  x tick label style={rotate=45, anchor=east, font=\small},
  ymajorgrids=true, grid style={popline},
  every axis plot/.append style={fill=popblue, draw=popdark}
]
\addplot coordinates {(Littoral,160) (Ouest,140) (Nord-Ouest,90) (Extrême-Nord,80) (Centre,50) (Nord,45) (Sud-Ouest,45) (Sud,25) (Adamaoua,18) (Est,10)};
\end{axis}
\end{tikzpicture}
\end{center}
\legende{Diagramme en barres des densités régionales de population au Cameroun (ordres de grandeur, hab/km\textsuperscript{2}). Les régions sont classées par densité décroissante.}
\end{popfigure}

\textbf{3. Explications.}
\begin{itemize}
\item \textbf{Ouest (forte densité) :} sols volcaniques très fertiles, climat d'altitude sain, peuplement ancien des Grassfields, agriculture vivrière et de rente intensive (café, maraîchage), réseau dense de marchés.
\item \textbf{Est (faible densité) :} forêt dense équatoriale difficile à exploiter, enclavement (routes rares), maladies endémiques, éloignement des grands axes et des ports ; l'économie (exploitation forestière, artisanat) offre peu d'emplois fixes.
\end{itemize}
\end{corrige}

\begin{corrige}[Exercice 8 — Croquis des migrations intérieures]
\textbf{1--3. Croquis, flèches, titre et légende.}

\begin{popfigure}
\begin{center}
\includegraphics[width=\linewidth,height=10cm,keepaspectratio]{N02/cmr_regions.jpg}
\end{center}
\legende{\textbf{Titre :} Les grands courants de migrations intérieures au Cameroun. Fond de carte : grands ensembles régionaux du Cameroun (nord, Adamaoua, hautes terres, littoral, plateau sud), grandes villes, fleuves et routes (noms en anglais). Le croquis attendu y situe : \textbf{1} : la zone de départ des monts Mandara (Extrême-Nord, pression foncière) ; \textbf{2} : la zone de départ des hautes terres de l'Ouest et du Nord-Ouest (Bamenda, Bafoussam) ; \textbf{3} : Douala, pôle portuaire et industriel ; \textbf{4} : Yaoundé, capitale administrative ; \textbf{5} : l'Adamaoua, zone d'accueil des éleveurs. Il y ajoute des flèches de migrations de travail de 1 et 2 vers 3 et 4, et une flèche d'éleveurs vers le Sud ; l'épaisseur est proportionnelle à l'importance du flux. \\ Source du fond de carte : Wikimedia Commons, Peter Fitzgerald, CC BY 3.0.}
\end{popfigure}

\textbf{4. Explication (cinq lignes).} Ces migrations mélangent durablement les peuples : Peuls, montagnards du Nord, Bamiléké et Bamoun de l'Ouest se retrouvent à Douala, à Yaoundé, dans les plantations et les nouvelles zones pionnières. La cohabitation quotidienne, les mariages mixtes, l'usage du français, de l'anglais et des langues véhiculaires créent des solidarités qui dépassent l'appartenance ethnique. Les villes deviennent ainsi des creusets d'intégration nationale, même si des tensions foncières et identitaires subsistent localement.
\end{corrige}

\begin{corrige}[Exercice 9 — Dissertation type Bac]
\textbf{Corrigé détaillé (plan et contenus attendus).}

\textbf{Introduction (4 pts).} Accroche : au marché Mokolo de Yaoundé ou au marché central de Douala, toutes les langues du pays se croisent. Définitions : \cle{grands groupes humains} (soudanais, bantou, semi-bantou, pygmées) ; \cle{intégration nationale} (construction d'une communauté unie au-delà des appartenances). Cadre : Cameroun indépendant, plus de 250 groupes ethniques, de 1960 à nos jours. Problématique : la diversité ethnique et l'inégale répartition de la population sont-elles un obstacle ou un atout pour l'unité nationale ? Annonce du plan en deux parties.

\textbf{Développement (12 pts).}

\emph{I. Une mosaïque humaine inégalement répartie (6 pts).}
\begin{itemize}
\item Quatre grands groupes localisés : soudanais au Nord (paléo-soudanais des monts Mandara et du Logone : Mafa, Kotoko ; néo-soudanais des plaines islamisées : Peuls, Arabes Choa) ; semi-bantou des Grassfields (Bamiléké, Bamoun, Tikar) ; bantou du Centre-Sud-Littoral (Bétis, Ewondo, Douala, Bassa) ; pygmées de la forêt (Baka, Bagyeli).
\item Des organisations socio-culturelles contrastées : lamidats centralisés et islamisés ; villages autonomes montagnards ; chefferies hiérarchisées des Grassfields ; sociétés lignagères bantou ; bandes nomades pygmées.
\item De forts contrastes de densité : Littoral (160), Ouest (140) contre Est (10) et Adamaoua (18) ; une diagonale quasi vide du Sud-Est.
\end{itemize}

\emph{II. Des dynamiques de brassage au service de l'intégration, mais avec des limites (6 pts).}
\begin{itemize}
\item Migrations intérieures : départs des monts Mandara et des hautes terres vers Douala, Yaoundé, les plantations, l'Adamaoua ; brassage urbain, mariages mixtes, langues véhiculaires.
\item Facteurs d'unité : État centralisé, école, français et anglais, économie de marché, équilibre régional recherché dans les nominations.
\item Limites : tensions foncières entre autochtones et allogènes, régionalisme politique, marginalisation des pygmées, exode rural qui vide les campagnes.
\end{itemize}

\textbf{Conclusion (2 pts).} La diversité est une richesse lorsque l'État et les migrations la transforment en brassage ; l'intégration reste un chantier permanent. Ouverture : le rôle des nouvelles technologies et de la jeunesse urbaine dans l'invention d'une identité camerounaise commune.

\textbf{Points de vigilance :} problématique explicitement formulée et reprise en conclusion ; exemples précis (peuples, régions, chiffres) dans chaque paragraphe ; transitions rédigées ; éviter le catalogue sans analyse.
\end{corrige}

\begin{corrige}[Exercice 10 — Analyse de documents]
\textbf{1. Localisation des quatre groupes (4 pts).} Les soudanais occupent le Nord : paléo-soudanais dans les monts Mandara et autour du Logone (Mafa, Mofu, Kapsiki, Kotoko), néo-soudanais dans les plaines islamisées (Peuls, Arabes Choa). Les semi-bantou occupent les Grassfields de l'Ouest et du Nord-Ouest (Bamiléké, Bamoun, Tikar). Les bantou couvrent le Centre, le Sud, le Littoral et l'Est (Bétis, Ewondo, Douala, Bassa, Bulu). Les pygmées vivent dans la forêt du Sud-Est (Baka, Bagyeli).

\textbf{2. Trois classes de densité (3 pts).}
\begin{itemize}
\item Densément peuplées ($> 75$) : Littoral (160), Ouest (140), Nord-Ouest (90), Extrême-Nord (80).
\item Moyennement peuplées (25--75) : Centre (50), Nord (45), Sud-Ouest (45), Sud (25).
\item Faiblement peuplées ($< 25$) : Adamaoua (18), Est (10).
\end{itemize}

\textbf{3. Fortes densités et foyers anciens (6 pts).} Les hautes terres de l'Ouest (foyer semi-bantou ancien) et l'Extrême-Nord (foyer paléo-soudanais des monts Mandara) figurent parmi les régions les plus denses : le peuplement s'est accumulé là où l'installation est ancienne. \emph{Explications historiques :} refuge dans les montagnes contre les razzias (concentration forcée) ; ancienneté des chefferies et des marchés des Grassfields, qui ont fixé et multiplié les populations. \emph{Explications naturelles :} sols volcaniques fertiles et climat d'altitude sain dans l'Ouest ; terres de montagne protégées et eau disponible dans les Mandara, plaines inondables propices à la pêche et aux cultures de décrue chez les Kotoko.

\textbf{4. Le cas du Littoral et du Centre (4 pts).} Ces régions doivent leurs fortes densités non à un peuplement ancien comparable, mais à la \emph{colonisation récente de l'espace} : Douala, premier port et capitale économique (industries, commerce, plantations), attire des migrants de tout le pays ; Yaoundé, capitale politique et administrative, concentre fonctionnaires, étudiants et services. Ce sont des densités \emph{acquises par les migrations}, et non héritées.

\textbf{5. Conclusion (3 pts).} Le peuplement historique explique les foyers denses du Nord et de l'Ouest ; les dynamiques contemporaines (urbanisation, migrations) expliquent ceux du Littoral et du Centre. Comprendre cette double origine permet d'organiser l'espace national : aménager les zones vides (Est, Adamaoua), désengorger les foyers surpeuplés, et faire du brassage un levier d'intégration.

\textbf{Points de vigilance :} citer les documents à l'appui (\og le document 2 montre que\ldots\fg) ; ne pas confondre foyer ancien et densité actuelle ; chiffres exacts repris sans invention.
\end{corrige}

\begin{corrige}[Exercice 11 — Construction cartographique : la carte des densités (TP 2)]
\textbf{1. Choix des classes (3 pts).} On retient trois classes aux seuils simples et justifiés par la répartition des valeurs (rupture nette entre 80 et 90 d'une part, entre 25 et 45 d'autre part) : \textbf{densément peuplé} : plus de 75 hab/km\textsuperscript{2} ; \textbf{moyennement peuplé} : 25 à 75 hab/km\textsuperscript{2} ; \textbf{faiblement peuplé} : moins de 25 hab/km\textsuperscript{2}.

\textbf{2--3. Carte, légende et titre (8 + 3 pts).}

\begin{popfigure}
\begin{center}
\includegraphics[width=\linewidth,height=10cm,keepaspectratio]{N02/cmr_dens.png}
\end{center}
\legende{\textbf{Titre :} La répartition de la population au Cameroun : classes de densité régionale. Fond de carte : densité de population par région en quatre classes (légende en anglais : plus de 50, 20 à 50, 10 à 20, moins de 10 hab/km\textsuperscript{2}). Le croquis attendu regroupe : \textbf{1} : les régions densément peuplées (bleu foncé : Extrême-Nord, Nord-Ouest, Ouest, Littoral) ; \textbf{2} : les régions moyennement peuplées (bleu clair et vert : Centre, Sud-Ouest, Nord) ; \textbf{3} : les régions faiblement peuplées (vert pâle : Adamaoua, Est, Sud). \\ Source : Wikimedia Commons, Austiger, CC BY-SA 4.0.}
\end{popfigure}

\textbf{4. Lecture de la carte (4 pts).} Trois ensembles se dégagent :
\begin{itemize}
\item le \textbf{Nord soudanais} (Extrême-Nord dense) : foyer ancien de peuplement (refuge des monts Mandara, plaines du Logone), malgré un milieu semi-aride ;
\item les \textbf{hautes terres de l'Ouest et du Nord-Ouest} : sols volcaniques fertiles, climat sain, chefferies anciennes et marchés dynamiques ;
\item l'\textbf{axe Douala--Yaoundé} (Littoral et Centre) : densités acquises par le port, l'industrie, l'administration et les migrations récentes.
\end{itemize}
En regard, la \textbf{diagonale vide Adamaoua--Est--Sud} s'explique par la forêt dense, l'enclavement, les maladies et l'éloignement des axes majeurs.

\textbf{5. Deux limites (2 pts).} Le choix des seuils est arbitraire : d'autres bornes modifieraient le classement de certaines régions. Surtout, l'échelle régionale masque les contrastes locaux : au sein de l'Est, les axes routiers sont plus peuplés que la forêt ; dans l'Ouest, certaines vallées sont moins denses que les plateaux. Une carte à l'échelle des départements affinerait l'analyse.

\textbf{Points de vigilance :} titre précis (quoi, où), légende numérotée avec les seuils chiffrés et les unités, nord fléché, échelle indicative ; ne jamais colorier sans légender.
\end{corrige}

\newpage

\coursec{Fiche bilan}

\begin{popfigure}
\begin{tikzpicture}[mindmap, grow cyclic, every node/.style={concept, font=\small},
  root concept/.append style={concept color=popblue, font=\bfseries},
  level 1/.append style={level distance=3.6cm, sibling angle=51.5},
  level 2/.append style={level distance=2.2cm, sibling angle=40, font=\scriptsize},
  text=white]
\node[root concept]{Les grands groupes humains du Cameroun}
  child[concept color=poporange]{ node {Soudanais}
    child { node {Paléo-soudanais : Kotoko, Masa, Toupouri} }
    child { node {Néo-soudanais : Peuls (Foulbé), arabes Choa} }
  }
  child[concept color=popgreen]{ node {Bantou}
    child { node {Centre, Sud, Littoral : Ewondo, Douala, Bassa} }
  }
  child[concept color=poppurple]{ node {Semi-bantou}
    child { node {Ouest et Nord-Ouest : Bamiléké, Bamoun, Tikars} }
  }
  child[concept color=poppink]{ node {Pygmées}
    child { node {Baka, Bagyeli, Bakola (forêt du Sud-Est)} }
  }
  child[concept color=popgold]{ node {Densités inégales}
    child { node {Fortes : Ouest, Littoral, Extrême-Nord} }
    child { node {Faibles : Est, Adamaoua, Sud} }
  }
  child[concept color=popblue!70]{ node {Mobilité et brassage}
    child { node {Migrations vers villes et plantations} }
    child { node {Intégration nationale} }
  };
\end{tikzpicture}
\legende{Carte mentale de la leçon : les grands groupes humains du Cameroun (localisation, organisations socio-culturelles, répartition de la population et brassage).}
\end{popfigure}

\begin{bilan}
\textbf{Définitions clés}
\begin{itemize}
  \item \cle{Groupe humain} : ensemble de populations partageant des traits linguistiques, culturels et des modes d'organisation sociale communs.
  \item \cle{Paléo-soudanais} : « anciens Soudanais », populations autochtones du Nord (Kotoko, Masa, Toupouri, Mousgoum), présentes avant les migrations peules.
  \item \cle{Néo-soudanais} : « nouveaux Soudanais », populations arrivées plus tardivement au Nord (Peuls/Foulbé au XIX\textsuperscript{e} siècle, arabes Choa).
  \item \cle{Bantou} : grand groupe linguistique du Centre, du Sud, du Littoral et d'une partie de l'Est (Ewondo, Bulu, Douala, Bassa, Fang).
  \item \cle{Semi-bantou} : groupe intermédiaire des hautes terres de l'Ouest et du Nord-Ouest (Bamiléké, Bamoun, Tikars, Widekum).
  \item \cle{Pygmées} : premiers habitants de la forêt du Sud-Est (Baka, Bagyeli, Bakola), traditionnellement chasseurs-cueilleurs.
  \item \cle{Densité de population} : nombre d'habitants par km\textsuperscript{2} ; \cle{brassage} : mélange des populations par les migrations.
\end{itemize}

\textbf{Repères chiffrés et localisations à connaître}
\begin{center}
\begin{tabularx}{\linewidth}{|l|X|X|}\hline
\rowcolor{popblueL} \textbf{Groupe} & \textbf{Localisation} & \textbf{Organisation dominante} \\ \hline
Paléo-soudanais & Extrême-Nord (plaine du Logone, monts Mandara) & Chefferies, agriculture et pêche, religions traditionnelles puis islam \\ \hline
Néo-soudanais & Nord, Adamaoua & Lamidats, islam, élevage bovin \\ \hline
Bantou & Centre, Sud, Littoral, Est & Sociétés lignagères, chefferies souvent peu centralisées \\ \hline
Semi-bantou & Ouest, Nord-Ouest & Chefferies puissantes et hiérarchisées (Foumban, Bandjoun) \\ \hline
Pygmées & Forêt du Sud-Est (région de l'Est) & Campements forestiers, chasse et cueillette \\ \hline
\end{tabularx}
\end{center}

\textbf{Densités (ordres de grandeur, recensements récents)}
\begin{itemize}
  \item Densité nationale moyenne : environ 50 hab/km\textsuperscript{2} (ordre de grandeur).
  \item Zones \cle{densément peuplées} : Ouest et Nord-Ouest (plus de 100 hab/km\textsuperscript{2} par endroits), Littoral (Douala), Extrême-Nord, Yaoundé.
  \item Zones \cle{moyennement peuplées} : Centre, Sud-Ouest, Nord.
  \item Zones \cle{faiblement peuplées} : Est et Adamaoua (souvent moins de 10 hab/km\textsuperscript{2}), Sud.
\end{itemize}

\textbf{Méthodes express}
\begin{itemize}
  \item \cle{Lire une carte de densités} : lire la légende, repérer les trois niveaux de peuplement, localiser, expliquer (relief, climat, histoire, activités).
  \item \cle{Construire la carte (TP 2)} : choisir un découpage en 3 classes, associer une gamme de couleurs (du foncé au clair), colorier les régions, rédiger titre, légende, orientation et échelle.
  \item \cle{Calculer une densité} : $d = \dfrac{\text{population}}{\text{superficie}}$ (hab/km\textsuperscript{2}).
\end{itemize}
\end{bilan}

\begin{aretenir}[Checklist avant le Bac]
\begin{itemize}
  \item $\square$ Je sais définir paléo-soudanais, néo-soudanais, bantou, semi-bantou et pygmées.
  \item $\square$ Je sais localiser chaque grand groupe humain sur une carte du Cameroun.
  \item $\square$ Je connais au moins deux ethnies représentatives par groupe (Kotoko, Peuls, Ewondo, Bamiléké, Baka).
  \item $\square$ Je sais décrire les organisations socio-culturelles : lamidats du Nord, chefferies de l'Ouest, sociétés lignagères du Sud.
  \item $\square$ Je connais l'ordre de grandeur de la densité moyenne nationale (environ 50 hab/km\textsuperscript{2}).
  \item $\square$ Je sais citer les régions densément peuplées (Ouest, Littoral, Extrême-Nord) et faiblement peuplées (Est, Adamaoua).
  \item $\square$ Je sais expliquer les contrastes de densité par des facteurs naturels et humains.
  \item $\square$ Je sais calculer une densité de population et un taux de variation.
  \item $\square$ Je sais construire une carte de répartition en trois classes avec titre, légende, nord et échelle.
  \item $\square$ Je sais expliquer le rôle des migrations et du brassage dans l'intégration nationale.
  \item $\square$ Je rédige mes réponses avec le vocabulaire géographique exact (densité, foyer de peuplement, vide humain).
\end{itemize}
\end{aretenir}

\findecours

\end{document}
